%% ============================================================
%%  ECON 621: Applied Econometrics — Textbook
%%  University of Northern British Columbia
%%  Prepared for ECON 621 (cross-listed senior UG / graduate)
%% ============================================================
\documentclass[11pt, twoside, openright]{book}

%% ── Encoding & fonts ──────────────────────────────────────
\usepackage[T1]{fontenc}
\usepackage[utf8]{inputenc}
% Declare Unicode math/punct chars that may slip through:
\DeclareUnicodeCharacter{2212}{$-$}         % − MINUS SIGN
\DeclareUnicodeCharacter{2248}{$\approx$}   % ≈
\DeclareUnicodeCharacter{2260}{$\neq$}      % ≠
\DeclareUnicodeCharacter{2261}{$\equiv$}    % ≡
\DeclareUnicodeCharacter{2264}{$\leq$}      % ≤
\DeclareUnicodeCharacter{2265}{$\geq$}      % ≥
\DeclareUnicodeCharacter{2192}{$\to$}       % →
\DeclareUnicodeCharacter{2190}{$\leftarrow$}% ←
\DeclareUnicodeCharacter{00D7}{$\times$}    % ×
\DeclareUnicodeCharacter{22A5}{$\perp$}     % ⊥
\DeclareUnicodeCharacter{2208}{$\in$}       % ∈
\DeclareUnicodeCharacter{2026}{\ldots}      % …
\DeclareUnicodeCharacter{00B7}{$\cdot$}     % ·
\DeclareUnicodeCharacter{2014}{---}         % — em dash
\DeclareUnicodeCharacter{2013}{--}          % – en dash
\DeclareUnicodeCharacter{201C}{``}          % "
\DeclareUnicodeCharacter{201D}{''}          % "
\DeclareUnicodeCharacter{2018}{`}           % '
\DeclareUnicodeCharacter{2019}{'}           % '
\usepackage{lmodern}
\usepackage{microtype}

%% ── Page geometry ─────────────────────────────────────────
\usepackage[
  paperwidth=7.5in, paperheight=10in,
  top=1.1in, bottom=1.1in,
  inner=1.1in, outer=0.85in,
  headsep=0.25in,
  headheight=14pt
]{geometry}

%% ── Mathematics ───────────────────────────────────────────
\usepackage{amsmath, amssymb, amsthm, mathtools, bm}
\usepackage{mathrsfs}

%% ── Theorem environments ─────────────────────────────────
\newtheoremstyle{econstyle}
  {8pt}{6pt}{\itshape}{}{\bfseries\color{econblue}}{.}{ }{}
\theoremstyle{econstyle}
\newtheorem{theorem}{Theorem}[chapter]
\newtheorem{lemma}[theorem]{Lemma}
\newtheorem{proposition}[theorem]{Proposition}
\newtheorem{corollary}[theorem]{Corollary}

\newtheoremstyle{defstyle}
  {8pt}{6pt}{\normalfont}{}{\bfseries\color{econblue}}{.}{ }{}
\theoremstyle{defstyle}
\newtheorem{definition}{Definition}[chapter]
\newtheorem{assumption}{Assumption}[chapter]
\newtheorem{example}{Example}[chapter]
\newtheorem{remark}{Remark}[chapter]

%% ── Colors ────────────────────────────────────────────────
\usepackage[dvipsnames, table]{xcolor}
\definecolor{econblue}{RGB}{0, 64, 128}
\definecolor{boxbg}{RGB}{235, 243, 255}
\definecolor{shadedbg}{RGB}{248, 248, 248}
\definecolor{codebg}{RGB}{248, 248, 248}
\definecolor{codeframe}{RGB}{180, 180, 180}

%% ── Boxes & callouts ──────────────────────────────────────
\usepackage{tcolorbox}
\tcbuselibrary{skins, breakable, theorems}

%% keybox: mandatory arg {title} – avoids pgfkeys parsing issues with
%%         math or commas in optional [title] arguments.
\newenvironment{keybox}[1]{%
  \begin{tcolorbox}[colback=boxbg, colframe=econblue,
    breakable, enhanced,
    before upper={\noindent\textbf{\color{econblue}#1}\par\smallskip}]
}{%
  \end{tcolorbox}%
}

\newenvironment{warningbox}[1]{%
  \begin{tcolorbox}[colback=yellow!10, colframe=orange!80!black,
    breakable, enhanced,
    before upper={\noindent\textbf{\color{orange!80!black}#1}\par\smallskip}]
}{%
  \end{tcolorbox}%
}

\newtcolorbox{notebox}{
  colback=shadedbg, colframe=gray!40,
  fontupper=\small, breakable
}

%% ── Code listings ─────────────────────────────────────────
\usepackage{listings}
\lstset{inputencoding=utf8, extendedchars=true,
  literate={←}{$\leftarrow$}{1}{→}{$\rightarrow$}{1}
           {≈}{$\approx$}{1}{≠}{$\neq$}{1}{≡}{$\equiv$}{1}
           {−}{$-$}{1}{×}{$\times$}{1}{≤}{$\leq$}{1}
           {≥}{$\geq$}{1}{∈}{$\in$}{1}{⊥}{$\perp$}{1}
           {…}{\ldots}{1}{·}{$\cdot$}{1}
}
\lstdefinestyle{Rstyle}{
  language=R,
  backgroundcolor=\color{codebg},
  basicstyle=\ttfamily\footnotesize,
  keywordstyle=\color{blue!70!black}\bfseries,
  commentstyle=\color{green!50!black}\itshape,
  stringstyle=\color{orange!80!black},
  numberstyle=\tiny\color{gray},
  frame=single, frameround=tttt,
  rulecolor=\color{codeframe},
  numbers=left, numbersep=6pt,
  breaklines=true, breakatwhitespace=true,
  showstringspaces=false,
  tabsize=2,
  captionpos=b,
  xleftmargin=8pt,
  morekeywords={library, data, ggplot, aes, geom_point, geom_line,
    geom_smooth, labs, theme_minimal, lm, summary, coef, vcov,
    residuals, fitted, predict, anova, rbind, cbind, matrix,
    solve, diag, t, nrow, ncol, round, cat, mean, sd, var,
    cor, set.seed, rnorm, rbinom, runif, replicate, map_dfr,
    as_tibble, mutate, filter, group_by, summarise, left_join,
    pivot_longer, pivot_wider, factor, case_when, ifelse,
    pdata.frame, plm, phtest, feols, iplot, ivreg, rdrobust,
    matchit, weightit, rddensity, rdplotdensity}
}
\lstset{style=Rstyle}

%% ── Tables ────────────────────────────────────────────────
\usepackage{booktabs, longtable, array, multirow, tabularx}
\renewcommand{\arraystretch}{1.3}

%% ── Graphics ──────────────────────────────────────────────
\usepackage{graphicx, float}
\usepackage{tikz}
\usepackage{pgfplots}
\pgfplotsset{compat=1.18}
\usetikzlibrary{arrows.meta, shapes.geometric, positioning, calc, decorations.pathreplacing, patterns, fit, backgrounds}
\tikzset{
  node/.style={circle, draw=econblue, fill=econblue!10, thick, minimum size=8mm, font=\small},
  arrow/.style={-{Stealth[scale=1.2]}, thick, econblue},
  dasharrow/.style={-{Stealth[scale=1.2]}, thick, dashed, gray},
}

%% ── Headers & footers ─────────────────────────────────────
\usepackage{fancyhdr}
\pagestyle{fancy}
\fancyhf{}
\fancyhead[LE]{\small\textcolor{econblue}{\itshape ECON 621: Applied Econometrics}}
\fancyhead[RO]{\small\textcolor{econblue}{\itshape\leftmark}}
\fancyfoot[LE,RO]{\small\thepage}
\renewcommand{\headrulewidth}{0.4pt}
\fancypagestyle{plain}{\fancyhf{}\fancyfoot[C]{\small\thepage}\renewcommand{\headrulewidth}{0pt}}

%% ── Chapter title style ───────────────────────────────────
\usepackage{titlesec}
\titleformat{\chapter}[display]
  {\normalfont\Huge\bfseries\color{econblue}}
  {\chaptertitlename\ \thechapter}{14pt}{\Huge}
\titlespacing*{\chapter}{0pt}{-10pt}{30pt}

\titleformat{\section}{\large\bfseries\color{econblue!80!black}}{\thesection}{0.8em}{}
\titleformat{\subsection}{\normalsize\bfseries\color{econblue!60!black}}{\thesubsection}{0.7em}{}
\titleformat{\subsubsection}{\normalsize\bfseries}{\thesubsubsection}{0.6em}{}

%% ── Table of contents ─────────────────────────────────────
\usepackage[titles]{tocloft}
\setcounter{tocdepth}{2}
\renewcommand{\cftchapfont}{\bfseries\color{econblue}}

%% ── Hyperlinks ────────────────────────────────────────────
\usepackage[
  colorlinks=true,
  linkcolor=econblue,
  citecolor=econblue!80!black,
  urlcolor=econblue!70!black,
  bookmarks=true,
  bookmarksnumbered=true
]{hyperref}

%% ── Bibliography ──────────────────────────────────────────
\usepackage{natbib}
\bibliographystyle{plainnat}

%% ── Miscellaneous ─────────────────────────────────────────
\usepackage{enumitem}
\setlist[itemize]{leftmargin=1.5em, itemsep=2pt}
\setlist[enumerate]{leftmargin=1.5em, itemsep=2pt}
\usepackage{parskip}
\setlength{\parskip}{4pt}
\usepackage{setspace}
\usepackage{epigraph}

%% ── Custom commands ───────────────────────────────────────
\newcommand{\E}{\mathbb{E}}
\newcommand{\V}{\mathrm{Var}}
\newcommand{\Cov}{\mathrm{Cov}}
\newcommand{\plim}{\operatorname{plim}}
\newcommand{\xrightarrowp}{\xrightarrow{p}}
\newcommand{\xrightarrowd}{\xrightarrow{d}}
\newcommand{\bhat}{\hat{\boldsymbol{\beta}}}
\newcommand{\bols}{\hat{\boldsymbol{\beta}}^{OLS}}
\newcommand{\Px}{\mathbf{P_X}}
\newcommand{\Mx}{\mathbf{M_X}}
\newcommand{\XX}{\mathbf{X}^T\mathbf{X}}
\newcommand{\Xy}{\mathbf{X}^T\mathbf{y}}
\newcommand{\MX}{\mathbf{M_X}}

%% ── Title page info ───────────────────────────────────────
\title{%
  \vspace{-1cm}
  {\color{econblue}\rule{\textwidth}{2pt}}\\[0.5em]
  {\Huge\bfseries\color{econblue} Applied Econometrics}\\[0.4em]
  {\large A Textbook for ECON 621}\\[0.3em]
  {\color{econblue}\rule{\textwidth}{1pt}}
}
\author{%
  University of Northern British Columbia\\[0.3em]
  School of Business and Economics\\[0.5em]
  \small Based on ECON 712: Applied Econometrics (Winter 2026)
}
\date{%
  \small Winter 2027\\[0.5em]
  \small\itshape Drawing on: Angrist \& Pischke (2009) $\cdot$
         Davidson \& MacKinnon (2004) $\cdot$
         Cunningham (2021) $\cdot$ Wooldridge (2019)
}

%% ===========================================================
\begin{document}

\frontmatter

%% ── Title page ────────────────────────────────────────────
\maketitle
\thispagestyle{empty}

%% ── Copyright / disclaimer ────────────────────────────────
\newpage
\thispagestyle{empty}
\vspace*{\fill}
\begin{center}
{\small
\textbf{ECON 621: Applied Econometrics}\\
University of Northern British Columbia\\
School of Business and Economics\\[1em]
These lecture notes are prepared for course use only.\\
They draw extensively on the following foundational texts:\\[0.5em]
Angrist, J.~D., and Pischke, J.~S. (2009).
\textit{Mostly Harmless Econometrics}. Princeton University Press.\\[0.2em]
Davidson, R., and MacKinnon, J.~G. (2004).
\textit{Econometric Theory and Methods}. Oxford University Press.\\[0.2em]
Cunningham, S. (2021). \textit{Causal Inference: The Mixtape}. Yale University Press.\\[0.2em]
Wooldridge, J.~M. (2019). \textit{Introductory Econometrics}, 7th ed. Cengage.\\[1em]
ECON 712 lecture materials by Prof.\ Leandro Freylejer, UNBC, Winter 2026.
}
\end{center}
\vspace*{\fill}

%% ── Table of Contents ─────────────────────────────────────
\tableofcontents

%% ── Preface ───────────────────────────────────────────────
\chapter*{Preface}
\addcontentsline{toc}{chapter}{Preface}

This textbook is written for ECON 621: Applied Econometrics at the University of Northern British Columbia.
The course is cross-listed as a senior undergraduate and graduate course, and these notes are designed to
serve both audiences: they develop the material rigorously enough for graduate students while maintaining
the accessibility and applied focus appropriate for advanced undergraduates.

The notes are organized in three parts.
\textbf{Part I: Foundations} (Chapters 1--5) covers the regression framework rigorously, from the
potential outcomes model through OLS algebra, finite-sample theory, and asymptotic inference with
robust standard errors. \textbf{Part II: The Regression Toolkit} (Chapters 6--7) extends the
framework with the Frisch-Waugh-Lovell theorem, omitted variable bias, and flexible specifications.
\textbf{Part III: Identification Strategies} (Chapters 8--13) covers instrumental variables,
the Local Average Treatment Effect, panel data, differences-in-differences, matching, and regression
discontinuity --- the core of modern applied econometrics.

Every chapter includes formal definitions and proofs, economic intuition, canonical empirical examples
from the literature, exercises at multiple levels, and a complete \textbf{R programming lab} with
working code that students can run directly.

The approach follows the \textit{Credibility Revolution} tradition: we are as concerned with
\emph{research design} as with estimation. Understanding what variation identifies an effect,
and whether that identification strategy is credible, is the central skill of applied econometrics.

\bigskip
\noindent\textit{Note to students:} Do not just read these notes --- work through the derivations
with pencil and paper, run the R code, and modify it to test your understanding. Applied econometrics
is learned by doing.

\mainmatter

%% ===========================================================
%%  PART I: FOUNDATIONS
%% ===========================================================
\part{Foundations}

\chapter{Chapter 1: The Quest for Causality --- An Introduction to Applied Econometrics}

\medskip\hrule\medskip

\section{Chapter Overview}

This chapter introduces the central challenge of applied econometrics: distinguishing \textbf{correlation} from \textbf{causation} in economic data. We begin by asking what kinds of questions economists care about, explain why answering them with observational data is hard, and introduce the conceptual framework --- the \textbf{potential outcomes model} --- that modern applied economists use to think clearly about causality.

By the end of this chapter you will understand why simply comparing group averages can be deeply misleading, what it means to "identify" a causal effect, and what tools the rest of this book will give you to do so.

\medskip\hrule\medskip

\section{Learning Objectives}

After completing this chapter, students will be able to:

\begin{enumerate}
\item Define applied econometrics and explain its goals.
\item Formulate a well-posed empirical question in terms of a causal relationship.
\item Write down the \textbf{potential outcomes} (counterfactual) model for a binary treatment.
\item Define the \textbf{Average Treatment Effect (ATE)}, the \textbf{Average Treatment Effect on the Treated (ATT)}, and relate them to the observable difference in means.
\item Identify and explain \textbf{selection bias} and why it contaminates simple comparisons.
\item Explain how \textbf{random assignment} solves the selection bias problem.
\item Describe, at a high level, the identification strategies studied in this book (regression, IV, DiD, matching, RDD).
\end{enumerate}

\medskip\hrule\medskip

\section{1.1 What Is Applied Econometrics?}

Economists study a wide variety of questions: Does access to health insurance improve health outcomes? Do minimum wage laws cause unemployment? Does attending a smaller class improve student test scores? Does going to university raise lifetime earnings?

What unites all of these questions is that they are \textbf{causal}. We do not merely want to know whether smaller class sizes and higher test scores tend to go together in data --- we want to know whether \textit{reducing} class size \textit{causes} test scores to rise, and by how much. The distinction matters enormously for policy. If better-resourced schools happen to have both smaller classes and better teachers, then we might find a strong correlation between class size and test scores even if class size has no causal effect whatsoever.

\textbf{Applied econometrics} is the branch of economics that uses \textbf{data and statistical methods} to estimate causal relationships and test economic theories. It sits at the intersection of economic theory, statistics, and data science, but its animating concern --- separating cause from correlation --- makes it distinctive.

\subsection{1.1.1 A Map of the Terrain}

Applied econometrics divides naturally into two broad activities:

\begin{enumerate}
\item \textbf{Estimating causal effects.} How large is the return to an additional year of schooling? Does a job-training program raise employment? By how much did a particular policy change crime rates?
\end{enumerate}

\begin{enumerate}
\item \textbf{Testing economic theories.} Do consumers respond to prices as standard theory predicts? Is the labour supply curve upward sloping? Are financial markets informationally efficient?
\end{enumerate}

This course focuses primarily on (1), but the tools we develop are equally useful for (2). Throughout, we will be careful to distinguish between \textbf{descriptive} statements (what is the world like?) and \textbf{causal} statements (what would happen if we intervened?).

\subsection{1.1.2 Data in Applied Econometrics}

Economists work with many types of data:

\begin{itemize}
\item \textbf{Cross-sectional data} --- observations on many units (individuals, firms, countries) at a single point in time.
\item \textbf{Time-series data} --- observations on a single unit over many time periods.
\item \textbf{Panel (longitudinal) data} --- observations on many units over multiple time periods. This is the richest structure and features prominently in this book.
\item \textbf{Experimental data} --- data collected through randomized controlled trials (RCTs).
\end{itemize}

Each data type raises its own challenges. Observational data --- data collected without deliberate experimental design --- is by far the most common, and learning to draw valid causal inferences from it is the central skill of applied econometrics.

\medskip\hrule\medskip

\section{1.2 Empirical Questions and the Causal Ideal}

The starting point of any empirical project is a \textbf{well-defined question}. Precision matters. Consider the following progression:

\begin{center}
\begin{tabular}{@{} l l @{}}
\toprule
\textbf{Vague Question} & \textbf{Precise Causal Question} \\
\midrule
Do hospitals make people healthier? & Does hospitalization \textit{cause} better health outcomes for elderly patients? \\
Does education matter? & Does an additional year of schooling \textit{cause} higher lifetime earnings? \\
Is the minimum wage bad for employment? & Does a \$1 increase in the minimum wage \textit{cause} a reduction in fast-food employment? \\
\bottomrule
\end{tabular}
\end{center}

The precise version specifies: (a) the \textbf{cause} (the treatment or intervention), (b) the \textbf{effect} (the outcome), and (c) the \textbf{population} of interest.

\subsection{1.2.1 Fundamental Concepts}

\textbf{Definition 1.1 (Treatment):} A \textit{treatment} is a variable $D_i$ that takes the value 1 if unit $i$ receives the intervention (is "treated") and 0 otherwise.

\textbf{Definition 1.2 (Outcome):} An \textit{outcome} $Y_i$ is the variable we care about --- earnings, health, test scores, employment.

\textbf{Definition 1.3 (Causal Effect):} The causal effect of the treatment on unit $i$ is the difference between the outcome they would experience \textit{with} treatment and the outcome they would experience \textit{without} treatment.

This definition immediately reveals the \textbf{fundamental problem of causal inference}: we can never observe the same unit simultaneously with and without the treatment. Whatever we observe is only one side of a counterfactual comparison.

\medskip\hrule\medskip

\section{1.3 The Potential Outcomes Framework}

The most productive way to think about causal questions is the \textbf{potential outcomes framework}, originally developed by Neyman (1923) and formalized by Rubin (1974). This framework is sometimes called the \textbf{Rubin Causal Model} or the \textbf{Neyman-Rubin framework}.

\subsection{1.3.1 Potential Outcomes Defined}

For each individual $i$ in the population, define two \textbf{potential outcomes}:

\[
Y_{0i} = \text{the outcome for unit } i \text{ if } D_i = 0 \text{ (not treated)}
\]

\[
Y_{1i} = \text{the outcome for unit } i \text{ if } D_i = 1 \text{ (treated)}
\]

\begin{keybox}{Key Insight:** Both $Y_{0i}$ and $Y_{1i}$ are defined for \textit{every} individual, but we can only ever \textit{observe} one of them --- the one corresponding to the treatment the individual actually received.}
\end{keybox}

The \textbf{individual causal effect} for unit $i$ is:

\[
\tau_i = Y_{1i} - Y_{0i}
\]

This is the difference that would result from "switching" the individual's treatment status, everything else held constant.

\begin{figure}[htbp]
  \centering
  \begin{tikzpicture}[scale=1.1]
    % Two columns: treated and untreated
    \node[font=\bfseries\color{econblue}] at (1.5, 3.8) {Treatment Group ($D_i=1$)};
    \node[font=\bfseries\color{econblue}] at (5.5, 3.8) {Control Group ($D_i=0$)};

    % Treated column
    \draw[fill=econblue!15, draw=econblue, thick, rounded corners] (0,0) rectangle (3,3.4);
    \node[font=\small] at (1.5, 2.8) {\textbf{Observed:} $Y_{1i} = Y_i(1)$};
    \node[font=\small] at (1.5, 2.0) {(outcome with treatment)};
    \draw[dashed, gray, thick, rounded corners] (0.2,0.2) rectangle (2.8,1.5);
    \node[font=\small, gray] at (1.5, 1.1) {\textbf{Counterfactual:} $Y_{0i}$};
    \node[font=\small, gray] at (1.5, 0.5) {(never observed)};

    % Control column
    \draw[fill=gray!10, draw=gray!50, thick, rounded corners] (4,0) rectangle (7,3.4);
    \node[font=\small] at (5.5, 2.8) {\textbf{Observed:} $Y_{0i} = Y_i(0)$};
    \node[font=\small] at (5.5, 2.0) {(outcome without treatment)};
    \draw[dashed, gray, thick, rounded corners] (4.2,0.2) rectangle (6.8,1.5);
    \node[font=\small, gray] at (5.5, 1.1) {\textbf{Counterfactual:} $Y_{1i}$};
    \node[font=\small, gray] at (5.5, 0.5) {(never observed)};

    % Causal effect arrow
    \draw[-{Stealth}, thick, econblue] (3.1, 1.8) -- (3.9, 1.8)
      node[midway, above, font=\small\bfseries, econblue] {$\delta_i = Y_{1i}-Y_{0i}$};
    \node[font=\tiny, gray] at (3.5, 1.5) {(never observed)};
  \end{tikzpicture}
  \caption{The fundamental problem of causal inference. For each unit, only one potential outcome is observed. The causal effect $\delta_i = Y_{1i} - Y_{0i}$ can never be computed directly.}
  \label{fig:ch01_potential_outcomes}
\end{figure}

\subsection{1.3.2 Population-Level Causal Estimands}

Because $\tau_i$ varies across individuals and we cannot observe it for any individual, economists focus on \textbf{average} causal effects:

\textbf{Definition 1.4 (Average Treatment Effect):}

\[
\text{ATE} = E[Y_{1i} - Y_{0i}] = E[\tau_i]
\]

The ATE is the average causal effect for a randomly chosen unit from the population.

\textbf{Definition 1.5 (Average Treatment Effect on the Treated):}

\[
\text{ATT} = E[Y_{1i} - Y_{0i} \mid D_i = 1]
\]

The ATT is the average causal effect for individuals who actually received the treatment. This is often the more policy-relevant quantity: when a government wants to know if a training program helped the participants, it cares about the ATT.

\textbf{Definition 1.6 (Average Treatment Effect on the Untreated):}

\[
\text{ATU} = E[Y_{1i} - Y_{0i} \mid D_i = 0]
\]

The ATU describes what the average effect would be if the untreated were brought into the program.

\begin{keybox}{Note on ATE vs ATT:** When treatment effects are heterogeneous (vary by individual) and treatment is not randomly assigned, ATE and ATT need not be equal. Understanding the distinction is important for interpreting estimates and for policy.}
\end{keybox}

\subsection{1.3.3 What We Observe}

In any dataset, we observe the \textit{realized} outcome:

\[
Y_i = Y_{0i} + (Y_{1i} - Y_{0i}) D_i = \begin{cases} Y_{1i} & \text{if } D_i = 1 \\ Y_{0i} & \text{if } D_i = 0 \end{cases}
\]

This is the \textbf{switching equation}. It links the observable data to the underlying potential outcomes.

\medskip\hrule\medskip

\section{1.4 Selection Bias: Why Simple Comparisons Mislead}

The most naive approach to estimating a causal effect is to compare the average outcomes of treated and untreated units. Let us see precisely why this fails.

\subsection{1.4.1 The Decomposition of Observed Differences}

Consider the difference in mean outcomes between treated and untreated groups:

\[
\underbrace{E[Y_i \mid D_i = 1] - E[Y_i \mid D_i = 0]}_{\text{Observed Difference in Means}}
\]

Substituting the switching equation:

\[
= E[Y_{1i} \mid D_i = 1] - E[Y_{0i} \mid D_i = 0]
\]

Now add and subtract $E[Y_{0i} \mid D_i = 1]$:

\[
= \underbrace{E[Y_{1i} - Y_{0i} \mid D_i = 1]}_{\text{ATT}} + \underbrace{E[Y_{0i} \mid D_i = 1] - E[Y_{0i} \mid D_i = 0]}_{\text{Selection Bias}}
\]

\textbf{This is the fundamental decomposition of observed differences into causal effects and bias.}

The first term is the ATT --- what we actually want to know.

The second term, \textbf{selection bias}, is the difference in average untreated outcomes between the treated and untreated groups. It captures the fact that people who select into treatment may be systematically different from those who do not.

\begin{keybox}{Equation 1.1: Bias Decomposition}
\end{keybox}
\[
\underbrace{E[Y_i|D_i=1] - E[Y_i|D_i=0]}_{\text{Observed difference}} = \underbrace{\text{ATT}}_{\text{Causal effect}} + \underbrace{E[Y_{0i}|D_i=1] - E[Y_{0i}|D_i=0]}_{\text{Selection bias}}
\]

\begin{figure}[htbp]
  \centering
  \begin{tikzpicture}[
    mynode/.style={circle, draw=econblue, fill=econblue!10, thick, minimum size=10mm, font=\small\bfseries},
    confnode/.style={circle, draw=red!70!black, fill=red!10, thick, minimum size=10mm, font=\small\bfseries},
    myarrow/.style={-{Stealth[scale=1.2]}, thick, econblue},
    confarrow/.style={-{Stealth[scale=1.2]}, thick, red!70!black, dashed},
    scale=1.2
  ]
    \node[mynode] (D) at (0, 0) {$D_i$};
    \node[mynode] (Y) at (4, 0) {$Y_i$};
    \node[confnode] (U) at (2, 2.2) {$U_i$};

    \draw[myarrow] (D) -- (Y) node[midway, below, font=\small] {causal effect};
    \draw[confarrow] (U) -- (D) node[midway, left, font=\small, red!70!black] {selection};
    \draw[confarrow] (U) -- (Y) node[midway, right, font=\small, red!70!black] {direct effect};

    \node[font=\footnotesize, gray] at (2, -0.6) {$\underbrace{E[Y_i|D_i=1]-E[Y_i|D_i=0]}_{\text{observed difference}} = \underbrace{E[Y_{1i}-Y_{0i}|D_i=1]}_{\text{ATT}} + \underbrace{E[Y_{0i}|D_i=1]-E[Y_{0i}|D_i=0]}_{\text{selection bias}}$};
  \end{tikzpicture}
  \caption{A directed acyclic graph (DAG) illustrating selection bias. The unobserved confounder $U_i$ (red, dashed) simultaneously affects treatment assignment $D_i$ and the outcome $Y_i$, creating a spurious association between treatment and outcome beyond the true causal effect.}
  \label{fig:ch01_selection_dag}
\end{figure}

\subsection{1.4.2 A Canonical Example: Hospitals and Health}

Consider the question: do hospitals make people healthier? The National Health Interview Survey (NHIS) data reveal a troubling fact --- individuals who were hospitalized in the past year report \textit{worse} health status on average than those who were not. Does this mean hospitals make you sick?

Of course not. The selection bias here runs in a clear direction: people who end up in hospitals are already, on average, much sicker than the general population. They have worse $Y_{0i}$ --- they would have had worse health \textit{even without} hospitalization. So:

\[
E[Y_{0i} \mid D_i = 1] < E[Y_{0i} \mid D_i = 0]
\]

The selection bias is \textit{negative}, masking what may well be a large positive causal effect of hospitalization on health. The observable comparison between hospitalized and non-hospitalized individuals is dominated by the fact that hospitalized people were sicker to begin with.

This example, from Angrist and Pischke (2009), illustrates how selection bias can flip the apparent sign of a causal relationship.

\subsection{1.4.3 Selection Bias in Labour Economics}

A similar issue arises in estimating returns to education. People who choose to obtain more education tend to be more able, more motivated, and from higher-income families. Even if education had \textit{zero} causal effect on earnings, we would observe a positive correlation between education and earnings because of these underlying differences.

Formally:

\[
E[Y_{0i} \mid D_i = 1] > E[Y_{0i} \mid D_i = 0]
\]

where $D_i = 1$ denotes more education. The selection bias is positive, and simple OLS regressions of wages on education will overstate the causal return to schooling. (We return to this in Chapters 6 and 8.)

\medskip\hrule\medskip

\section{1.5 The Gold Standard: Randomized Experiments}

The cleanest solution to the selection bias problem is \textbf{random assignment}. In a randomized controlled trial (RCT), whether an individual receives treatment is determined by a coin flip (or random number generator), not by any characteristic of the individual.

\subsection{1.5.1 Why Randomization Works}

If $D_i$ is assigned randomly, then it is independent of potential outcomes:

\[
\{Y_{0i}, Y_{1i}\} \perp\!\!\!\perp D_i
\]

This independence condition ensures that:

\[
E[Y_{0i} \mid D_i = 1] = E[Y_{0i} \mid D_i = 0] = E[Y_{0i}]
\]

That is, the untreated outcomes of treated and untreated individuals are the same in expectation. \textbf{Selection bias is zero.}

As a result:

\[
E[Y_i \mid D_i = 1] - E[Y_i \mid D_i = 0] = E[Y_{1i} - Y_{0i}] = \text{ATE}
\]

The simple difference in group means identifies the Average Treatment Effect.

\subsection{1.5.2 Linking Randomization to Regression}

To connect experiments to the regression framework, write:

\[
Y_i = \alpha + \rho D_i + \epsilon_i
\]

where $\alpha = E[Y_{0i}]$, $\rho = Y_{1i} - Y_{0i}$ (assumed constant here), and $\epsilon_i = Y_{0i} - E[Y_{0i}]$.

Under random assignment, $D_i$ is independent of $Y_{0i}$, so $D_i$ is uncorrelated with $\epsilon_i$. The OLS estimator of $\rho$ from this regression recovers the causal effect:

\[
\hat{\rho}^{OLS} \xrightarrow{p} \rho
\]

This is the bridge between the potential outcomes framework and regression. Regression with a randomly assigned treatment variable identifies causal effects. The problem --- and most of this course --- is about what to do when treatment is \textit{not} randomly assigned.

\subsection{1.5.3 Example: Project STAR}

One of the most celebrated experiments in economics is Project STAR (Student Teacher Achievement Ratio), conducted in Tennessee from 1985 to 1989. Over 11,000 students and their teachers were randomly assigned to small classes (13--17 students) or regular classes (22--26 students) from kindergarten through grade 3.

Because assignment was random, differences in test scores between students in small and large classes can be attributed to class size, not to preexisting differences in student ability or family background. Krueger (1999) found that students assigned to smaller classes scored roughly 4--5 percentile points higher on standardized tests --- and that these effects were larger for minority and disadvantaged students.

This example illustrates two things: the power of randomization to answer causal questions cleanly, and the policy relevance of the answers.

\begin{keybox}{Box 1.1 --- Fundamental Questions (FUQs)}

Not every empirical question can be answered even in principle. A question is a \textit{Fundamental Unidentified Question} (FUQ) if it requires observing a counterfactual that can never be defined. For example: "What would Canada's GDP have been if it had never adopted a central bank?" The counterfactual here is so distant from any observed data that no statistical method can recover it. The potential outcomes framework helps identify which questions are FUQs and which, with clever design, can be answered.
\end{keybox}

\medskip\hrule\medskip

\section{1.6 When Experiments Are Impossible: Identification Strategies}

In most economic settings, running experiments is not feasible --- for ethical, practical, or financial reasons. We cannot randomly assign people to different levels of schooling, different countries to different policies, or firms to different market structures.

Applied econometrics has developed a set of \textbf{identification strategies} --- methods for estimating causal effects from non-experimental (observational) data. The key insight behind each strategy is to find a source of variation in treatment that is "as good as random," even if treatment itself is not randomly assigned.

The major identification strategies covered in this book are:

\begin{center}
\begin{tabular}{@{} l l l @{}}
\toprule
\textbf{Strategy} & \textbf{Core Idea} & \textbf{Chapters} \\
\midrule
\textbf{Multiple Regression / OLS} & Control for observed confounders & 2--7 \\
\textbf{Instrumental Variables (IV)} & Use exogenous variation in a "third variable" & 8--9 \\
\textbf{Panel Data / Fixed Effects} & Control for time-invariant unobservables & 10 \\
\textbf{Differences-in-Differences} & Compare changes across treated and untreated groups & 11 \\
\textbf{Matching} & Match treated and untreated units on observable characteristics & 12 \\
\textbf{Regression Discontinuity} & Exploit arbitrary thresholds in treatment assignment & 13 \\
\bottomrule
\end{tabular}
\end{center}

Each strategy rests on a \textbf{key identifying assumption} that must be plausible given the research question and context. A major theme of this book is that we must be transparent about these assumptions and subject them to scrutiny.

\subsection{1.6.1 The Credibility Revolution}

Beginning in the 1990s, empirical economics underwent a methodological transformation sometimes called the \textbf{Credibility Revolution}. Economists increasingly insisted that empirical claims rest on transparent identification assumptions that could be argued on substantive grounds, rather than on parametric assumptions that are difficult to verify.

This shift was associated with economists like Joshua Angrist, David Card, Guido Imbens, and Alan Krueger (and recognized by the 2021 Nobel Prize in Economic Sciences, awarded to Card, Angrist, and Imbens). Their work emphasized the use of natural experiments --- situations where policy changes, institutional rules, or other features of the world generate quasi-random variation in treatment --- to identify causal effects.

This textbook teaches econometrics in that tradition. We will be as concerned with \textit{research design} (which variation are we using?) as with \textit{estimation} (how are we fitting curves to data?).

\medskip\hrule\medskip

\section{1.7 Roadmap of This Textbook}

This book is organized into three parts.

\textbf{Part I: Foundations (Chapters 2--5)} develops the regression framework rigorously. We cover the Conditional Expectation Function, OLS algebra and geometry, finite-sample and asymptotic statistical properties, and modern robust inference (heteroskedasticity- and cluster-robust standard errors). This is the toolkit you will use throughout.

\textbf{Part II: The Regression Toolkit (Chapters 6--7)} extends the basic framework with essential tools: the Frisch-Waugh-Lovell theorem, omitted variable bias, dummy variables, interaction terms, and non-linearities.

\textbf{Part III: Identification Strategies (Chapters 8--13)} covers each major strategy in turn: instrumental variables, LATE, panel data, differences-in-differences, matching, and regression discontinuity. For each, we develop the theory, discuss the key assumptions, and study canonical empirical applications.

\textbf{Appendices} provide reviews of probability, matrix algebra, and R programming.

\medskip\hrule\medskip

\section{Chapter Summary}

\begin{itemize}
\item Applied econometrics aims to estimate \textbf{causal relationships} from data.
\item The \textbf{potential outcomes framework} defines the causal effect of a treatment for each individual as the difference between two hypothetical outcomes --- one with treatment, one without.
\item Because we cannot observe both potential outcomes for the same individual, causal inference requires comparing different groups.
\item \textbf{Selection bias} arises when treated and untreated units differ in their untreated potential outcomes.
\item The simple observed difference in means equals the ATT plus selection bias.
\item \textbf{Random assignment} eliminates selection bias by ensuring treatment is independent of potential outcomes.
\item When randomization is infeasible, applied econometricians use \textbf{identification strategies} --- research designs that generate quasi-random variation in treatment.
\item The \textbf{Credibility Revolution} has reoriented empirical economics around transparent, design-based identification.
\end{itemize}

\medskip\hrule\medskip

\section{Key Terms}

\textbf{Applied econometrics} --- the use of statistical methods and data to estimate economic relationships and test economic theories.

\textbf{Potential outcomes} --- the hypothetical outcomes $Y_{1i}$ and $Y_{0i}$ that a unit would experience under treatment and control, respectively.

\textbf{Causal effect} --- the individual-level difference $\tau_i = Y_{1i} - Y_{0i}$.

\textbf{ATE (Average Treatment Effect)} --- $E[Y_{1i} - Y_{0i}]$; the average causal effect across all individuals.

\textbf{ATT (Average Treatment Effect on the Treated)} --- $E[Y_{1i} - Y_{0i} \mid D_i = 1]$.

\textbf{Selection bias} --- $E[Y_{0i} \mid D_i = 1] - E[Y_{0i} \mid D_i = 0]$; the difference in counterfactual outcomes between treated and untreated groups.

\textbf{Switching equation} --- $Y_i = Y_{0i} + (Y_{1i} - Y_{0i})D_i$.

\textbf{Randomized controlled trial (RCT)} --- a study in which treatment is assigned randomly.

\textbf{Identification strategy} --- a research design that exploits quasi-random variation in treatment to estimate causal effects.

\textbf{Fundamental problem of causal inference} --- the impossibility of observing both $Y_{1i}$ and $Y_{0i}$ for the same unit.

\medskip\hrule\medskip

\section{Exercises}

\subsection{Conceptual Questions}

\textbf{1.1} Distinguish between correlation and causation. Give an example (not in this chapter) where a strong correlation does not reflect a causal relationship.

\textbf{1.2} Suppose we observe that people who wear helmets while cycling have a higher rate of cycling injuries than those who don't. Does this mean helmets cause injuries? Use the potential outcomes framework to explain the selection bias present in this comparison.

\textbf{1.3} Write down the switching equation for the following scenario: $D_i = 1$ if individual $i$ participates in a job-training program; $Y_i$ is earnings one year after the program.

\textbf{1.4} In your own words, explain the difference between the ATE and the ATT. In what situations would these two quantities be very different? In what situations would they be approximately equal?

\textbf{1.5} Explain why, in a randomized experiment, the simple difference in mean outcomes between treated and control groups identifies the ATE.

\textbf{1.6} Consider Project STAR. Explain why the fact that some students switched class-size assignments after randomization (non-compliance) complicates the interpretation of the results.

\subsection{Analytical Questions}

\textbf{1.7} Suppose:
\begin{itemize}
\item $E[Y_{1i}] = 10$
\item $E[Y_{0i}] = 6$
\item $E[Y_{0i} \mid D_i = 1] = 8$
\item $E[Y_{0i} \mid D_i = 0] = 5$
\end{itemize}

(a) Compute the ATE.
(b) Compute the selection bias.
(c) Compute the observed difference in means $E[Y_i \mid D_i = 1] - E[Y_i \mid D_i = 0]$, assuming constant treatment effects.
(d) Does the observed difference in means overstate or understate the ATE? Why?

\textbf{1.8} Show algebraically that when $D_i \perp \{Y_{0i}, Y_{1i}\}$ (random assignment), the selection bias term is zero and the observed difference in means equals the ATE.

\textbf{1.9} Consider a government program that provides subsidized childcare to low-income families. Describe the selection bias you would expect if you estimated the effect of the program on children's school readiness by comparing children who received subsidized care to those who did not.

\subsection{Applied Questions}

\textbf{1.10} The following table shows average earnings (in \$000/year) for workers who attended university and those who did not, in a hypothetical survey:

\begin{center}
\begin{tabular}{@{} l l l l @{}}
\toprule
\textbf{} & \textbf{Attended University} & \textbf{Did Not Attend} & \textbf{Difference} \\
\midrule
Earnings & \$65 & \$40 & \$25 \\
\bottomrule
\end{tabular}
\end{center}

(a) What assumption would you need to make to interpret this \$25 difference as the causal return to university education?
(b) Is this assumption plausible? Discuss the likely direction of selection bias.
(c) Propose a research design that would allow you to estimate the causal effect more credibly.

\textbf{1.11} Project STAR involved four conditions: small class, regular class, regular class with a teacher's aide. Briefly describe how you would estimate the effect of small classes relative to regular classes, assuming random assignment was perfect.

\medskip\hrule\medskip

\section{R Lab 1: Introduction to R and Causal Thinking}

\subsection{Objectives}
\begin{itemize}
\item Learn basic R syntax and data manipulation
\item Simulate potential outcomes data
\item Illustrate the selection bias problem with simulated data
\item Perform a simple comparison of means (as a naive estimator)
\end{itemize}

\subsection{Setup}

\begin{lstlisting}
# Install and load necessary packages
# install.packages(c("tidyverse", "haven", "stargazer"))
library(tidyverse)
set.seed(42)
\end{lstlisting}

\subsection{Exercise 1.1: Simulating Potential Outcomes}

\begin{lstlisting}
# Simulate a population of 1000 individuals
n <- 1000

# Each individual has two potential outcomes
# Y0: outcome without treatment (e.g., earnings without training)
# Y1: outcome with treatment (e.g., earnings with training)

# True individual-level treatment effect is 5 (in $000)
# But individuals with higher ability (theta) also have higher Y0
theta <- rnorm(n, mean = 10, sd = 3)   # unobserved ability

Y0 <- 30 + 2 * theta + rnorm(n, sd = 5)  # untreated earnings
Y1 <- Y0 + 5                              # treated earnings (constant +5)

# True ATE
true_ATE <- mean(Y1 - Y0)
cat("True ATE:", true_ATE, "\n")   # Should be approximately 5
\end{lstlisting}

\subsection{Exercise 1.2: Simulating Selection Bias}

\begin{lstlisting}
# Now simulate NON-random treatment:
# High-ability individuals are more likely to select into treatment
prob_treat <- plogis(-5 + 0.5 * theta)   # probability of selecting treatment
D <- rbinom(n, 1, prob_treat)            # treatment indicator

# Observed outcome
Y <- D * Y1 + (1 - D) * Y0

# Naive comparison of means
mean_treated   <- mean(Y[D == 1])
mean_untreated <- mean(Y[D == 0])
naive_estimate <- mean_treated - mean_untreated

cat("Mean outcome, treated:", round(mean_treated, 2), "\n")
cat("Mean outcome, untreated:", round(mean_untreated, 2), "\n")
cat("Naive difference in means:", round(naive_estimate, 2), "\n")
cat("True ATE:", round(true_ATE, 2), "\n")
cat("Selection bias:", round(naive_estimate - true_ATE, 2), "\n")
\end{lstlisting}

\textbf{Questions:}
\begin{itemize}
\item How large is the selection bias in your simulation?
\item In what direction does it go? Why?
\item What would happen if you ran an OLS regression of Y on D? Would you recover the true ATE?
\end{itemize}

\subsection{Exercise 1.3: Simulating a Randomized Experiment}

\begin{lstlisting}
# Now randomly assign treatment
D_random <- rbinom(n, 1, 0.5)
Y_random  <- D_random * Y1 + (1 - D_random) * Y0

# Compare means
rct_estimate <- mean(Y_random[D_random == 1]) - mean(Y_random[D_random == 0])
cat("RCT estimate:", round(rct_estimate, 2), "\n")
cat("True ATE:", round(true_ATE, 2), "\n")
cat("Bias under RCT:", round(rct_estimate - true_ATE, 2), "\n")

# OLS regression under random assignment
lm_rct <- lm(Y_random ~ D_random)
summary(lm_rct)
\end{lstlisting}

\textbf{Questions:}
\begin{itemize}
\item How does the RCT estimate compare to the true ATE?
\item What is the role of the constant in the regression?
\item Run the simulation 1000 times (use \texttt{replicate()}) and show that the RCT estimate is unbiased on average.
\end{itemize}

\subsection{Exercise 1.4: Visualizing Selection Bias}

\begin{lstlisting}
# Create data frame
df <- data.frame(Y0 = Y0, Y1 = Y1, D = factor(D, labels = c("Untreated", "Treated")))

# Plot distributions of Y0 by treatment status
ggplot(df, aes(x = Y0, fill = D)) +
  geom_density(alpha = 0.5) +
  labs(
    title = "Distribution of Untreated Potential Outcome (Y0) by Treatment Status",
    subtitle = "Selection bias is visible: treated group has higher Y0 on average",
    x = "Untreated Potential Outcome (Y0)",
    y = "Density",
    fill = "Treatment Status"
  ) +
  theme_minimal()
\end{lstlisting}

\textbf{Interpretation:} If treatment were randomly assigned, the two distributions would overlap almost perfectly. The fact that they are separated visually confirms the presence of selection bias.

\medskip\hrule\medskip

\textit{End of Chapter 1}

\medskip\hrule\medskip

\textbf{References}

Angrist, J. D., and Pischke, J. (2009). \textit{Mostly Harmless Econometrics: An Empiricist's Companion}. Princeton University Press.

Cunningham, S. (2021). \textit{Causal Inference: The Mixtape}. Yale University Press.

Imbens, G. W., and Rubin, D. B. (2015). \textit{Causal Inference for Statistics, Social, and Biomedical Sciences}. Cambridge University Press.

Krueger, A. B. (1999). Experimental estimates of education production functions. \textit{Quarterly Journal of Economics}, 114(2), 497--532.

Neyman, J. (1923). On the application of probability theory to agricultural experiments. Translated in \textit{Statistical Science}, 5(4), 465--472.

Rubin, D. B. (1974). Estimating causal effects of treatments in randomized and nonrandomized studies. \textit{Journal of Educational Psychology}, 66(5), 688--701.

\chapter{Chapter 2: The Conditional Expectation Function and Linear Regression}

\medskip\hrule\medskip

\section{Chapter Overview}

Before we can understand why regression works --- and when it doesn't --- we need a clear statistical framework for describing relationships between variables. This chapter develops the \textbf{Conditional Expectation Function (CEF)}, the population object that regression is trying to estimate. We show that the CEF has deep and useful mathematical properties, that OLS regression is the best linear approximation to the CEF, and that the two coincide when the CEF happens to be linear.

This chapter also connects the regression framework to the causal questions introduced in Chapter 1. Understanding the relationship between regression and the CEF is essential for interpreting regression output correctly and avoiding common mistakes.

\medskip\hrule\medskip

\section{Learning Objectives}

After completing this chapter, students will be able to:

\begin{enumerate}
\item Define the \textbf{Conditional Expectation Function (CEF)} and explain what it represents.
\item State and apply the \textbf{Law of Iterated Expectations (LIE)}.
\item Explain the \textbf{CEF Decomposition Property} and its implications.
\item Prove that the CEF is the \textbf{best predictor} of $Y$ given $X$ in the mean-squared-error sense.
\item Explain the \textbf{ANOVA Theorem} (variance decomposition).
\item State the \textbf{OLS linear regression problem} and derive the population parameter $\boldsymbol{\beta}^{OLS}$.
\item Explain why OLS is the \textbf{best linear approximation} to the CEF even when the CEF is non-linear.
\item Interpret regression coefficients as statements about the CEF.
\item Apply these concepts to real economic data using R.
\end{enumerate}

\medskip\hrule\medskip

\section{2.1 Describing Relationships: From Scatter Plots to Functions}

Consider the question: how do wages vary with education? Figure 2.1 (described below) shows a scatter plot of log weekly earnings against years of schooling, constructed from Current Population Survey (CPS) data. There is a clear positive relationship, but with enormous variation: two people with identical schooling can have very different earnings. What summary of this cloud of points do we want?

The obvious answer is: we want to know the \textbf{average wage for each level of education}. If we could calculate the mean of $Y_i$ (log earnings) for every value of $X_i$ (schooling), we would trace out a curve through the scatter plot. This curve is the \textbf{Conditional Expectation Function}.

\begin{keybox}{Figure 2.1 Description:** Scatter plot of log weekly earnings ($Y$) against years of schooling ($X$) for a random sample of workers. The cloud of points slopes upward on average. Within each level of schooling, wages are spread over a wide range. The line running through the scatter plot represents $E[Y_i \mid X_i = x]$ for each value of $x$.}
\end{keybox}

\medskip\hrule\medskip

\section{2.2 The Conditional Expectation Function (CEF)}

\textbf{Definition 2.1 (CEF):} Let $Y_i$ be a random variable (the outcome) and $\mathbf{X}_i$ be a vector of conditioning variables (regressors). The \textbf{Conditional Expectation Function} is:

\[
E[Y_i \mid \mathbf{X}_i = \mathbf{x}]
\]

This is a function of $\mathbf{x}$ that gives the mean of $Y_i$ in the subpopulation where $\mathbf{X}_i$ equals $\mathbf{x}$.

\subsection{2.2.1 The CEF with a Single Binary Variable}

The simplest case has $X_i \in \{0, 1\}$:

\[
E[Y_i \mid X_i = 1] = \text{mean outcome among treated units}
\]
\[
E[Y_i \mid X_i = 0] = \text{mean outcome among untreated units}
\]

These are the two conditional means we discussed in Chapter 1. The CEF here takes only two values.

\subsection{2.2.2 The CEF with a Continuous Variable}

When $X_i$ is continuous, $E[Y_i \mid X_i = x]$ traces a potentially complex curve. For the education-wage example, it might look roughly linear but could be non-linear (e.g., the return to the 12th year of schooling may differ from the return to the 16th).

\subsection{2.2.3 The Multivariate CEF}

In practice, we condition on multiple variables $\mathbf{X}_i = (X_{1i}, X_{2i}, \ldots, X_{ki})^T$:

\[
E[Y_i \mid \mathbf{X}_i = \mathbf{x}] = E\left[Y_i \mid X_{1i} = x_1, X_{2i} = x_2, \ldots, X_{ki} = x_k\right]
\]

This is a function from $\mathbb{R}^k$ to $\mathbb{R}$. It answers the question: among all individuals with \textit{exactly} these values of all $k$ covariates, what is the average value of $Y$?

\begin{figure}[htbp]
  \centering
  \begin{tikzpicture}
    \begin{axis}[
      width=10cm, height=7cm,
      xlabel={$X$ (e.g., years of schooling)},
      ylabel={$Y$ (e.g., log wages)},
      xmin=0, xmax=6,
      ymin=0, ymax=5,
      axis lines=left,
      tick style={draw=none},
      xtick={1,2,3,4,5},
      xticklabels={10,12,14,16,18},
      ytick={1,2,3,4},
      yticklabels={2.0,2.5,3.0,3.5},
      legend style={at={(0.95,0.15)}, anchor=south east, font=\small},
      grid=none,
    ]
      % Scatter points (simulated data)
      \addplot[only marks, mark=*, mark size=1.5pt, color=gray!60, opacity=0.7]
        coordinates {
          (1,0.8)(1,1.2)(1,1.5)(1,0.5)(1,1.8)
          (2,1.3)(2,1.8)(2,2.1)(2,1.0)(2,2.5)(2,1.6)
          (3,2.0)(3,2.6)(3,2.2)(3,3.1)(3,1.8)(3,2.8)
          (4,2.5)(4,3.2)(4,2.8)(4,3.5)(4,2.1)(4,3.8)
          (5,3.2)(5,3.8)(5,3.5)(5,4.2)(5,2.9)(5,4.5)
        };

      % CEF (step function)
      \addplot[thick, econblue, const plot]
        coordinates {(0.7,1.16)(1.7,1.16)(1.7,1.88)(2.7,1.88)(2.7,2.5)(3.7,2.5)(3.7,3.18)(4.7,3.18)(4.7,3.82)(5.7,3.82)};

      % OLS regression line
      \addplot[thick, red!70!black, dashed, domain=0.7:5.7]
        {0.5 + 0.666*x};

      \legend{Data, CEF $E[Y|X=x]$, OLS regression line}
    \end{axis}
  \end{tikzpicture}
  \caption{The Conditional Expectation Function (CEF) and the OLS regression line. The CEF $E[Y_i|X_i=x]$ (solid blue step function) gives the population mean of $Y$ at each value of $X$. The OLS regression line (dashed red) is the best linear approximation to the CEF in the minimum mean-squared-error sense.}
  \label{fig:ch02_cef_vs_ols}
\end{figure}

\medskip\hrule\medskip

\section{2.3 The Law of Iterated Expectations}

The Law of Iterated Expectations (LIE) is one of the most useful identities in probability and statistics. We will use it repeatedly throughout this course.

\textbf{Theorem 2.1 (Law of Iterated Expectations):}

\[
E[Y_i] = E\left[E[Y_i \mid \mathbf{X}_i]\right]
\]

\textit{Interpretation:} You can compute the unconditional expectation of $Y_i$ in two steps: first compute $E[Y_i \mid \mathbf{X}_i = \mathbf{x}]$ for each value $\mathbf{x}$, then average this over the distribution of $\mathbf{X}_i$.

\textbf{Proof (continuous case):} Let $g(\mathbf{x}) = E[Y_i \mid \mathbf{X}_i = \mathbf{x}]$ and let $f_X(\mathbf{x})$ be the marginal density of $\mathbf{X}_i$. Then:

\[
E\left[E[Y_i \mid \mathbf{X}_i]\right] = \int g(\mathbf{x}) f_X(\mathbf{x})\, d\mathbf{x} = \int E[Y_i \mid \mathbf{X}_i = \mathbf{x}] f_X(\mathbf{x})\, d\mathbf{x} = E[Y_i] \quad \square
\]

\subsection{2.3.1 Practical Application of the LIE}

\textbf{Example 2.1:} Suppose we want to compute average wages in an economy where workers are either male (M) or female (F). The LIE tells us:

\[
E[\text{Wage}] = E[\text{Wage} \mid \text{Male}] \cdot P(\text{Male}) + E[\text{Wage} \mid \text{Female}] \cdot P(\text{Female})
\]

This is intuitive: the average wage is a probability-weighted average of average wages within each group.

\textbf{Corollary 2.1:} $E[h(\mathbf{X}_i) \epsilon_i] = 0$ for any function $h$, where $\epsilon_i = Y_i - E[Y_i \mid \mathbf{X}_i]$.

\textit{Proof:} By the LIE, $E[h(\mathbf{X}_i) \epsilon_i] = E\left[E[h(\mathbf{X}_i) \epsilon_i \mid \mathbf{X}_i]\right] = E\left[h(\mathbf{X}_i) E[\epsilon_i \mid \mathbf{X}_i]\right] = 0$, since $E[\epsilon_i \mid \mathbf{X}_i] = 0$ by construction. $\square$

\begin{figure}[htbp]
  \centering
  \begin{tikzpicture}[scale=1.1,
    box/.style={draw=econblue, fill=econblue!8, thick, rounded corners, minimum width=3.5cm, minimum height=1cm, align=center, font=\small}
  ]
    \node[box] (outer) at (0, 0) {$E[Y_i]$\\(unconditional mean)};
    \node[box] (middle) at (0, -2.2) {$E\bigl[E[Y_i|X_i]\bigr]$\\(average of conditional means)};
    \node[box] (inner) at (0, -4.4) {$E[Y_i|X_i=x]$\\(conditional mean at $x$)};

    \draw[-{Stealth}, thick, econblue] (outer) -- (middle)
      node[midway, right, font=\small] {$= $ (LIE)};
    \draw[-{Stealth}, thick, econblue] (middle) -- (inner)
      node[midway, right, font=\small] {average over $X$};

    \node[font=\footnotesize, gray, align=center] at (4.5, -2.2)
      {$E[Y_i] = E\bigl[E[Y_i | X_i]\bigr]$\\``the expectation of\\a conditional expectation\\equals the unconditional expectation''};
  \end{tikzpicture}
  \caption{The Law of Iterated Expectations (LIE). The unconditional mean equals the expectation of the conditional mean function taken over the distribution of $X$.}
  \label{fig:ch02_lie}
\end{figure}

\medskip\hrule\medskip

\section{2.4 The CEF Decomposition and Its Properties}

The CEF defines a natural \textbf{decomposition} of any random variable into a systematic part and a residual part.

\subsection{2.4.1 The CEF Decomposition}

\textbf{Proposition 2.1 (CEF Decomposition):} Any random variable $Y_i$ can be written as:

\[
Y_i = \underbrace{E[Y_i \mid \mathbf{X}_i]}_{\text{Systematic part}} + \underbrace{\epsilon_i}_{\text{Residual part}}
\]

where the residual $\epsilon_i = Y_i - E[Y_i \mid \mathbf{X}_i]$ satisfies:

\begin{enumerate}
\item \textbf{Mean Independence:} $E[\epsilon_i \mid \mathbf{X}_i] = 0$ for all values of $\mathbf{X}_i$.
\item \textbf{Zero Unconditional Mean:} $E[\epsilon_i] = 0$ (by the LIE).
\item \textbf{Uncorrelated with Any Function of $\mathbf{X}_i$:} $E[h(\mathbf{X}_i) \epsilon_i] = 0$ for any function $h$.
\end{enumerate}

\textit{Proof:} Property (1) follows directly from the definition: $E[\epsilon_i \mid \mathbf{X}_i] = E[Y_i \mid \mathbf{X}_i] - E[Y_i \mid \mathbf{X}_i] = 0$. Properties (2) and (3) follow from (1) and the LIE. $\square$

\begin{keybox}{This decomposition is fundamental.** It says that any variable can be split into a part predictable from $\mathbf{X}_i$ (the CEF) and a part that is, in the strongest possible sense, unpredictable given $\mathbf{X}_i$ (the residual $\epsilon_i$).}
\end{keybox}

\subsection{2.4.2 The CEF is the Best Predictor}

\textbf{Theorem 2.2 (CEF as Best Predictor):} Among \textit{all} functions $m(\mathbf{X}_i)$, the CEF $E[Y_i \mid \mathbf{X}_i]$ minimizes the \textbf{Mean Squared Prediction Error (MSPE)}:

\[
E[Y_i \mid \mathbf{X}_i] = \underset{m(\cdot)}{\operatorname{argmin}} \ E\!\left[(Y_i - m(\mathbf{X}_i))^2\right]
\]

\textit{Proof:} For any function $m$, write:

\[
E\left[(Y_i - m(\mathbf{X}_i))^2\right] = E\left[(Y_i - E[Y_i \mid \mathbf{X}_i] + E[Y_i \mid \mathbf{X}_i] - m(\mathbf{X}_i))^2\right]
\]

\[
= E[\epsilon_i^2] + E\left[(E[Y_i \mid \mathbf{X}_i] - m(\mathbf{X}_i))^2\right] + 2E\!\left[\epsilon_i (E[Y_i \mid \mathbf{X}_i] - m(\mathbf{X}_i))\right]
\]

The cross term is zero by Corollary 2.1. The second term is non-negative and equals zero only when $m(\mathbf{X}_i) = E[Y_i \mid \mathbf{X}_i]$. $\square$

\textit{Intuition:} $E[Y_i]$ is the best constant predictor of $Y_i$. The CEF does better by conditioning on available information: by the LIE, $E[E[Y_i \mid \mathbf{X}_i]] = E[Y_i]$, so the CEF is as good on average and better in specific strata.

\subsection{2.4.3 The ANOVA Theorem}

\textbf{Theorem 2.3 (Variance Decomposition / ANOVA):}

\[
\underbrace{V(Y_i)}_{\text{Total variance}} = \underbrace{V(E[Y_i \mid \mathbf{X}_i])}_{\text{Explained variance}} + \underbrace{E[V(Y_i \mid \mathbf{X}_i)]}_{\text{Unexplained variance}}
\]

\textit{Interpretation:} The total variation in $Y_i$ decomposes into:
\begin{itemize}
\item The variation in the conditional mean (what $\mathbf{X}$ can "explain")
\item The average within-group variation (what remains after conditioning on $\mathbf{X}$)
\end{itemize}

This is the population analogue of the sample $R^2$ decomposition $\text{TSS} = \text{ESS} + \text{SSR}$.

\textbf{Example 2.2 (ANOVA in Wage Data):} In the CPS data, approximately 25--30\% of the variation in log wages is explained by education, age, gender, and occupation. The remaining 70--75\% is within-group variation: workers with the same observable characteristics still have very different wages.

\medskip\hrule\medskip

\section{2.5 Why Linear Regression?}

The CEF can take any shape --- it need not be linear. So why do economists almost always estimate \textit{linear} regressions? This section provides two justifications.

\subsection{2.5.1 Justification 1: The CEF is Linear}

If the true population CEF is linear:

\[
E[Y_i \mid \mathbf{X}_i] = \mathbf{X}_i^T \boldsymbol{\beta}^*
\]

then by the CEF decomposition theorem:

\[
E\left[\mathbf{X}_i(Y_i - \mathbf{X}_i^T \boldsymbol{\beta}^*)\right] = \mathbf{0}
\]

Solving for $\boldsymbol{\beta}^*$:

\[
\boldsymbol{\beta}^* = E[\mathbf{X}_i \mathbf{X}_i^T]^{-1} E[\mathbf{X}_i Y_i] \equiv \boldsymbol{\beta}^{OLS}
\]

So when the CEF is linear, the OLS population parameter $\boldsymbol{\beta}^{OLS}$ equals the true parameter $\boldsymbol{\beta}^*$ of the CEF. The linear regression model is correctly specified.

\textbf{When is the CEF exactly linear?} The CEF is linear in $X$ when $(Y, X)$ follow a jointly normal distribution (because $E[Y \mid X] = \mu_Y + \rho \frac{\sigma_Y}{\sigma_X}(X - \mu_X)$ for the bivariate normal). But this is a very strong assumption.

\subsection{2.5.2 Justification 2: OLS Provides the Best Linear Approximation}

What if the CEF is \textit{not} linear? Does OLS still make sense? The answer is yes.

\textbf{Theorem 2.4 (OLS as Best Linear Approximation):} Even when $E[Y_i \mid \mathbf{X}_i]$ is not linear, $\boldsymbol{\beta}^{OLS}$ minimizes the mean squared error of linear prediction:

\[
\boldsymbol{\beta}^{OLS} = \underset{\mathbf{b}}{\operatorname{argmin}} \ E\!\left[(Y_i - \mathbf{X}_i^T \mathbf{b})^2\right]
\]

Moreover, the linear projection $\mathbf{X}_i^T \boldsymbol{\beta}^{OLS}$ is the \textbf{best linear approximation to the CEF}:

\[
\boldsymbol{\beta}^{OLS} = \underset{\mathbf{b}}{\operatorname{argmin}} \ E\!\left[(E[Y_i \mid \mathbf{X}_i] - \mathbf{X}_i^T \mathbf{b})^2\right]
\]

\textit{Proof of the second result:} Taking the first-order condition:

\[
E\!\left[\mathbf{X}_i(E[Y_i \mid \mathbf{X}_i] - \mathbf{X}_i^T \boldsymbol{\beta})\right] = \mathbf{0}
\]

\[
\Rightarrow \boldsymbol{\beta} = E[\mathbf{X}_i \mathbf{X}_i^T]^{-1} E[\mathbf{X}_i E[Y_i \mid \mathbf{X}_i]]
\]

By the CEF decomposition, $E[\mathbf{X}_i E[Y_i \mid \mathbf{X}_i]] = E[\mathbf{X}_i Y_i] - E[\mathbf{X}_i \epsilon_i] = E[\mathbf{X}_i Y_i]$. 

Therefore: $\boldsymbol{\beta} = E[\mathbf{X}_i \mathbf{X}_i^T]^{-1} E[\mathbf{X}_i Y_i] = \boldsymbol{\beta}^{OLS}$. $\square$

\begin{keybox}{Key Takeaway:** OLS provides the best linear fit to the data, regardless of whether the CEF is linear or not. When the CEF is not linear, the OLS regression is the best \textit{linear} approximation to the CEF --- it fits a straight line through a potentially curved cloud as well as any straight line can.}
\end{keybox}

\subsection{2.5.3 Visualizing the Linear Approximation}

Consider the education-wage CEF. The true CEF may curve (with higher returns to college years than to early secondary years). But the OLS regression line through the scatter plot is the straight line that minimizes the sum of squared vertical distances to all data points. This line is the best linear summary of the relationship.

For many purposes --- especially for computing conditional means at typical values of the covariates --- the linear approximation is excellent. For prediction at extreme values or for studying non-linearities, we might prefer more flexible methods.

\medskip\hrule\medskip

\section{2.6 OLS as the Best Linear Approximation to the CEF}

Let us be more explicit about the \textbf{population linear projection model}. 

\textbf{Definition 2.2 (Linear Projection / Population OLS):} The linear projection of $Y_i$ on $\mathbf{X}_i$ is:

\[
L(Y_i \mid \mathbf{X}_i) = \mathbf{X}_i^T \boldsymbol{\beta}^{OLS}
\]

where $\boldsymbol{\beta}^{OLS}$ solves:

\[
\boldsymbol{\beta}^{OLS} = \left(E[\mathbf{X}_i \mathbf{X}_i^T]\right)^{-1} E[\mathbf{X}_i Y_i]
\]

This is the population parameter vector that OLS estimates from a sample. The sample OLS estimator (Chapter 3) converges to this population quantity as sample size grows.

\textbf{Interpreting the Components:} Consider the simple bivariate case with $\mathbf{X}_i = (1, X_i)^T$ and $\boldsymbol{\beta} = (\alpha, \beta)^T$. Then:

\[
\beta^{OLS} = \frac{\text{Cov}(X_i, Y_i)}{V(X_i)}
\]

This is the \textbf{population regression slope}: the ratio of the covariance of $X$ and $Y$ to the variance of $X$. It equals the change in the conditional mean of $Y$ for a one-unit change in $X$, under the linear approximation.

\subsection{2.6.1 The Linear Regression Equation}

Given $\boldsymbol{\beta}^{OLS}$, we can always write:

\[
Y_i = \mathbf{X}_i^T \boldsymbol{\beta}^{OLS} + \underbrace{e_i}_{\text{Population OLS residual}}
\]

where $e_i = Y_i - \mathbf{X}_i^T \boldsymbol{\beta}^{OLS}$ is the \textbf{population OLS residual} (not the CEF residual $\epsilon_i$, though they coincide when the CEF is linear). By construction:

\[
E[\mathbf{X}_i e_i] = \mathbf{0}
\]

This \textbf{zero covariance condition} (or \textbf{orthogonality condition}) is the key property of the OLS estimator --- it is weaker than strict mean independence $E[e_i \mid \mathbf{X}_i] = 0$.

\subsection{2.6.2 When Does Regression Identify Causal Effects?}

Regression identifies the CEF (up to a linear approximation). But the CEF is not necessarily a causal relationship. Whether the regression coefficient $\beta$ has a causal interpretation depends on \textit{why} $E[Y_i \mid X_i = x]$ varies with $x$.

Three distinct reasons for $E[Y_i \mid X_i = x]$ to vary with $x$:

\begin{enumerate}
\item \textbf{Causality:} changing $X$ for a given individual changes their $Y$.
\item \textbf{Selection:} people with higher $X$ would have had higher $Y$ even without the higher $X$ (like the wage-education example with ability selection).
\item \textbf{Reverse causality:} $Y$ causes $X$ (e.g., healthier people exercise more).
\end{enumerate}

OLS recovers the CEF regardless of the reason. The \textbf{CEF is causal} only if $X$ is randomly assigned or if we can credibly argue that, conditional on $\mathbf{X}$, the treatment variable of interest is "as good as randomly assigned." This is the identification problem we discussed in Chapter 1.

\medskip\hrule\medskip

\section{2.7 Empirical Illustration: Education and Wages}

Let us apply these concepts to a classic economic question: how does schooling affect wages?

\subsection{2.7.1 Data and Setting}

We use data from the Current Population Survey (CPS), a representative survey of the US labour market. The outcome variable $Y_i$ is the log of weekly earnings; the regressor $X_i$ is years of completed schooling.

\subsection{2.7.2 The Wage-Schooling CEF}

When we calculate $E[\log(\text{wage}) \mid \text{schooling} = s]$ for each value of $s = 0, 1, \ldots, 20$, we trace out the wage-schooling CEF. Key features:

\begin{itemize}
\item The CEF is approximately \textbf{linear} over most of the schooling range.
\item There is a notable \textbf{kink} at 12 years (high school completion) and at 16 years (college completion) --- the "sheepskin" effects.
\item The overall slope is approximately \textbf{0.08--0.11} in OLS regressions, implying each additional year of schooling is associated with an 8--11\% increase in wages.
\end{itemize}

\subsection{2.7.3 The Linear Regression Approximation}

The OLS regression line:

\[
\widehat{\log(\text{wage})_i} = \hat{\alpha} + \hat{\beta} \times \text{schooling}_i
\]

provides the best linear fit to the wage-schooling scatter. The slope $\hat{\beta} \approx 0.10$ says: on average, each additional year of schooling is associated with a 10\% increase in weekly wages.

\textbf{Is this causal?} Not necessarily. People who stay in school longer differ in many ways --- cognitive ability, family background, patience --- from those who leave school earlier. The OLS regression does not distinguish between the causal effect of schooling and these selection differences. We return to this in Chapters 6 (omitted variable bias) and 8 (instrumental variables).

\medskip\hrule\medskip

\section{Chapter Summary}

\begin{itemize}
\item The \textbf{Conditional Expectation Function (CEF)} $E[Y_i \mid \mathbf{X}_i]$ is the best predictor of $Y_i$ given $\mathbf{X}_i$ in the mean-squared-error sense.
\item The \textbf{CEF Decomposition} writes $Y_i = E[Y_i \mid \mathbf{X}_i] + \epsilon_i$, where $\epsilon_i$ has zero mean conditional on $\mathbf{X}_i$.
\item The \textbf{Law of Iterated Expectations} relates the unconditional mean to an average of conditional means.
\item The \textbf{ANOVA Theorem} decomposes total variance into explained and unexplained components.
\item The OLS population coefficient vector $\boldsymbol{\beta}^{OLS}$ equals the CEF parameter when the CEF is linear.
\item When the CEF is non-linear, OLS still provides the \textbf{best linear approximation} to the CEF.
\item Whether the OLS regression has a \textbf{causal interpretation} depends on the research design, not the statistical method.
\end{itemize}

\medskip\hrule\medskip

\section{Key Terms}

\textbf{Conditional Expectation Function (CEF)} --- the function $E[Y_i \mid \mathbf{X}_i = \mathbf{x}]$ giving the expected value of $Y$ for each value of $\mathbf{X}$.

\textbf{Law of Iterated Expectations (LIE)} --- $E[Y_i] = E[E[Y_i \mid \mathbf{X}_i]]$.

\textbf{CEF Decomposition} --- $Y_i = E[Y_i \mid \mathbf{X}_i] + \epsilon_i$ where $E[\epsilon_i \mid \mathbf{X}_i] = 0$.

\textbf{Mean Independence} --- the condition $E[\epsilon_i \mid \mathbf{X}_i] = 0$.

\textbf{ANOVA Theorem} --- $V(Y_i) = V(E[Y_i \mid \mathbf{X}_i]) + E[V(Y_i \mid \mathbf{X}_i)]$.

\textbf{Best linear approximation} --- the linear function of $\mathbf{X}_i$ that is closest to the CEF in the mean-squared-error sense.

\textbf{Population OLS parameter} --- $\boldsymbol{\beta}^{OLS} = (E[\mathbf{X}_i \mathbf{X}_i^T])^{-1} E[\mathbf{X}_i Y_i]$.

\textbf{Linear projection} --- $L(Y_i \mid \mathbf{X}_i) = \mathbf{X}_i^T \boldsymbol{\beta}^{OLS}$.

\textbf{Orthogonality condition} --- $E[\mathbf{X}_i e_i] = \mathbf{0}$, where $e_i$ is the population OLS residual.

\medskip\hrule\medskip

\section{Exercises}

\subsection{Conceptual Questions}

\textbf{2.1} Explain in plain language what the CEF $E[Y_i \mid X_i = x]$ means. Give an example from economics where the CEF is likely non-linear.

\textbf{2.2} State the Law of Iterated Expectations. Give an example where you would use it.

\textbf{2.3} Is it possible for $E[Y_i \mid X_i = x]$ to be decreasing in $x$ even though the population regression coefficient $\beta^{OLS}$ is positive? Give an example or prove impossibility.

\textbf{2.4} Explain the difference between the CEF residual $\epsilon_i = Y_i - E[Y_i \mid \mathbf{X}_i]$ and the population OLS residual $e_i = Y_i - \mathbf{X}_i^T \boldsymbol{\beta}^{OLS}$. When are they the same?

\textbf{2.5} "OLS regression always identifies the causal effect of $X$ on $Y$." Is this statement true or false? Explain.

\subsection{Analytical Questions}

\textbf{2.6} Let $Y_i \in \{0, 1\}$ (binary outcome) and $X_i \in \{0, 1\}$ (binary treatment), with joint distribution:

\begin{center}
\begin{tabular}{@{} l l l @{}}
\toprule
\textbf{} & \textbf{$X_i = 0$} & \textbf{$X_i = 1$} \\
\midrule
$Y_i = 0$ & 0.40 & 0.10 \\
$Y_i = 1$ & 0.20 & 0.30 \\
\bottomrule
\end{tabular}
\end{center}

(a) Compute $E[Y_i]$, $E[Y_i \mid X_i = 0]$, and $E[Y_i \mid X_i = 1]$.  
(b) Compute $\beta^{OLS}$ from the bivariate regression of $Y_i$ on $X_i$ using $\beta^{OLS} = \text{Cov}(X,Y)/V(X)$.  
(c) Interpret the slope coefficient. Does it have a causal interpretation here?

\textbf{2.7} Using the ANOVA theorem, show that the population $R^2 = V(E[Y_i \mid \mathbf{X}_i]) / V(Y_i)$ lies between 0 and 1.

\textbf{2.8} Suppose $E[Y_i \mid X_i] = \alpha + \beta X_i + \gamma X_i^2$ (a quadratic CEF). You run the OLS regression of $Y_i$ on $X_i$ only (without the quadratic term). Show that the OLS estimate of the slope will generally not equal $\beta$, and derive the formula for the bias in terms of $\gamma$ and the moments of $X_i$.

\textbf{2.9} Prove Corollary 2.1: if $\epsilon_i = Y_i - E[Y_i \mid \mathbf{X}_i]$, then $E[h(\mathbf{X}_i)\epsilon_i] = 0$ for any function $h(\cdot)$.

\textbf{2.10} Consider the regression of $Y_i$ on $\mathbf{X}_i$. Show that adding a variable $Z_i$ to the regression does not change the population OLS coefficient on $X_{1i}$ if and only if $Z_i$ is uncorrelated with $X_{1i}$ after partialling out the other regressors.

\subsection{Applied Questions}

\textbf{2.11} Describe how you would construct the wage-schooling CEF from survey data. What would you do for each value of schooling from 0 to 20?

\textbf{2.12} The CEF of health status on income slopes upward: people with higher income have better health on average. List at least three reasons this positive association might hold, and classify each as: (a) causal (income causes health), (b) reverse causality (health causes income), (c) confounding.

\textbf{2.13} In your own words, explain why the fact that "OLS provides the best linear approximation to the CEF" is a useful result even when the CEF is non-linear.

\medskip\hrule\medskip

\section{R Lab 2: Estimating the CEF and OLS Regression}

\subsection{Objectives}
\begin{itemize}
\item Compute the empirical CEF from data
\item Compare the CEF to OLS regression
\item Visualize the relationship between a continuous variable and an outcome
\end{itemize}

\subsection{Setup}

\begin{lstlisting}
library(tidyverse)
# install.packages("quantreg")
# install.packages("AER")
library(AER)  # contains CPS1988 data

# Load CPS data on wages and education
data("CPS1988")
head(CPS1988)
\end{lstlisting}

\subsection{Exercise 2.1: Exploring the Data}

\begin{lstlisting}
# Summary statistics
summary(CPS1988[, c("wage", "education", "experience", "ethnicity")])

# Log wages
CPS1988 <- CPS1988 %>%
  mutate(log_wage = log(wage))

# Distribution of education
ggplot(CPS1988, aes(x = education)) +
  geom_histogram(bins = 16, fill = "steelblue", color = "white") +
  labs(title = "Distribution of Years of Schooling", x = "Years of Education", y = "Count") +
  theme_minimal()
\end{lstlisting}

\subsection{Exercise 2.2: Estimating the Empirical CEF}

\begin{lstlisting}
# Compute the empirical CEF: mean log wage for each education level
cef_data <- CPS1988 %>%
  group_by(education) %>%
  summarise(
    mean_log_wage  = mean(log_wage),
    median_log_wage = median(log_wage),
    n = n(),
    se_log_wage = sd(log_wage) / sqrt(n)
  )

print(cef_data)

# Plot the empirical CEF with confidence intervals
ggplot(cef_data, aes(x = education, y = mean_log_wage)) +
  geom_point(aes(size = n), color = "steelblue") +
  geom_line(color = "steelblue") +
  geom_ribbon(aes(ymin = mean_log_wage - 1.96 * se_log_wage,
                  ymax = mean_log_wage + 1.96 * se_log_wage),
              alpha = 0.2, fill = "steelblue") +
  labs(
    title = "Empirical CEF: Mean Log Wage by Years of Education",
    subtitle = "The CEF shows a generally increasing, approximately linear relationship",
    x = "Years of Education",
    y = "Mean Log Wage",
    size = "Sample Size"
  ) +
  theme_minimal()
\end{lstlisting}

\subsection{Exercise 2.3: OLS Regression vs. the CEF}

\begin{lstlisting}
# Run OLS regression
ols_model <- lm(log_wage ~ education, data = CPS1988)
summary(ols_model)

# Add OLS fitted line to the CEF plot
ggplot(cef_data, aes(x = education, y = mean_log_wage)) +
  geom_point(color = "steelblue", size = 3) +
  geom_line(color = "steelblue", linetype = "dashed") +
  geom_smooth(
    data = CPS1988,
    aes(x = education, y = log_wage),
    method = "lm",
    se = TRUE,
    color = "red"
  ) +
  labs(
    title = "Empirical CEF vs. OLS Regression Line",
    subtitle = "Blue: Empirical CEF; Red: OLS best linear fit",
    x = "Years of Education",
    y = "Log Wage"
  ) +
  theme_minimal()
\end{lstlisting}

\textbf{Questions:}
\begin{enumerate}
\item Is the empirical CEF approximately linear?
\item Where does it deviate most from linearity? (Hint: look at 12 and 16 years of schooling.)
\item How does the slope of the OLS regression compare to the slope of the CEF in the linear range?
\end{enumerate}

\subsection{Exercise 2.4: Residuals and the ANOVA Decomposition}

\begin{lstlisting}
# Compute OLS residuals
CPS1988 <- CPS1988 %>%
  mutate(
    fitted_ols = predict(ols_model),
    resid_ols  = residuals(ols_model)
  )

# ANOVA decomposition
TSS <- var(CPS1988$log_wage)
ESS <- var(CPS1988$fitted_ols)
SSR_mean <- mean(CPS1988$resid_ols^2)

cat("Total Variance (TSS equivalent):", round(TSS, 4), "\n")
cat("Explained Variance (ESS):", round(ESS, 4), "\n")
cat("Average Residual Variance:", round(SSR_mean, 4), "\n")
cat("R-squared:", round(summary(ols_model)$r.squared, 4), "\n")
cat("Check (ESS/TSS):", round(ESS/TSS, 4), "\n")

# Does E[X * resid] = 0?
cat("Cov(education, residuals):", round(cov(CPS1988$education, CPS1988$resid_ols), 6), "\n")
\end{lstlisting}

\subsection{Exercise 2.5: Multiple Regression CEF}

\begin{lstlisting}
# Multiple regression: education + experience
ols_multi <- lm(log_wage ~ education + experience + I(experience^2), data = CPS1988)
summary(ols_multi)

# Partial effect of education (controlling for experience)
# Using the FWL theorem (preview):
# Regress education on experience, get residuals
educ_resid <- residuals(lm(education ~ experience + I(experience^2), data = CPS1988))
wage_resid  <- residuals(lm(log_wage  ~ experience + I(experience^2), data = CPS1988))
fwl_model   <- lm(wage_resid ~ educ_resid)
cat("FWL estimate of education coefficient:", round(coef(fwl_model)[2], 4), "\n")
cat("Multiple regression estimate:", round(coef(ols_multi)["education"], 4), "\n")
\end{lstlisting}

\textbf{Interpretation:} The two estimates should be identical, illustrating the Frisch-Waugh-Lovell Theorem (covered in Chapter 6).

\medskip\hrule\medskip

\textit{End of Chapter 2}

\medskip\hrule\medskip

\textbf{References}

Angrist, J. D., and Pischke, J. (2009). \textit{Mostly Harmless Econometrics}. Princeton University Press. Chapter 3, Sections 3.1.1--3.1.4.

Davidson, R., and MacKinnon, J. G. (2004). \textit{Econometric Theory and Methods}. Oxford University Press. Chapter 1.

Wooldridge, J. M. (2019). \textit{Introductory Econometrics: A Modern Approach} (7th ed.). Cengage. Chapters 2--3.

\chapter{Chapter 3: Ordinary Least Squares --- Algebra and Geometry}

\medskip\hrule\medskip

\section{Chapter Overview}

This chapter develops the OLS estimator from the ground up. We begin with the multiple regression model written in matrix notation, derive the OLS estimator through two routes (method-of-moments and minimization), and then provide a geometric interpretation using linear algebra. This geometric perspective is not merely elegant --- it reveals deep structural properties of OLS that are impossible to see from the algebraic derivation alone. We end with the Frisch-Waugh-Lovell insight, goodness-of-fit measures, and invariance properties of OLS.

Throughout this chapter, we work at the \textbf{sample} level: we have $n$ observations and we are fitting an estimator to data, not analyzing a population parameter.

\medskip\hrule\medskip

\section{Learning Objectives}

After completing this chapter, students will be able to:

\begin{enumerate}
\item Write the multiple regression model in \textbf{matrix notation}.
\item Derive the \textbf{OLS estimator} using the method-of-moments and the least-squares criterion.
\item State the conditions under which $(\mathbf{X}^T\mathbf{X})$ is invertible (full rank condition).
\item Interpret OLS geometrically as an \textbf{orthogonal projection} onto the column space of $\mathbf{X}$.
\item Define the \textbf{projection matrices} $\mathbf{P_X}$ and $\mathbf{M_X}$ and state their key properties.
\item State and apply the \textbf{Pythagorean decomposition} of sum of squares.
\item Compute and interpret \textbf{$R^2$} (centered and uncentered).
\item State and apply the \textbf{Frisch-Waugh-Lovell (FWL) Theorem}.
\item Discuss how OLS coefficients change with linear transformations of regressors.
\end{enumerate}

\medskip\hrule\medskip

\section{3.1 The Multiple Regression Model}

We observe $n$ units (individuals, firms, countries), each described by:
\begin{itemize}
\item an \textbf{outcome variable} $y_i \in \mathbb{R}$
\item a $k \times 1$ vector of \textbf{regressors} (including a constant in the first position): $\mathbf{x}_i = (1, x_{i2}, \ldots, x_{ik})^T$
\end{itemize}

The multiple regression equation for unit $i$ is:

\[
y_i = \beta_1 + \beta_2 x_{i2} + \cdots + \beta_k x_{ik} + u_i = \mathbf{x}_i^T \boldsymbol{\beta} + u_i
\]

where $\boldsymbol{\beta} = (\beta_1, \beta_2, \ldots, \beta_k)^T$ is the $k \times 1$ vector of \textbf{coefficients} and $u_i$ is the \textbf{error term}.

\subsection{3.1.1 Matrix Notation}

Stack all $n$ observations:

\[
\underset{n \times 1}{\mathbf{y}} = \underset{n \times k}{\mathbf{X}} \underset{k \times 1}{\boldsymbol{\beta}} + \underset{n \times 1}{\mathbf{u}}
\]

where:

\[
\mathbf{y} = \begin{pmatrix} y_1 \\ y_2 \\ \vdots \\ y_n \end{pmatrix}, \quad
\mathbf{X} = \begin{pmatrix} \mathbf{x}_1^T \\ \mathbf{x}_2^T \\ \vdots \\ \mathbf{x}_n^T \end{pmatrix} = \begin{pmatrix} 1 & x_{12} & \cdots & x_{1k} \\ 1 & x_{22} & \cdots & x_{2k} \\ \vdots & \vdots & \ddots & \vdots \\ 1 & x_{n2} & \cdots & x_{nk} \end{pmatrix}, \quad
\boldsymbol{\beta} = \begin{pmatrix} \beta_1 \\ \beta_2 \\ \vdots \\ \beta_k \end{pmatrix}, \quad
\mathbf{u} = \begin{pmatrix} u_1 \\ u_2 \\ \vdots \\ u_n \end{pmatrix}
\]

The matrix $\mathbf{X}$ is called the \textbf{design matrix}. Each row is an observation; each column is a regressor.

\begin{keybox}{Convention:** We use bold lowercase letters for vectors ($\mathbf{y}$, $\mathbf{u}$) and bold uppercase letters for matrices ($\mathbf{X}$, $\mathbf{P}$, $\mathbf{M}$). Scalars are unbolded.}
\end{keybox}

\textbf{Example 3.1 (Wage Regression):} In a regression of log wage on years of education and years of experience:

\[
\log(\text{wage}_i) = \beta_1 + \beta_2 \cdot \text{educ}_i + \beta_3 \cdot \text{exp}_i + u_i
\]

Here $n$ is the number of workers, $k = 3$, and the design matrix $\mathbf{X}$ has three columns: a column of ones, a column of education values, and a column of experience values.

\medskip\hrule\medskip

\section{3.2 The OLS Estimator: Derivation via Method of Moments}

The \textbf{Method of Moments (MM)} estimator is derived by finding $\hat{\boldsymbol{\beta}}$ such that the \textbf{sample moment conditions} are satisfied.

From the population orthogonality condition $E[\mathbf{x}_i u_i] = \mathbf{0}$, the sample analogue is:

\[
\frac{1}{n} \sum_{i=1}^n \mathbf{x}_i (y_i - \mathbf{x}_i^T \hat{\boldsymbol{\beta}}) = \mathbf{0}
\]

In matrix notation:

\[
\mathbf{X}^T (\mathbf{y} - \mathbf{X}\hat{\boldsymbol{\beta}}) = \mathbf{0}
\]

These are the \textbf{$k$ Normal Equations} (one for each regressor):

\[
\mathbf{X}^T \mathbf{X} \hat{\boldsymbol{\beta}} = \mathbf{X}^T \mathbf{y}
\]

If $\mathbf{X}^T \mathbf{X}$ is \textbf{invertible} (i.e., $\mathbf{X}$ has full column rank $k$), we can solve:

\[
\boxed{\hat{\boldsymbol{\beta}}^{OLS} = (\mathbf{X}^T \mathbf{X})^{-1} \mathbf{X}^T \mathbf{y}}
\]

This is the \textbf{OLS estimator} --- the central formula of classical econometrics.

\subsection{3.2.1 The Full Rank Condition}

The matrix $\mathbf{X}$ has \textbf{full column rank} (rank $k$) if and only if no column of $\mathbf{X}$ can be written as a linear combination of the other columns. Violations include:

\begin{enumerate}
\item \textbf{Perfect multicollinearity:} two regressors are proportional or one is an exact linear combination of others.
\item \textbf{Dummy variable trap:} including a dummy for every category while also including an intercept.
\item \textbf{Too many variables relative to observations:} $k > n$.
\end{enumerate}

\textbf{Example 3.2:} If we include both a male dummy and a female dummy (along with an intercept), then (male dummy) + (female dummy) = (column of ones). The three columns are linearly dependent; $\mathbf{X}^T \mathbf{X}$ is singular and the OLS estimator does not exist.

\medskip\hrule\medskip

\section{3.3 OLS Derivation via Minimization}

Alternatively, define the \textbf{sum of squared residuals (SSR)}:

\[
SSR(\hat{\boldsymbol{\beta}}) = (\mathbf{y} - \mathbf{X}\hat{\boldsymbol{\beta}})^T (\mathbf{y} - \mathbf{X}\hat{\boldsymbol{\beta}}) = \sum_{i=1}^n (y_i - \mathbf{x}_i^T \hat{\boldsymbol{\beta}})^2 = \|\mathbf{y} - \mathbf{X}\hat{\boldsymbol{\beta}}\|^2
\]

The \textbf{Ordinary Least Squares} estimator minimizes the SSR:

\[
\hat{\boldsymbol{\beta}}^{OLS} = \underset{\hat{\boldsymbol{\beta}}}{\operatorname{argmin}} \ \|\mathbf{y} - \mathbf{X}\hat{\boldsymbol{\beta}}\|^2
\]

\textbf{Derivation:} Expand the objective function:

\[
SSR = \mathbf{y}^T\mathbf{y} - 2\hat{\boldsymbol{\beta}}^T \mathbf{X}^T\mathbf{y} + \hat{\boldsymbol{\beta}}^T \mathbf{X}^T\mathbf{X} \hat{\boldsymbol{\beta}}
\]

Taking the derivative with respect to $\hat{\boldsymbol{\beta}}$ and setting to zero:

\[
\frac{\partial SSR}{\partial \hat{\boldsymbol{\beta}}} = -2\mathbf{X}^T\mathbf{y} + 2\mathbf{X}^T\mathbf{X}\hat{\boldsymbol{\beta}} = \mathbf{0}
\]

This gives exactly the normal equations. The second-order condition requires $\mathbf{X}^T\mathbf{X}$ to be positive definite, which holds when $\mathbf{X}$ has full rank. $\square$

\begin{keybox}{The equivalence of MM and OLS** follows because the moment conditions and the first-order conditions of the least-squares problem are identical. This is a special property of linear models.}
\end{keybox}

\medskip\hrule\medskip

\section{3.4 Linear Regression as Orthogonal Projection}

The geometric interpretation of OLS is the most insightful way to understand what the estimator is doing. This section requires thinking of vectors and matrices in $\mathbb{R}^n$.

\subsection{3.4.1 Vectors in $\mathbb{R}^n$}

Think of $\mathbf{y} = (y_1, \ldots, y_n)^T$ as a single point (or arrow from the origin) in $n$-dimensional Euclidean space $E^n$. Similarly, each column of $\mathbf{X}$ is a vector in $E^n$.

The \textbf{scalar product} (inner product) of two vectors $\mathbf{a}, \mathbf{b} \in E^n$ is:

\[
\langle \mathbf{a}, \mathbf{b} \rangle = \mathbf{a}^T \mathbf{b} = \sum_{i=1}^n a_i b_i
\]

The \textbf{length} (norm) of a vector is:

\[
\|\mathbf{a}\| = \sqrt{\langle \mathbf{a}, \mathbf{a} \rangle} = \sqrt{\sum_{i=1}^n a_i^2}
\]

Two vectors are \textbf{orthogonal} if $\langle \mathbf{a}, \mathbf{b} \rangle = 0$. The \textbf{Cauchy-Schwarz inequality} states: $|\langle \mathbf{a}, \mathbf{b} \rangle| \leq \|\mathbf{a}\| \|\mathbf{b}\|$.

\subsection{3.4.2 The Column Space of X}

The \textbf{column space} of $\mathbf{X}$, denoted $\mathcal{S}(\mathbf{X})$, is the set of all vectors that can be written as a linear combination of the columns of $\mathbf{X}$:

\[
\mathcal{S}(\mathbf{X}) = \left\{\mathbf{z} \in E^n : \mathbf{z} = \mathbf{X}\mathbf{b} \text{ for some } \mathbf{b} \in \mathbb{R}^k\right\}
\]

When $\mathbf{X}$ has rank $k$, $\mathcal{S}(\mathbf{X})$ is a $k$-dimensional \textbf{subspace} of $E^n$.

\textbf{Key insight:} The fitted values $\mathbf{X}\hat{\boldsymbol{\beta}}$ lie in $\mathcal{S}(\mathbf{X})$ for \textit{any} choice of $\hat{\boldsymbol{\beta}}$.

\subsection{3.4.3 OLS as Closest Point in the Column Space}

The OLS problem asks: \textbf{which point in $\mathcal{S}(\mathbf{X})$ is closest to $\mathbf{y}$?}

Formally: find $\hat{\mathbf{y}} = \mathbf{X}\hat{\boldsymbol{\beta}} \in \mathcal{S}(\mathbf{X})$ that minimizes $\|\mathbf{y} - \hat{\mathbf{y}}\|^2$.

\textbf{The answer is the orthogonal projection of $\mathbf{y}$ onto $\mathcal{S}(\mathbf{X})$.}

This is the point in $\mathcal{S}(\mathbf{X})$ such that the residual vector $\hat{\mathbf{u}} = \mathbf{y} - \hat{\mathbf{y}}$ is \textbf{orthogonal} to every vector in $\mathcal{S}(\mathbf{X})$:

\[
\mathbf{X}^T \hat{\mathbf{u}} = \mathbf{0}
\]

This is exactly the normal equation, confirming the geometric interpretation.

\begin{keybox}{Figure 3.1 (Geometric Interpretation):** Imagine $n = 3$ observations and $k = 2$ regressors (including the constant). The column space $\mathcal{S}(\mathbf{X})$ is a 2-dimensional plane in 3-dimensional space. The vector $\mathbf{y}$ is a point in 3D space, generally not lying in the plane. The OLS fitted value $\hat{\mathbf{y}}$ is the point in the plane closest to $\mathbf{y}$, obtained by dropping a perpendicular from $\mathbf{y}$ onto the plane. The residual vector $\hat{\mathbf{u}} = \mathbf{y} - \hat{\mathbf{y}}$ is the perpendicular.}
\end{keybox}

\medskip\hrule\medskip

\section{3.5 Projection Matrices: $\mathbf{P_X}$ and $\mathbf{M_X}$}

\subsection{3.5.1 The Hat Matrix $\mathbf{P_X}$}

The \textbf{orthogonal projection matrix} onto $\mathcal{S}(\mathbf{X})$ is:

\[
\mathbf{P_X} = \mathbf{X}(\mathbf{X}^T \mathbf{X})^{-1} \mathbf{X}^T
\]

This is also called the \textbf{hat matrix} (because it "puts a hat on" $\mathbf{y}$):

\[
\hat{\mathbf{y}} = \mathbf{X}\hat{\boldsymbol{\beta}} = \mathbf{X}(\mathbf{X}^T\mathbf{X})^{-1}\mathbf{X}^T \mathbf{y} = \mathbf{P_X} \mathbf{y}
\]

\subsection{3.5.2 The Annihilator Matrix $\mathbf{M_X}$}

The \textbf{orthogonal projection onto the orthogonal complement} of $\mathcal{S}(\mathbf{X})$ is:

\[
\mathbf{M_X} = \mathbf{I}_n - \mathbf{P_X}
\]

This matrix produces the OLS residuals:

\[
\hat{\mathbf{u}} = \mathbf{y} - \hat{\mathbf{y}} = \mathbf{y} - \mathbf{P_X}\mathbf{y} = (\mathbf{I} - \mathbf{P_X})\mathbf{y} = \mathbf{M_X}\mathbf{y}
\]

$\mathbf{M_X}$ is called the \textbf{annihilator} because it destroys (annihilates) any vector in $\mathcal{S}(\mathbf{X})$: $\mathbf{M_X}\mathbf{X} = \mathbf{0}$.

\subsection{3.5.3 Key Properties of $\mathbf{P_X}$ and $\mathbf{M_X}$}

\textbf{Property 3.1 (Symmetry):} $\mathbf{P_X}^T = \mathbf{P_X}$ and $\mathbf{M_X}^T = \mathbf{M_X}$.

\textbf{Property 3.2 (Idempotency):} $\mathbf{P_X}^2 = \mathbf{P_X}$ and $\mathbf{M_X}^2 = \mathbf{M_X}$.

\textit{Proof:} $\mathbf{P_X}^2 = \mathbf{X}(\mathbf{X}^T\mathbf{X})^{-1}\mathbf{X}^T \cdot \mathbf{X}(\mathbf{X}^T\mathbf{X})^{-1}\mathbf{X}^T = \mathbf{X}(\mathbf{X}^T\mathbf{X})^{-1}(\mathbf{X}^T\mathbf{X})(\mathbf{X}^T\mathbf{X})^{-1}\mathbf{X}^T = \mathbf{P_X}$. $\square$

\textit{Interpretation:} Projecting a vector twice onto the same subspace returns it to where it was after the first projection. Idempotency is the algebraic expression of this geometric fact.

\textbf{Property 3.3 (Annihilation):} $\mathbf{P_X}\mathbf{M_X} = \mathbf{M_X}\mathbf{P_X} = \mathbf{0}$.

\textit{Proof:} $\mathbf{P_X}\mathbf{M_X} = \mathbf{P_X}(\mathbf{I} - \mathbf{P_X}) = \mathbf{P_X} - \mathbf{P_X}^2 = \mathbf{P_X} - \mathbf{P_X} = \mathbf{0}$. $\square$

\textit{Interpretation:} The subspace $\mathcal{S}(\mathbf{X})$ and its orthogonal complement are perpendicular; projecting onto one annihilates any component in the other.

\textbf{Property 3.4 (Trace):} $\text{tr}(\mathbf{P_X}) = k$ and $\text{tr}(\mathbf{M_X}) = n - k$.

\textit{Proof:} $\text{tr}(\mathbf{P_X}) = \text{tr}(\mathbf{X}(\mathbf{X}^T\mathbf{X})^{-1}\mathbf{X}^T) = \text{tr}((\mathbf{X}^T\mathbf{X})^{-1}\mathbf{X}^T\mathbf{X}) = \text{tr}(\mathbf{I}_k) = k$, using the cyclic property of trace. $\square$

\textit{Interpretation:} The trace of a projection matrix equals the dimension of the subspace it projects onto. The hat matrix projects onto a $k$-dimensional space; the annihilator projects onto an $(n-k)$-dimensional space.

\textbf{Property 3.5 (Invariance of projection):} $\mathbf{P_X}\mathbf{X} = \mathbf{X}$.

\textit{Proof:} Any vector in $\mathcal{S}(\mathbf{X})$ maps to itself under the projection. $\mathbf{X}$ is in $\mathcal{S}(\mathbf{X})$ (each column is). $\square$

\textbf{Property 3.6 (OLS residuals and regressors are orthogonal):} $\mathbf{X}^T\hat{\mathbf{u}} = \mathbf{0}$.

\textit{Proof:} $\mathbf{X}^T\hat{\mathbf{u}} = \mathbf{X}^T\mathbf{M_X}\mathbf{y} = (\mathbf{M_X}\mathbf{X})^T\mathbf{y} = \mathbf{0}^T\mathbf{y} = \mathbf{0}$. $\square$

\begin{figure}[htbp]
  \centering
  \begin{tikzpicture}[scale=1.3,
    vec/.style={-{Stealth[scale=1.2]}, very thick},
    dashvec/.style={-{Stealth[scale=1.1]}, thick, dashed}
  ]
    % Column space plane
    \fill[econblue!8, draw=econblue!40, thick]
      (-0.5,-0.5) -- (4.5,-0.5) -- (3.5,1.5) -- (-1.5,1.5) -- cycle;
    \node[font=\small\itshape, econblue!60!black] at (3.5, 0.2) {col$(\mathbf{X})$};

    % Origin
    \filldraw[black] (1,0.5) circle (1.5pt) node[below left, font=\small] {$\mathbf{0}$};

    % y vector
    \draw[vec, red!70!black] (1,0.5) -- (2.5,3.5) node[right, font=\small\bfseries] {$\mathbf{y}$};

    % y_hat (projection)
    \draw[vec, econblue, very thick] (1,0.5) -- (2.5,0.9) node[below right, font=\small\bfseries] {$\hat{\mathbf{y}} = \mathbf{P_X y}$};

    % residual vector
    \draw[dashvec, gray!70!black] (2.5,0.9) -- (2.5,3.5) node[right, font=\small\bfseries] {$\hat{\mathbf{u}} = \mathbf{M_X y}$};

    % Right angle mark
    \draw[gray, thick] (2.5,0.9) ++(0,0.2) -- ++(0.2,0) -- ++(0,-0.2);

    \node[font=\footnotesize, gray] at (1, -1.0) {$\mathbf{y} = \underbrace{\hat{\mathbf{y}}}_{\in\,\text{col}(\mathbf{X})} + \underbrace{\hat{\mathbf{u}}}_{\perp\,\text{col}(\mathbf{X})}$};
  \end{tikzpicture}
  \caption{The geometry of OLS. The OLS estimator projects $\mathbf{y}$ orthogonally onto the column space of $\mathbf{X}$. The fitted values $\hat{\mathbf{y}} = \mathbf{P_X y}$ lie in col$(\mathbf{X})$, and the residual vector $\hat{\mathbf{u}} = \mathbf{M_X y}$ is perpendicular to col$(\mathbf{X})$ (i.e., $\mathbf{X}^T\hat{\mathbf{u}} = \mathbf{0}$).}
  \label{fig:ch03_projection}
\end{figure}

\medskip\hrule\medskip

\section{3.6 The Pythagorean Theorem and Goodness of Fit ($R^2$)}

\subsection{3.6.1 Sum of Squares Decomposition}

By the geometry of projections (Pythagorean theorem in $E^n$):

\[
\|\mathbf{y}\|^2 = \|\mathbf{P_X}\mathbf{y}\|^2 + \|\mathbf{M_X}\mathbf{y}\|^2
\]

In more familiar notation (when an intercept is included in $\mathbf{X}$):

\[
\underbrace{\|\mathbf{y} - \bar{y}\boldsymbol{\iota}\|^2}_{TSS} = \underbrace{\|\hat{\mathbf{y}} - \bar{y}\boldsymbol{\iota}\|^2}_{ESS} + \underbrace{\|\hat{\mathbf{u}}\|^2}_{SSR}
\]

where:
\begin{itemize}
\item $TSS = \sum_{i=1}^n (y_i - \bar{y})^2$ is the \textbf{Total Sum of Squares}
\item $ESS = \sum_{i=1}^n (\hat{y}_i - \bar{y})^2$ is the \textbf{Explained Sum of Squares}
\item $SSR = \sum_{i=1}^n \hat{u}_i^2$ is the \textbf{Sum of Squared Residuals}
\end{itemize}

\textit{Proof that $TSS = ESS + SSR$:} Write $\mathbf{y} - \bar{y}\boldsymbol{\iota} = (\hat{\mathbf{y}} - \bar{y}\boldsymbol{\iota}) + \hat{\mathbf{u}}$. Squaring: $\|\mathbf{y} - \bar{y}\boldsymbol{\iota}\|^2 = \|\hat{\mathbf{y}} - \bar{y}\boldsymbol{\iota}\|^2 + \|\hat{\mathbf{u}}\|^2 + 2(\hat{\mathbf{y}} - \bar{y}\boldsymbol{\iota})^T\hat{\mathbf{u}}$. The cross term is zero because $\hat{\mathbf{u}} \perp \mathcal{S}(\mathbf{X})$ and $(\hat{\mathbf{y}} - \bar{y}\boldsymbol{\iota}) \in \mathcal{S}(\mathbf{X})$. $\square$

\subsection{3.6.2 The Coefficient of Determination ($R^2$)}

\textbf{Definition 3.1 (Centered $R^2$):}

\[
R_c^2 = \frac{ESS}{TSS} = 1 - \frac{SSR}{TSS} = \frac{\|\mathbf{P_X}\mathbf{M_\iota}\mathbf{y}\|^2}{\|\mathbf{M_\iota}\mathbf{y}\|^2}
\]

where $\mathbf{M_\iota} = \mathbf{I} - \frac{1}{n}\boldsymbol{\iota}\boldsymbol{\iota}^T$ is the demeaning matrix.

\textbf{Definition 3.2 (Uncentered $R^2$):}

\[
R_u^2 = \frac{\|\mathbf{P_X}\mathbf{y}\|^2}{\|\mathbf{y}\|^2}
\]

\textbf{Properties of $R^2$:}

\begin{enumerate}
\item $0 \leq R^2 \leq 1$ (when an intercept is included).
\item $R^2$ is the square of the sample correlation between $\hat{y}_i$ and $y_i$.
\item Adding more regressors never decreases $R^2$ (but may be pure noise fitting).
\item $R^2$ measures \textbf{in-sample fit}, not predictive accuracy or causal validity.
\end{enumerate}

\begin{keybox}{Warning:** A high $R^2$ does not mean the regression is correctly specified or that the coefficients are causally interpretable. A low $R^2$ does not mean the regression is useless --- in many empirical settings, including program evaluation, $R^2$ values of 0.05--0.20 are common even when the estimates are valid.}
\end{keybox}

\textbf{Definition 3.3 (Adjusted $R^2$):} To penalize for the number of regressors:

\[
\bar{R}^2 = 1 - \frac{SSR/(n-k)}{TSS/(n-1)} = 1 - (1 - R^2)\frac{n-1}{n-k}
\]

The adjusted $R^2$ can decrease when an irrelevant variable is added.

\begin{figure}[htbp]
  \centering
  \begin{tikzpicture}
    \begin{axis}[
      width=9cm, height=6cm,
      xlabel={$X$}, ylabel={$Y$},
      xmin=0, xmax=6, ymin=0, ymax=5,
      axis lines=left,
      tick style={draw=none},
      xtick=\empty, ytick={2.5},
      yticklabels={$\bar{y}$},
    ]
      % Data points
      \addplot[only marks, mark=*, mark size=2pt, color=gray!50]
        coordinates {(1,1)(2,2.2)(3,2.8)(4,3.5)(5,4.2)};

      % OLS line
      \addplot[thick, econblue, domain=0.5:5.5] {0.65*x + 0.6};

      % Mean line
      \addplot[thick, dashed, gray, domain=0.3:5.7] {2.5};

      % Decomposition for one point (x=4, y=3.5, yhat~3.2, ybar=2.5)
      \draw[<->, thick, red!70!black] (axis cs:4.2, 2.5) -- (axis cs:4.2, 3.5)
        node[midway, right, font=\footnotesize, red!70!black] {$y_i - \bar{y}$ (TSS)};
      \draw[<->, thick, econblue] (axis cs:3.8, 2.5) -- (axis cs:3.8, 3.2)
        node[midway, left, font=\footnotesize, econblue] {$\hat{y}_i-\bar{y}$ (ESS)};
      \draw[<->, thick, orange!80!black] (axis cs:4, 3.2) -- (axis cs:4, 3.5)
        node[midway, right, font=\footnotesize, orange!80!black] {$\hat{u}_i$ (RSS)};

      \node[font=\footnotesize, gray] at (axis cs:1.5, 4.5) {$R^2 = \text{ESS}/\text{TSS}$};
    \end{axis}
  \end{tikzpicture}
  \caption{Decomposition of variation in the OLS model. For each observation, the total deviation from the mean $y_i - \bar{y}$ (TSS component) is split into the explained part $\hat{y}_i - \bar{y}$ (ESS) and the residual $\hat{u}_i$ (RSS). $R^2 = \text{ESS}/\text{TSS}$ measures the share of variance explained.}
  \label{fig:ch03_r_squared}
\end{figure}

\medskip\hrule\medskip

\section{3.7 Linear Transformations and Invariance Results}

\subsection{3.7.1 Rescaling Regressors}

If we replace $\mathbf{X}$ with $\mathbf{Z} = \mathbf{X}\mathbf{A}$ where $\mathbf{A}$ is a $k \times k$ full-rank matrix, then the OLS fitted values are unchanged:

\[
\mathbf{P_Z} = \mathbf{P_X}
\]

and the coefficient estimates satisfy:

\[
\mathbf{A}\hat{\boldsymbol{\alpha}} = \hat{\boldsymbol{\beta}}
\]

where $\hat{\boldsymbol{\alpha}}$ is the OLS estimate from regressing $\mathbf{y}$ on $\mathbf{Z}$.

\textbf{Implication:} OLS coefficients change predictably when units change. If we measure income in \$000 instead of dollars, the coefficient on income is multiplied by 1000. If we take the log of income, the coefficient changes entirely. But the fitted values, residuals, and $R^2$ are invariant to linear transformations.

\subsection{3.7.2 Demeaning as a Linear Transformation}

An important special case: subtracting the mean $\bar{x}$ from each regressor (demeaning).

\[
Z_i = X_i - \bar{X}
\]

Under this transformation:
\begin{itemize}
\item The intercept changes but all slope coefficients are unchanged.
\item The $R^2$ and residuals are unchanged.
\item The demeaned regression $\tilde{y}_i = \mathbf{M_\iota}\mathbf{y}$ on $\tilde{X}_i = \mathbf{M_\iota}\mathbf{X}$ gives the same slopes as the original regression.
\end{itemize}

This is a special case of the \textbf{Frisch-Waugh-Lovell (FWL) Theorem}.

\subsection{3.7.3 The Frisch-Waugh-Lovell Theorem}

This theorem is one of the most useful results in regression analysis. It shows that the coefficients on a subset of regressors can be obtained by "partialling out" the other regressors.

\textbf{Theorem 3.1 (Frisch-Waugh-Lovell):} Partition the design matrix as $\mathbf{X} = [\mathbf{X_1} \: \mathbf{X_2}]$, where $\mathbf{X_1}$ is $n \times k_1$ and $\mathbf{X_2}$ is $n \times k_2$, with $k_1 + k_2 = k$. The OLS estimator of $\boldsymbol{\beta_2}$ from the full regression $\mathbf{y} = \mathbf{X_1}\boldsymbol{\beta_1} + \mathbf{X_2}\boldsymbol{\beta_2} + \mathbf{u}$ equals the OLS estimator from the "residualized" regression:

\[
\hat{\boldsymbol{\beta}}_2^{FWL} = (\tilde{\mathbf{X}}_2^T \tilde{\mathbf{X}}_2)^{-1} \tilde{\mathbf{X}}_2^T \tilde{\mathbf{y}}
\]

where $\tilde{\mathbf{y}} = \mathbf{M_{X_1}}\mathbf{y}$ and $\tilde{\mathbf{X}}_2 = \mathbf{M_{X_1}}\mathbf{X_2}$ are the residuals from regressing $\mathbf{y}$ and $\mathbf{X_2}$, respectively, on $\mathbf{X_1}$.

Moreover, $\hat{\boldsymbol{\beta}}_2^{FWL} = \hat{\boldsymbol{\beta}}_2^{OLS}$ and the residuals from both regressions are identical.

\textbf{Proof sketch:} Let $\tilde{\mathbf{y}} = \mathbf{M_{X_1}}\mathbf{y}$ and $\tilde{\mathbf{X}}_2 = \mathbf{M_{X_1}}\mathbf{X}_2$. Then $\tilde{\mathbf{X}}_2^T\tilde{\mathbf{y}} = \mathbf{X}_2^T\mathbf{M_{X_1}}\mathbf{y}$ (since $\mathbf{M_{X_1}}$ is symmetric). One can show by substitution that the normal equations of the FWL regression are equivalent to the normal equations of the full regression projected onto the $\boldsymbol{\beta}_2$ block. The residuals are the same because $\mathbf{M_X} = \mathbf{M_{X_1}}\mathbf{M}_{[\mathbf{M_{X_1}}\mathbf{X_2}]}$. (See Davidson and MacKinnon (2004), Chapter 2 for a complete proof.) $\square$

\begin{keybox}{Why the FWL Theorem Matters}

1. \textbf{Interpreting multiple regression.} The coefficient on $X_{2i}$ in a multiple regression is the effect of $X_{2i}$ on $y_i$ after "controlling for" $X_{1i}$ --- specifically, after removing the part of both $y_i$ and $X_{2i}$ that is linearly predictable from $X_{1i}$.

2. \textbf{Partial effects.} The FWL theorem shows that OLS coefficients measure \textit{partial} effects: the marginal association between $y$ and one regressor, holding the other regressors fixed.

3. \textbf{Omitted variable bias.} If we omit $\mathbf{X_1}$ from the regression, $\hat{\boldsymbol{\beta}}_2$ will be biased if $\mathbf{X_1}$ is correlated with $\mathbf{X_2}$ (since we won't residualize $\mathbf{X_2}$ on $\mathbf{X_1}$).
\end{keybox}

\textbf{Example 3.3 (FWL in the Wage Regression):} In the wage regression with education and experience, the FWL theorem says:
\begin{enumerate}
\item Regress log wages on experience; save residuals $\tilde{y}_i$.
\item Regress education on experience; save residuals $\tilde{\text{educ}}_i$.
\item Regress $\tilde{y}_i$ on $\tilde{\text{educ}}_i$ (no intercept needed); the slope is the education coefficient from the full regression.
\end{enumerate}

The slope in step 3 measures the return to education \textit{after} removing the variation in both wages and education that is due to experience.

\medskip\hrule\medskip

\section{Chapter Summary}

\begin{itemize}
\item The multiple regression model in matrix notation is $\mathbf{y} = \mathbf{X}\boldsymbol{\beta} + \mathbf{u}$.
\item The OLS estimator $\hat{\boldsymbol{\beta}} = (\mathbf{X}^T\mathbf{X})^{-1}\mathbf{X}^T\mathbf{y}$ is derived by minimizing the SSR or equating sample moments to zero.
\item OLS is geometrically an \textbf{orthogonal projection} of $\mathbf{y}$ onto the column space of $\mathbf{X}$.
\item The \textbf{projection matrices} $\mathbf{P_X} = \mathbf{X}(\mathbf{X}^T\mathbf{X})^{-1}\mathbf{X}^T$ and $\mathbf{M_X} = \mathbf{I} - \mathbf{P_X}$ are symmetric, idempotent, and annihilate each other.
\item The \textbf{Pythagorean decomposition} $TSS = ESS + SSR$ follows from the geometry of projections.
\item The \textbf{$R^2$} measures in-sample fit as the fraction of variance explained by the regressors.
\item The \textbf{FWL Theorem} shows that OLS coefficients measure partial effects after residualizing on the other regressors.
\end{itemize}

\medskip\hrule\medskip

\section{Key Terms}

\textbf{Design matrix} ($\mathbf{X}$) --- the $n \times k$ matrix of observations on the regressors.

\textbf{Normal equations} --- $\mathbf{X}^T\mathbf{X}\hat{\boldsymbol{\beta}} = \mathbf{X}^T\mathbf{y}$; the first-order conditions for the OLS estimator.

\textbf{Full rank} --- $\mathbf{X}$ has full column rank $k$ (no perfect multicollinearity).

\textbf{OLS estimator} --- $\hat{\boldsymbol{\beta}} = (\mathbf{X}^T\mathbf{X})^{-1}\mathbf{X}^T\mathbf{y}$.

\textbf{Column space} ($\mathcal{S}(\mathbf{X})$) --- the set of all vectors of the form $\mathbf{X}\mathbf{b}$ for $\mathbf{b} \in \mathbb{R}^k$.

\textbf{Projection matrix} ($\mathbf{P_X}$) --- $\mathbf{X}(\mathbf{X}^T\mathbf{X})^{-1}\mathbf{X}^T$; projects onto $\mathcal{S}(\mathbf{X})$.

\textbf{Annihilator matrix} ($\mathbf{M_X}$) --- $\mathbf{I} - \mathbf{P_X}$; projects onto the orthogonal complement of $\mathcal{S}(\mathbf{X})$.

\textbf{Idempotent} --- a matrix $\mathbf{A}$ such that $\mathbf{A}^2 = \mathbf{A}$.

\textbf{Frisch-Waugh-Lovell (FWL) Theorem} --- the coefficient on $\mathbf{X_2}$ in the full regression equals the coefficient from regressing residualized $\mathbf{y}$ on residualized $\mathbf{X_2}$.

\textbf{$R^2$} --- fraction of total variation in $y$ explained by $\mathbf{X}$; $R^2 = ESS/TSS$.

\medskip\hrule\medskip

\section{Exercises}

\subsection{Conceptual Questions}

\textbf{3.1} Explain in your own words why the OLS estimator can be interpreted as an orthogonal projection.

\textbf{3.2} What is the geometric interpretation of the OLS residuals $\hat{\mathbf{u}}$? What does it mean for $\hat{\mathbf{u}}$ to be orthogonal to the columns of $\mathbf{X}$?

\textbf{3.3} "If we add a variable that is completely uncorrelated with all existing regressors and the outcome variable, $R^2$ does not change." Is this true? Explain.

\textbf{3.4} You run a regression of test scores on class size and find $R^2 = 0.02$. A colleague says the regression is worthless because $R^2$ is so low. How would you respond?

\textbf{3.5} State the FWL Theorem. Give an intuitive explanation of why the "residualized regression" gives the same coefficient as the full regression.

\subsection{Analytical Questions}

\textbf{3.6} Prove that the hat matrix $\mathbf{P_X}$ is symmetric and idempotent.

\textbf{3.7} Show that $\hat{\boldsymbol{\beta}}$ from a regression of $\mathbf{y}$ on $\mathbf{X}$ and $\hat{\boldsymbol{\beta}}$ from a regression of $\mathbf{M_\iota}\mathbf{y}$ on $\mathbf{M_\iota}\mathbf{X}$ (i.e., demeaned regression) give the same slope coefficients.

\textbf{3.8} Consider the simple bivariate regression $y_i = \alpha + \beta x_i + u_i$. 

(a) Show that $\hat{\beta} = \frac{\sum_{i=1}^n (x_i - \bar{x})(y_i - \bar{y})}{\sum_{i=1}^n (x_i - \bar{x})^2}$.

(b) Show that $\hat{\alpha} = \bar{y} - \hat{\beta}\bar{x}$.

(c) Show that the OLS regression line passes through the point $(\bar{x}, \bar{y})$.

\textbf{3.9} Consider the regression of $y_i$ on $x_i$ and $z_i$. If $x_i$ and $z_i$ are orthogonal in the sample ($\sum x_i z_i = 0$ after demeaning), show that the OLS coefficient on $x_i$ from the multiple regression equals the OLS coefficient on $x_i$ from the simple regression of $y_i$ on $x_i$ alone.

\textbf{3.10} Show that $\text{tr}(\mathbf{M_X}) = n - k$.

\textbf{3.11} Let $\mathbf{X}_1 = \boldsymbol{\iota}$ (column of ones). Show that $\mathbf{M_{X_1}} = \mathbf{M_\iota}$ is the demeaning matrix: $\mathbf{M_\iota}\mathbf{y} = \mathbf{y} - \bar{y}\boldsymbol{\iota}$.

\subsection{Applied Questions}

\textbf{3.12} You have data on $n = 500$ workers with: log wages ($y$), years of education ($x_1$), and a dummy for urban residence ($x_2$). 

(a) Write out the full OLS estimator in matrix form.
(b) Describe what the FWL theorem says about the coefficient on education in this regression.
(c) If education is positively correlated with urban residence, and urban workers earn higher wages, what happens to the education coefficient if you omit the urban dummy?

\textbf{3.13} Using the data from R Lab 2, verify the FWL Theorem numerically: (a) run the full regression of log wages on education and experience; (b) residualize log wages and education on experience; (c) regress the residualized wages on residualized education; (d) confirm the slopes in (a) and (c) match.

\medskip\hrule\medskip

\section{R Lab 3: OLS in Matrix Form and the Geometry of Regression}

\subsection{Objectives}
\begin{itemize}
\item Compute OLS directly using matrix algebra
\item Construct and verify projection matrices
\item Verify the FWL theorem numerically
\item Visualize the sum-of-squares decomposition
\end{itemize}

\subsection{Setup}

\begin{lstlisting}
library(tidyverse)
library(AER)
data("CPS1988")
CPS1988 <- CPS1988 %>% mutate(log_wage = log(wage))
\end{lstlisting}

\subsection{Exercise 3.1: OLS via Matrix Algebra}

\begin{lstlisting}
# Create the design matrix
y <- as.matrix(CPS1988$log_wage)
X <- cbind(1, CPS1988$education, CPS1988$experience, CPS1988$experience^2)
colnames(X) <- c("intercept", "education", "experience", "experience_sq")

n <- nrow(X)
k <- ncol(X)

# OLS estimator: beta_hat = (X'X)^{-1} X'y
XtX     <- t(X) %*% X
Xty     <- t(X) %*% y
beta_hat <- solve(XtX) %*% Xty
print(round(beta_hat, 4))

# Compare with lm()
lm_model <- lm(log_wage ~ education + experience + I(experience^2), data = CPS1988)
print(round(coef(lm_model), 4))
\end{lstlisting}

\subsection{Exercise 3.2: Projection Matrices}

\begin{lstlisting}
# Hat matrix P_X
P_X <- X %*% solve(XtX) %*% t(X)

# Annihilator M_X
M_X <- diag(n) - P_X

# Verify properties
cat("Is P_X idempotent?", all.equal(P_X %*% P_X, P_X, check.attributes = FALSE), "\n")
cat("Is M_X idempotent?", all.equal(M_X %*% M_X, M_X, check.attributes = FALSE), "\n")

# Fitted values and residuals
y_hat <- P_X %*% y        # fitted values
u_hat <- M_X %*% y        # residuals

# Check: X'u_hat = 0 (should be numerically zero)
cat("Max |X'u_hat|:", max(abs(t(X) %*% u_hat)), "\n")

# Trace of projection matrices
cat("tr(P_X) =", round(sum(diag(P_X)), 2), "(should be", k, ")\n")
cat("tr(M_X) =", round(sum(diag(M_X)), 2), "(should be", n - k, ")\n")
\end{lstlisting}

\subsection{Exercise 3.3: Sum of Squares Decomposition}

\begin{lstlisting}
# Demeaning matrix
iota <- matrix(1, n, 1)
M_iota <- diag(n) - (1/n) * (iota %*% t(iota))

TSS <- as.numeric(t(y)   %*% M_iota %*% y)   # Total SS (demeaned)
ESS <- as.numeric(t(y_hat) %*% M_iota %*% y_hat)  # Explained SS
SSR <- as.numeric(t(u_hat) %*% u_hat)          # Residual SS

cat("TSS:", round(TSS, 4), "\n")
cat("ESS:", round(ESS, 4), "\n")
cat("SSR:", round(SSR, 4), "\n")
cat("ESS + SSR:", round(ESS + SSR, 4), "(should equal TSS)\n")

R_squared <- ESS / TSS
cat("R-squared:", round(R_squared, 4), "\n")
cat("R-squared from lm():", round(summary(lm_model)$r.squared, 4), "\n")
\end{lstlisting}

\subsection{Exercise 3.4: Verifying the FWL Theorem}

\begin{lstlisting}
# Full regression: log_wage ~ education + experience + experience^2
# FWL for the education coefficient

# Step 1: Partition X1 (intercept + experience) and X2 (education)
X1 <- cbind(1, CPS1988$experience, CPS1988$experience^2)
X2 <- as.matrix(CPS1988$education)

# Step 2: Residualize y and X2 on X1
M_X1 <- diag(n) - X1 %*% solve(t(X1) %*% X1) %*% t(X1)
y_tilde  <- M_X1 %*% y
X2_tilde <- M_X1 %*% X2

# Step 3: FWL regression
beta_FWL <- solve(t(X2_tilde) %*% X2_tilde) %*% t(X2_tilde) %*% y_tilde
cat("FWL estimate of education coefficient:", round(beta_FWL, 6), "\n")
cat("Full OLS estimate of education coefficient:", round(beta_hat["education",], 6), "\n")

# Confirm residuals are the same
u_FWL <- y_tilde - X2_tilde %*% beta_FWL
cat("Max difference in residuals:", max(abs(u_hat - u_FWL)), "\n")
\end{lstlisting}

\subsection{Exercise 3.5: Visualizing OLS Geometry (2D Approximation)}

\begin{lstlisting}
# Simplified example: 3 observations, 2 regressors (for visualization)
set.seed(123)
n_small <- 3
x_small <- c(1, 2, 4)
y_small <- c(2, 3, 5)
X_small <- cbind(1, x_small)
beta_small <- solve(t(X_small) %*% X_small) %*% t(X_small) %*% y_small
y_hat_small <- X_small %*% beta_small
u_small <- y_small - y_hat_small

cat("Observations:", y_small, "\n")
cat("Fitted values:", round(y_hat_small, 3), "\n")
cat("Residuals:", round(u_small, 3), "\n")
cat("OLS coefficients:", round(beta_small, 3), "\n")
cat("SSR:", round(sum(u_small^2), 4), "\n")

# Verify orthogonality
cat("X'u_hat:", round(t(X_small) %*% u_small, 6), "\n")
\end{lstlisting}

\medskip\hrule\medskip

\textit{End of Chapter 3}

\medskip\hrule\medskip

\textbf{References}

Davidson, R., and MacKinnon, J. G. (2004). \textit{Econometric Theory and Methods}. Oxford University Press. Chapter 2, Sections 2.1--2.4.

Angrist, J. D., and Pischke, J. (2009). \textit{Mostly Harmless Econometrics}. Princeton University Press. Chapter 3.

Wooldridge, J. M. (2019). \textit{Introductory Econometrics} (7th ed.). Cengage. Chapters 3--4.

Frisch, R., and Waugh, F. V. (1933). Partial time regressions as compared with individual trends. \textit{Econometrica}, 1(4), 387--401.

Lovell, M. C. (1963). Seasonal adjustment of economic time series and multiple regression analysis. \textit{Journal of the American Statistical Association}, 58(304), 993--1010.

\chapter{Chapter 4: Statistical Properties of OLS --- Finite Sample Theory}

\medskip\hrule\medskip

\section{Chapter Overview}

Chapter 3 established the \textit{algebraic} properties of OLS --- results that hold in any sample, without any distributional assumptions. This chapter turns to the \textit{statistical} properties of OLS: is $\hat{\boldsymbol{\beta}}$ close to the true parameter $\boldsymbol{\beta}$? How precise is it? How do we conduct valid inference (hypothesis tests and confidence intervals)?

We answer these questions first in the \textbf{finite-sample} framework, which requires stronger assumptions but gives exact distributional results. We state the classical \textbf{Gauss-Markov assumptions}, prove \textbf{unbiasedness} and the \textbf{Gauss-Markov Theorem} (OLS is BLUE), and derive inference procedures under normality. Chapter 5 relaxes the normality and homoskedasticity assumptions using large-sample (asymptotic) theory.

\medskip\hrule\medskip

\section{Learning Objectives}

After completing this chapter, students will be able to:

\begin{enumerate}
\item State the \textbf{classical (Gauss-Markov) assumptions} for the linear regression model.
\item Prove that OLS is \textbf{unbiased} under strict exogeneity.
\item Derive the \textbf{variance-covariance matrix} of the OLS estimator.
\item State and prove the \textbf{Gauss-Markov Theorem}.
\item Derive the \textbf{$t$-statistic} and construct \textbf{confidence intervals} for individual coefficients.
\item Construct and interpret the \textbf{$F$-statistic} for joint hypothesis tests.
\item Explain the \textbf{estimated variance} of OLS and why $\hat{\sigma}^2 = SSR/(n-k)$.
\item Discuss the sources of estimator variance (sample size, multicollinearity, error variance).
\end{enumerate}

\medskip\hrule\medskip

\section{4.1 Setting Up the Framework: Classical Assumptions}

To derive statistical properties of OLS, we need assumptions about the data-generating process (DGP). The \textbf{Classical Linear Regression Model (CLRM)} rests on five assumptions.

\textbf{Assumption 4.1 (Linearity):} The population model is correctly specified as:
\[
\mathbf{y} = \mathbf{X}\boldsymbol{\beta} + \mathbf{u}
\]
where $\boldsymbol{\beta}$ is the true (unknown) parameter vector.

\textbf{Assumption 4.2 (Full Rank):} The design matrix $\mathbf{X}$ has full column rank $k$ (no perfect multicollinearity).

\textbf{Assumption 4.3 (Strict Exogeneity):}
\[
E[\mathbf{u} \mid \mathbf{X}] = \mathbf{0}
\]
This says that the conditional expectation of the error vector, given the entire design matrix, is zero. This implies $E[u_i \mid \mathbf{x}_1, \mathbf{x}_2, \ldots, \mathbf{x}_n] = 0$: the error for observation $i$ is mean-zero conditional on \textit{all} observations' covariates.

\textbf{Assumption 4.4 (Spherical Errors / Homoskedasticity and No Serial Correlation):}
\[
E[\mathbf{u}\mathbf{u}^T \mid \mathbf{X}] = \sigma^2 \mathbf{I}_n
\]
This combines two conditions:
\begin{itemize}
\item \textit{Homoskedasticity:} $E[u_i^2 \mid \mathbf{X}] = \sigma^2$ for all $i$ (constant variance).
\item \textit{No serial correlation:} $E[u_i u_j \mid \mathbf{X}] = 0$ for all $i \neq j$ (uncorrelated errors).
\end{itemize}

\textbf{Assumption 4.5 (Normality):} The errors are jointly normally distributed:
\[
\mathbf{u} \mid \mathbf{X} \sim N(\mathbf{0}, \sigma^2 \mathbf{I}_n)
\]

\begin{keybox}{Note on Assumptions:\textbf{ Assumptions 4.1--4.4 are the }Gauss-Markov assumptions**. Assumption 4.5 (normality) is additional and is needed only for \textit{exact} finite-sample inference. As we show in Chapter 5, Assumption 4.5 can be dropped in large samples.}
\end{keybox}

\subsection{4.1.1 Discussion of Strict Exogeneity}

Assumption 4.3 is the critical identifying assumption. It requires that the error $u_i$ is uncorrelated with \textit{all} regressors in \textit{all} observations. This is stronger than merely requiring $E[u_i \mid \mathbf{x}_i] = 0$ (conditional mean zero for observation $i$'s own regressors).

\textbf{When strict exogeneity fails:}
\begin{itemize}
\item \textbf{Omitted variable bias:} An important determinant of $y$ that is correlated with $\mathbf{x}_i$ is left out of the model.
\item \textbf{Reverse causality:} $y$ causes $\mathbf{x}$ as well as vice versa.
\item \textbf{Measurement error:} $\mathbf{x}_i$ is measured with error that is correlated with the true regressor.
\item \textbf{Panel data with lagged dependent variables:} $E[u_{it} \mid x_{i,t-1}]$ may be non-zero even if $E[u_{it} \mid x_{it}] = 0$.
\end{itemize}

When Assumption 4.3 fails, OLS is \textbf{biased} and potentially \textbf{inconsistent}. Addressing this is the focus of Chapters 6--13.

\medskip\hrule\medskip

\section{4.2 Unbiasedness of OLS}

\textbf{Theorem 4.1 (Unbiasedness of OLS):} Under Assumptions 4.1--4.3, the OLS estimator is \textbf{unbiased}: $E[\hat{\boldsymbol{\beta}} \mid \mathbf{X}] = \boldsymbol{\beta}$.

\textbf{Proof:}

\[
\hat{\boldsymbol{\beta}} = (\mathbf{X}^T\mathbf{X})^{-1}\mathbf{X}^T\mathbf{y} = (\mathbf{X}^T\mathbf{X})^{-1}\mathbf{X}^T(\mathbf{X}\boldsymbol{\beta} + \mathbf{u}) = \boldsymbol{\beta} + (\mathbf{X}^T\mathbf{X})^{-1}\mathbf{X}^T\mathbf{u}
\]

Taking the conditional expectation given $\mathbf{X}$:

\[
E[\hat{\boldsymbol{\beta}} \mid \mathbf{X}] = \boldsymbol{\beta} + (\mathbf{X}^T\mathbf{X})^{-1}\mathbf{X}^T \underbrace{E[\mathbf{u} \mid \mathbf{X}]}_{= \mathbf{0}} = \boldsymbol{\beta}
\]

By the \textbf{Law of Iterated Expectations}: $E[\hat{\boldsymbol{\beta}}] = E[E[\hat{\boldsymbol{\beta}} \mid \mathbf{X}]] = E[\boldsymbol{\beta}] = \boldsymbol{\beta}$. $\square$

\textbf{Interpretation:} On average (over repeated samples), the OLS estimator hits the true parameter. In any given sample, $\hat{\boldsymbol{\beta}} \neq \boldsymbol{\beta}$ due to sampling variation, but there is no systematic error.

\subsection{4.2.1 Importance of Strict Exogeneity}

The proof relies critically on Assumption 4.3. If $E[\mathbf{u} \mid \mathbf{X}] \neq \mathbf{0}$, then:

\[
E[\hat{\boldsymbol{\beta}} \mid \mathbf{X}] = \boldsymbol{\beta} + (\mathbf{X}^T\mathbf{X})^{-1}\mathbf{X}^T E[\mathbf{u} \mid \mathbf{X}] \neq \boldsymbol{\beta}
\]

The OLS estimator is biased. The bias is:

\[
\text{Bias}(\hat{\boldsymbol{\beta}}) = E[\hat{\boldsymbol{\beta}}] - \boldsymbol{\beta} = (\mathbf{X}^T\mathbf{X})^{-1}\mathbf{X}^T E[\mathbf{u}]
\]

\medskip\hrule\medskip

\section{4.3 The Variance-Covariance Matrix of OLS}

\textbf{Theorem 4.2 (Variance of OLS under Spherical Errors):} Under Assumptions 4.1--4.4:

\[
\text{Var}(\hat{\boldsymbol{\beta}} \mid \mathbf{X}) = \sigma^2 (\mathbf{X}^T\mathbf{X})^{-1}
\]

\textbf{Proof:} From the proof of Theorem 4.1: $\hat{\boldsymbol{\beta}} - \boldsymbol{\beta} = (\mathbf{X}^T\mathbf{X})^{-1}\mathbf{X}^T\mathbf{u}$.

\[
\text{Var}(\hat{\boldsymbol{\beta}} \mid \mathbf{X}) = E\left[(\hat{\boldsymbol{\beta}} - \boldsymbol{\beta})(\hat{\boldsymbol{\beta}} - \boldsymbol{\beta})^T \mid \mathbf{X}\right]
\]

\[
= (\mathbf{X}^T\mathbf{X})^{-1}\mathbf{X}^T \underbrace{E[\mathbf{u}\mathbf{u}^T \mid \mathbf{X}]}_{= \sigma^2 \mathbf{I}_n} \mathbf{X}(\mathbf{X}^T\mathbf{X})^{-1}
\]

\[
= \sigma^2 (\mathbf{X}^T\mathbf{X})^{-1}\mathbf{X}^T\mathbf{X}(\mathbf{X}^T\mathbf{X})^{-1} = \sigma^2 (\mathbf{X}^T\mathbf{X})^{-1} \quad \square
\]

\subsection{4.3.1 The Variance of Individual Coefficients}

The variance of the $j$-th coefficient $\hat{\beta}_j$ is the $j$-th diagonal element of $\sigma^2(\mathbf{X}^T\mathbf{X})^{-1}$:

\[
\text{Var}(\hat{\beta}_j \mid \mathbf{X}) = \sigma^2 [(\mathbf{X}^T\mathbf{X})^{-1}]_{jj}
\]

Using the FWL theorem, one can show:

\[
\text{Var}(\hat{\beta}_j \mid \mathbf{X}) = \frac{\sigma^2}{\text{SST}_j (1 - R^2_j)}
\]

where $\text{SST}_j = \sum_{i=1}^n (x_{ij} - \bar{x}_j)^2$ is the total variation in $x_{ij}$, and $R^2_j$ is the $R^2$ from regressing $x_{ij}$ on all other regressors.

\textbf{Three sources of estimator variance:}

\begin{center}
\begin{tabular}{@{} l l l @{}}
\toprule
\textbf{Factor} & \textbf{Effect on $\text{Var}(\hat{\beta}_j)$} & \textbf{Intuition} \\
\midrule
Larger $\sigma^2$ (more noise) & Increases & More noise $\to$ harder to estimate signal \\
Larger $\text{SST}_j$ (more variation in $x_j$) & Decreases & More variation $\to$ more information \\
Larger $R^2_j$ (more multicollinearity) & Increases & Collinearity $\to$ $x_j$ explains little net information \\
\bottomrule
\end{tabular}
\end{center}

\subsection{4.3.2 Estimating $\sigma^2$}

The population error variance $\sigma^2 = E[u_i^2]$ is unknown. We estimate it using the OLS residuals:

\[
\hat{\sigma}^2 = \frac{SSR}{n - k} = \frac{\hat{\mathbf{u}}^T\hat{\mathbf{u}}}{n - k}
\]

This is an \textbf{unbiased estimator} of $\sigma^2$ under Assumptions 4.1--4.4.

\textbf{Why divide by $n - k$ instead of $n$?}

\[
E[SSR \mid \mathbf{X}] = E[\hat{\mathbf{u}}^T\hat{\mathbf{u}} \mid \mathbf{X}] = E[\mathbf{u}^T\mathbf{M_X}\mathbf{u} \mid \mathbf{X}]
\]

\[
= \text{tr}(\mathbf{M_X} E[\mathbf{u}\mathbf{u}^T \mid \mathbf{X}]) = \sigma^2 \text{tr}(\mathbf{M_X}) = \sigma^2 (n - k)
\]

Therefore $E[\hat{\sigma}^2] = E[SSR/(n-k)] = \sigma^2$. $\square$

The degrees of freedom correction $n - k$ accounts for the fact that OLS "uses up" $k$ degrees of freedom to estimate $k$ parameters. With $k$ parameters perfectly fitting $k$ observations, there are no degrees of freedom left for estimation.

\medskip\hrule\medskip

\section{4.4 The Gauss-Markov Theorem: OLS is BLUE}

\textbf{Theorem 4.3 (Gauss-Markov Theorem):} Under Assumptions 4.1--4.4, the OLS estimator $\hat{\boldsymbol{\beta}}$ is the \textbf{Best Linear Unbiased Estimator (BLUE)} of $\boldsymbol{\beta}$. That is, among all linear unbiased estimators $\tilde{\boldsymbol{\beta}} = \mathbf{C}\mathbf{y}$ (where $\mathbf{C}$ is any non-random matrix with $E[\tilde{\boldsymbol{\beta}}] = \boldsymbol{\beta}$), OLS has the smallest variance matrix in the sense that:

\[
\text{Var}(\tilde{\boldsymbol{\beta}} \mid \mathbf{X}) - \text{Var}(\hat{\boldsymbol{\beta}} \mid \mathbf{X}) = \text{Var}(\tilde{\boldsymbol{\beta}} \mid \mathbf{X}) - \sigma^2 (\mathbf{X}^T\mathbf{X})^{-1} \geq 0 \quad \text{(positive semi-definite)}
\]

\textbf{Proof:} Let $\tilde{\boldsymbol{\beta}} = \mathbf{C}\mathbf{y}$ be any linear unbiased estimator. For unbiasedness: $E[\tilde{\boldsymbol{\beta}} \mid \mathbf{X}] = \mathbf{C}\mathbf{X}\boldsymbol{\beta} = \boldsymbol{\beta}$ for all $\boldsymbol{\beta}$, which requires $\mathbf{C}\mathbf{X} = \mathbf{I}_k$.

Write $\mathbf{C} = (\mathbf{X}^T\mathbf{X})^{-1}\mathbf{X}^T + \mathbf{D}$ where $\mathbf{D} = \mathbf{C} - (\mathbf{X}^T\mathbf{X})^{-1}\mathbf{X}^T$. Unbiasedness requires $\mathbf{D}\mathbf{X} = \mathbf{0}$.

Then:

\[
\text{Var}(\tilde{\boldsymbol{\beta}} \mid \mathbf{X}) = \mathbf{C} \cdot \sigma^2 \mathbf{I}_n \cdot \mathbf{C}^T = \sigma^2 \mathbf{C}\mathbf{C}^T
\]

\[
= \sigma^2 \left[(\mathbf{X}^T\mathbf{X})^{-1}\mathbf{X}^T\mathbf{X}(\mathbf{X}^T\mathbf{X})^{-1} + \mathbf{D}\mathbf{X}(\mathbf{X}^T\mathbf{X})^{-1} + (\mathbf{X}^T\mathbf{X})^{-1}\mathbf{X}^T\mathbf{D}^T + \mathbf{D}\mathbf{D}^T\right]
\]

Since $\mathbf{D}\mathbf{X} = \mathbf{0}$, the cross terms vanish:

\[
\text{Var}(\tilde{\boldsymbol{\beta}} \mid \mathbf{X}) = \sigma^2 (\mathbf{X}^T\mathbf{X})^{-1} + \sigma^2 \mathbf{D}\mathbf{D}^T = \text{Var}(\hat{\boldsymbol{\beta}} \mid \mathbf{X}) + \sigma^2 \mathbf{D}\mathbf{D}^T
\]

Since $\sigma^2 \mathbf{D}\mathbf{D}^T \geq 0$ (positive semi-definite), we have $\text{Var}(\tilde{\boldsymbol{\beta}}) \geq \text{Var}(\hat{\boldsymbol{\beta}})$. $\square$

\begin{keybox}{Interpreting Gauss-Markov:** The theorem says OLS is the most efficient estimator among all linear unbiased estimators. No other linear combination of the data produces unbiased coefficient estimates with smaller standard errors. This is a strong optimality result within the class of linear estimators.}
\end{keybox}

\textbf{Important caveats:}
\begin{enumerate}
\item Gauss-Markov only compares OLS to \textit{linear} estimators. Non-linear estimators (e.g., maximum likelihood under non-normal errors) may be more efficient.
\item The theorem requires spherical errors (Assumption 4.4). Under heteroskedasticity, \textbf{Feasible Generalized Least Squares (FGLS)} beats OLS in efficiency.
\item The theorem says nothing about whether OLS is \textit{consistent} (consistent requires Assumption 4.3 to be weakened, as we show in Chapter 5).
\end{enumerate}

\begin{figure}[htbp]
  \centering
  \begin{tikzpicture}
    \begin{axis}[
      width=10cm, height=6cm,
      xlabel={$\hat{\beta}$},
      ylabel={Density},
      xmin=-3, xmax=5,
      ymin=0, ymax=0.65,
      axis lines=left,
      tick style={draw=none},
      xtick={1},
      xticklabels={$\beta_0$ (true value)},
      ytick=\empty,
      legend style={at={(0.97,0.95)}, anchor=north east, font=\small},
    ]
      % OLS sampling distribution (efficient, centered)
      \addplot[very thick, econblue, domain=-3:5, samples=100]
        {0.5*exp(-0.5*(x-1)^2/0.6^2)/sqrt(2*pi*0.6^2)};

      % Alternative (less efficient)
      \addplot[thick, red!70!black, dashed, domain=-3:5, samples=100]
        {0.5*exp(-0.5*(x-1)^2/1.3^2)/sqrt(2*pi*1.3^2)};

      % Biased estimator
      \addplot[thick, orange!80!black, dotted, domain=-3:5, samples=100]
        {0.5*exp(-0.5*(x-2.2)^2/0.6^2)/sqrt(2*pi*0.6^2)};

      % True value line
      \draw[gray, dashed, thick] (axis cs:1,0) -- (axis cs:1,0.65);

      \legend{OLS (BLUE), Unbiased but inefficient, Biased estimator}
    \end{axis}
  \end{tikzpicture}
  \caption{Sampling distributions of three estimators of $\beta_0$ (true value marked by dashed vertical line). The OLS estimator (solid blue) is both centered on the truth (unbiased) and has the smallest variance among linear unbiased estimators (BLUE). An alternative unbiased estimator (dashed red) has higher variance. A biased estimator (dotted orange) is centered away from $\beta_0$.}
  \label{fig:ch04_sampling_dist}
\end{figure}

\medskip\hrule\medskip

\section{4.5 Inference Under Normality}

Under all five Gauss-Markov assumptions including normality (Assumption 4.5):

\[
\hat{\boldsymbol{\beta}} \mid \mathbf{X} \sim N\left(\boldsymbol{\beta}, \sigma^2(\mathbf{X}^T\mathbf{X})^{-1}\right)
\]

This follows immediately from the fact that $\hat{\boldsymbol{\beta}} = \boldsymbol{\beta} + (\mathbf{X}^T\mathbf{X})^{-1}\mathbf{X}^T\mathbf{u}$ is a linear transformation of $\mathbf{u}$, which is normal.

For any individual coefficient $\hat{\beta}_j$:

\[
\hat{\beta}_j \mid \mathbf{X} \sim N\left(\beta_j, \sigma^2 [(\mathbf{X}^T\mathbf{X})^{-1}]_{jj}\right)
\]

\subsection{4.5.1 The $t$-Statistic}

Since $\sigma^2$ is unknown, we replace it with $\hat{\sigma}^2$. The resulting statistic follows a $t$-distribution:

\[
t_j = \frac{\hat{\beta}_j - \beta_j^0}{\widehat{SE}(\hat{\beta}_j)} \sim t_{n-k}
\]

where $\beta_j^0$ is the value under the null hypothesis, and $\widehat{SE}(\hat{\beta}_j) = \hat{\sigma}\sqrt{[(\mathbf{X}^T\mathbf{X})^{-1}]_{jj}}$ is the estimated standard error.

\textbf{Proof:} We have $(\hat{\beta}_j - \beta_j) / (\sigma \sqrt{[(\mathbf{X}^T\mathbf{X})^{-1}]_{jj}}) \sim N(0,1)$ and $(n-k)\hat{\sigma}^2/\sigma^2 \sim \chi^2_{n-k}$ (under normality), and these two are independent. The ratio of a standard normal to the square root of an independent chi-squared divided by its degrees of freedom has a $t$-distribution. $\square$

\textbf{Hypothesis testing:} To test $H_0: \beta_j = 0$ against $H_1: \beta_j \neq 0$:

\begin{enumerate}
\item Compute $t_j = \hat{\beta}_j / \widehat{SE}(\hat{\beta}_j)$.
\item Compare to $t_{n-k, \alpha/2}$ (critical value from $t$-distribution with $n-k$ degrees of freedom).
\item Reject $H_0$ if $|t_j| > t_{n-k, \alpha/2}$.
\item The $p$-value is $P(|T| > |t_j|)$ where $T \sim t_{n-k}$.
\end{enumerate}

For large $n$, $t_{n-k}$ is well approximated by $N(0,1)$, and the "rule of thumb" for significance at the 5\% level is $|t| > 1.96$.

\subsection{4.5.2 Confidence Intervals}

A $(1-\alpha)$ confidence interval for $\beta_j$ is:

\[
\hat{\beta}_j \pm t_{n-k, \alpha/2} \cdot \widehat{SE}(\hat{\beta}_j)
\]

\textbf{Interpretation:} In repeated sampling, $(1-\alpha)$\% of such intervals will contain the true $\beta_j$. \textbf{It is not true} that the probability the true $\beta_j$ falls in a specific realized interval is $(1-\alpha)$\% --- the true $\beta_j$ is fixed; the interval is random.

\subsection{4.5.3 Practical Significance vs. Statistical Significance}

A statistically significant coefficient is not necessarily economically significant. A study with $n = 100{,}000$ observations can detect very small effects with high precision. Conversely, a study with $n = 50$ may fail to reject the null even for economically large effects.

\textbf{Good practice:} Report both the coefficient estimate (magnitude of economic effect) and the standard error (precision), not just the $t$-statistic.

\medskip\hrule\medskip

\section{4.6 Joint Hypothesis Tests: The $F$-Statistic}

Often we want to test joint hypotheses about multiple coefficients. For example:
\begin{itemize}
\item $H_0: \beta_2 = \beta_3 = 0$ (neither experience nor experience$^{2}$ matters)
\item $H_0: \beta_2 = \beta_3$ (two effects are equal)
\end{itemize}

\textbf{General setup:} Consider the null hypothesis $H_0: \mathbf{R}\boldsymbol{\beta} = \mathbf{r}$, where $\mathbf{R}$ is a $q \times k$ matrix of linear restrictions (rank $q$) and $\mathbf{r}$ is a $q \times 1$ vector.

\textbf{Definition 4.1 (F-Statistic):}

\[
F = \frac{(\mathbf{R}\hat{\boldsymbol{\beta}} - \mathbf{r})^T \left[\mathbf{R}(\mathbf{X}^T\mathbf{X})^{-1}\mathbf{R}^T\right]^{-1} (\mathbf{R}\hat{\boldsymbol{\beta}} - \mathbf{r})}{q \hat{\sigma}^2}
\]

Under $H_0$ and Assumptions 4.1--4.5: $F \sim F_{q, n-k}$.

\textbf{Equivalent Formulation via Restricted and Unrestricted Regressions:}

\[
F = \frac{(SSR_R - SSR_U)/q}{SSR_U/(n-k)}
\]

where $SSR_R$ is from the \textbf{restricted regression} (imposing $H_0$) and $SSR_U$ is from the \textbf{unrestricted regression}.

\textbf{Intuition:} If the restrictions are true, imposing them should not increase the SSR by much. If the $F$-statistic is large, the restrictions are inconsistent with the data.

\textbf{Special case:} When $q = 1$ and $H_0: \beta_j = 0$, $F = t_j^2$ (the $F$-statistic is the square of the $t$-statistic).

\textbf{The overall $F$-test:} With $\mathbf{R} = [\mathbf{0} \mid \mathbf{I}_{k-1}]$ and $\mathbf{r} = \mathbf{0}_{k-1}$ (testing all slope coefficients jointly), the $F$-statistic tests whether \textit{any} regressor matters:

\[
F = \frac{R^2/(k-1)}{(1-R^2)/(n-k)}
\]

\medskip\hrule\medskip

\section{Chapter Summary}

\begin{itemize}
\item Under Assumptions 4.1--4.3 (including strict exogeneity), OLS is \textbf{unbiased}.
\item Under Assumptions 4.1--4.4 (adding spherical errors), the variance of OLS is $\sigma^2(\mathbf{X}^T\mathbf{X})^{-1}$, and OLS is the \textbf{Best Linear Unbiased Estimator (BLUE)} --- the \textbf{Gauss-Markov Theorem}.
\item The variance of $\hat{\beta}_j$ depends inversely on the variation in $x_j$ and positively on the error variance and multicollinearity ($R^2_j$).
\item $\hat{\sigma}^2 = SSR/(n-k)$ is an unbiased estimator of $\sigma^2$; the degrees of freedom are $n - k$.
\item Under normality (Assumption 4.5), $\hat{\boldsymbol{\beta}}$ is normally distributed, enabling exact $t$-tests and $F$-tests.
\item The $t$-statistic tests single hypotheses; the $F$-statistic tests joint hypotheses.
\end{itemize}

\medskip\hrule\medskip

\section{Key Terms}

\textbf{Strict exogeneity} --- $E[\mathbf{u} \mid \mathbf{X}] = \mathbf{0}$.

\textbf{Homoskedasticity} --- $E[u_i^2 \mid \mathbf{X}] = \sigma^2$ (constant error variance).

\textbf{Spherical errors} --- $E[\mathbf{u}\mathbf{u}^T \mid \mathbf{X}] = \sigma^2 \mathbf{I}_n$.

\textbf{Unbiasedness} --- $E[\hat{\boldsymbol{\beta}}] = \boldsymbol{\beta}$.

\textbf{Variance-covariance matrix of OLS} --- $\text{Var}(\hat{\boldsymbol{\beta}} \mid \mathbf{X}) = \sigma^2(\mathbf{X}^T\mathbf{X})^{-1}$.

\textbf{Gauss-Markov Theorem} --- OLS is BLUE (Best Linear Unbiased Estimator) under Assumptions 4.1--4.4.

\textbf{BLUE} --- Best (minimum variance) Linear Unbiased Estimator.

\textbf{Degrees of freedom} --- $n - k$; the number of observations minus the number of estimated parameters.

\textbf{$t$-statistic} --- $t_j = \hat{\beta}_j / \widehat{SE}(\hat{\beta}_j) \sim t_{n-k}$ under $H_0: \beta_j = 0$.

\textbf{$F$-statistic} --- statistic for joint hypothesis tests; $F \sim F_{q, n-k}$ under $H_0$.

\medskip\hrule\medskip

\section{Exercises}

\subsection{Conceptual Questions}

\textbf{4.1} Explain the difference between \textbf{unbiasedness} and \textbf{consistency}. Can an estimator be unbiased but inconsistent? Consistent but biased?

\textbf{4.2} Assumption 4.3 (strict exogeneity) requires $E[u_i \mid \mathbf{x}_1, \ldots, \mathbf{x}_n] = 0$, not just $E[u_i \mid \mathbf{x}_i] = 0$. Give an example where the weaker condition holds but strict exogeneity fails.

\textbf{4.3} Suppose we run OLS on data where $V(u_i) = \sigma^2 x_i^2$ (variance increases with $x_i$). The Gauss-Markov assumptions are violated. What specific assumption fails? Is OLS still unbiased? Is it still BLUE?

\textbf{4.4} Explain the degrees of freedom correction in $\hat{\sigma}^2 = SSR/(n-k)$. Why is dividing by $n$ instead of $n - k$ problematic?

\textbf{4.5} A researcher has a sample of $n = 1000$ and runs a regression with 3 regressors. She obtains $\hat{\beta}_2 = 0.5$ with $\widehat{SE}(\hat{\beta}_2) = 0.4$ ($t$-stat = 1.25). She concludes "the effect of $x_2$ is economically large but statistically insignificant." Comment on this interpretation.

\subsection{Analytical Questions}

\textbf{4.6} Under Assumptions 4.1--4.5, derive the distribution of $\hat{\boldsymbol{\beta}} \mid \mathbf{X}$.

\textbf{4.7} Prove that $E[SSR \mid \mathbf{X}] = (n-k)\sigma^2$, showing that $\hat{\sigma}^2 = SSR/(n-k)$ is unbiased.

\textbf{4.8} Let $c$ be a known $k \times 1$ vector. Find the variance of the linear combination $\mathbf{c}^T \hat{\boldsymbol{\beta}}$ under Assumptions 4.1--4.4. (This is useful for constructing confidence intervals for marginal effects or predictions.)

\textbf{4.9} Derive the $F$-statistic for the overall test of significance ($H_0$: all slope coefficients = 0) and show that it equals $\frac{R^2/(k-1)}{(1-R^2)/(n-k)}$.

\textbf{4.10} Suppose you add an irrelevant variable $z$ (with $\text{Cov}(z, u) = 0$ in the population) to a correctly specified regression. What happens to: (a) unbiasedness of existing coefficient estimates? (b) the variances of existing coefficient estimates? (c) $R^2$? (d) adjusted $R^2$?

\subsection{Applied Questions}

\textbf{4.11} Using the CPS data from previous labs, run a regression of log wages on education, experience, and experience$^2$. 

(a) Test $H_0: \beta_{\text{edu}} = 0$ at the 5\% level. Report the $t$-statistic, $p$-value, and conclusion.  
(b) Construct a 95\% confidence interval for $\beta_{\text{edu}}$.  
(c) Test the joint hypothesis $H_0: \beta_{\text{exp}} = \beta_{\text{exp}^2} = 0$.  
(d) What is the estimated peak of the experience-earnings profile?

\textbf{4.12} Conduct a Monte Carlo study to verify unbiasedness of OLS:
\begin{itemize}
\item Generate data: $y_i = 2 + 3x_i + u_i$, $x_i \sim N(0,1)$, $u_i \sim N(0,1)$, $n = 50$.
\item Repeat 1000 times; in each repetition, compute $\hat{\beta}_1$ and $\hat{\beta}_2$.
\item Report the mean of $\hat{\beta}_1$ and $\hat{\beta}_2$ across replications.
\item Show that the sampling distributions are centred on the true values.
\end{itemize}

\medskip\hrule\medskip

\section{R Lab 4: Statistical Properties and Inference}

\subsection{Objectives}
\begin{itemize}
\item Verify unbiasedness via Monte Carlo simulation
\item Compute standard errors and construct confidence intervals
\item Test hypotheses using t and F statistics
\item Explore the effects of multicollinearity
\end{itemize}

\subsection{Exercise 4.1: Verifying Unbiasedness via Monte Carlo}

\begin{lstlisting}
set.seed(42)
n_sims <- 2000
n      <- 100
beta_true <- c(2, 3, -1)   # intercept, beta_1, beta_2

# Storage
beta_hat_store <- matrix(NA, nrow = n_sims, ncol = 3)

for (s in 1:n_sims) {
  x1 <- rnorm(n)
  x2 <- rnorm(n)
  u  <- rnorm(n)
  y  <- beta_true[1] + beta_true[2] * x1 + beta_true[3] * x2 + u
  
  fit <- lm(y ~ x1 + x2)
  beta_hat_store[s, ] <- coef(fit)
}

# Are OLS estimates centred on true values?
colMeans(beta_hat_store)   # should be close to (2, 3, -1)
apply(beta_hat_store, 2, sd)  # sampling standard deviations

# Plot sampling distribution of beta_1
ggplot(data.frame(beta1 = beta_hat_store[, 2]), aes(x = beta1)) +
  geom_histogram(aes(y = ..density..), bins = 40, fill = "steelblue", alpha = 0.7) +
  geom_vline(xintercept = beta_true[2], color = "red", linewidth = 1) +
  stat_function(fun = dnorm,
                args = list(mean = beta_true[2], sd = sd(beta_hat_store[, 2])),
                color = "darkred", linewidth = 1) +
  labs(title = "Sampling Distribution of OLS Estimator",
       subtitle = paste0("True beta_1 = ", beta_true[2], " (red line); n_sims = ", n_sims),
       x = "Estimated beta_1", y = "Density") +
  theme_minimal()
\end{lstlisting}

\subsection{Exercise 4.2: Inference in R}

\begin{lstlisting}
library(AER)
data("CPS1988")
CPS1988 <- CPS1988 %>% mutate(log_wage = log(wage))

model <- lm(log_wage ~ education + experience + I(experience^2), data = CPS1988)
summary(model)

# Manual confidence interval
alpha <- 0.05
n <- nrow(CPS1988)
k <- length(coef(model))
t_crit <- qt(1 - alpha/2, df = n - k)

coef_table <- data.frame(
  Estimate = coef(model),
  SE = sqrt(diag(vcov(model)))
) %>%
  mutate(
    t_stat  = Estimate / SE,
    CI_lower = Estimate - t_crit * SE,
    CI_upper = Estimate + t_crit * SE,
    p_value  = 2 * pt(-abs(t_stat), df = n - k)
  )
print(round(coef_table, 4))
\end{lstlisting}

\subsection{Exercise 4.3: F-Test for Joint Significance}

\begin{lstlisting}
# Test H0: beta_experience = beta_experience^2 = 0
# Method 1: Built-in F-test via anova()
model_restricted <- lm(log_wage ~ education, data = CPS1988)
model_unrestricted <- model
anova(model_restricted, model_unrestricted)

# Method 2: Manual F-statistic
SSR_R <- sum(residuals(model_restricted)^2)
SSR_U <- sum(residuals(model_unrestricted)^2)
q <- 2   # number of restrictions
F_stat <- ((SSR_R - SSR_U) / q) / (SSR_U / (n - k))
p_val  <- pf(F_stat, df1 = q, df2 = n - k, lower.tail = FALSE)
cat("F-statistic:", round(F_stat, 3), "\n")
cat("p-value:", round(p_val, 4), "\n")
\end{lstlisting}

\subsection{Exercise 4.4: Effect of Multicollinearity}

\begin{lstlisting}
# Simulate data with varying levels of correlation between x1 and x2
set.seed(42)
n <- 200
beta_true <- c(2, 3, -1)

results <- map_dfr(c(0, 0.5, 0.8, 0.95, 0.99), function(rho) {
  sims <- replicate(1000, {
    SIGMA <- matrix(c(1, rho, rho, 1), 2, 2)
    X_sim <- MASS::mvrnorm(n, mu = c(0, 0), Sigma = SIGMA)
    u_sim <- rnorm(n)
    y_sim <- beta_true[1] + beta_true[2] * X_sim[,1] + beta_true[3] * X_sim[,2] + u_sim
    coef(lm(y_sim ~ X_sim))[2]  # coefficient on x1
  })
  data.frame(rho = rho, se_beta1 = sd(sims))
})

print(results)

ggplot(results, aes(x = rho, y = se_beta1)) +
  geom_line(color = "steelblue", linewidth = 1.2) +
  geom_point(size = 3, color = "steelblue") +
  labs(
    title = "Effect of Multicollinearity on Standard Errors",
    subtitle = "As correlation between x1 and x2 increases, SE of beta_1 explodes",
    x = "Correlation between x1 and x2",
    y = "Standard Deviation of OLS Estimator (Monte Carlo)"
  ) +
  theme_minimal()
\end{lstlisting}

\medskip\hrule\medskip

\textit{End of Chapter 4}

\medskip\hrule\medskip

\textbf{References}

Davidson, R., and MacKinnon, J. G. (2004). \textit{Econometric Theory and Methods}. Oxford University Press. Chapter 3.

Wooldridge, J. M. (2019). \textit{Introductory Econometrics} (7th ed.). Cengage. Chapters 4--5.

Angrist, J. D., and Pischke, J. (2009). \textit{Mostly Harmless Econometrics}. Princeton University Press. Chapter 3.

\chapter{Chapter 5: Asymptotic Theory and Robust Inference}

\medskip\hrule\medskip

\section{Chapter Overview}

Chapter 4 derived exact finite-sample results under the strong assumption that errors are normally distributed. In practice, economic data are rarely normally distributed: wages are right-skewed, binary outcomes are Bernoulli, error variances vary across observations. This chapter develops \textbf{asymptotic (large-sample) theory} --- statistical results that hold approximately when $n$ is large, without requiring normality or homoskedasticity.

We derive the \textbf{asymptotic distribution of OLS}, show how to estimate the asymptotic variance under \textbf{heteroskedasticity} (non-constant variance) and \textbf{within-cluster correlation}, and provide practical guidance on when and how to use robust standard errors. These are the standard errors used in virtually all applied economic research today.

\medskip\hrule\medskip

\section{Learning Objectives}

After completing this chapter, students will be able to:

\begin{enumerate}
\item Explain why asymptotic theory is needed in applied econometrics.
\item State and apply the \textbf{Law of Large Numbers (LLN)} and \textbf{Central Limit Theorem (CLT)} to OLS.
\item Prove \textbf{consistency} of OLS under weaker conditions than Gauss-Markov.
\item Derive the \textbf{asymptotic distribution} of the OLS estimator.
\item Explain what \textbf{heteroskedasticity} is and why it invalidates classical standard errors.
\item Compute and interpret \textbf{heteroskedasticity-robust (HC) standard errors}.
\item Explain \textbf{cluster-robust standard errors} and when clustering is necessary.
\item Provide practical guidance on standard error choice in applied work.
\end{enumerate}

\medskip\hrule\medskip

\section{5.1 Why Asymptotic Theory?}

The Gauss-Markov assumptions, particularly normality (Assumption 4.5), are often implausible:

\begin{itemize}
\item \textbf{Discrete outcomes:} Earnings rounded to the nearest dollar, binary indicators of employment, counts of accidents. Error distributions are clearly non-normal.
\item \textbf{Non-spherical errors:} Wages vary more for high earners (heteroskedasticity); students in the same class share correlated outcomes (clustering); annual firm data have autocorrelated shocks.
\item \textbf{Misspecified functional form:} When the CEF is non-linear, OLS residuals systematically deviate from normality.
\end{itemize}

When the classical assumptions fail, we need \textbf{asymptotic theory} --- results that hold as $n \to \infty$ under weaker conditions. The key tools are the \textbf{Law of Large Numbers} (sample averages converge to population means) and the \textbf{Central Limit Theorem} (standardized sample averages are approximately normal for large $n$).

\medskip\hrule\medskip

\section{5.2 Consistency of OLS}

\textbf{Definition 5.1 (Consistency):} An estimator $\hat{\boldsymbol{\beta}}_n$ is \textbf{consistent} for $\boldsymbol{\beta}$ if $\hat{\boldsymbol{\beta}}_n \xrightarrow{p} \boldsymbol{\beta}$ as $n \to \infty$. That is, for any $\epsilon > 0$:
\[
\lim_{n \to \infty} P(\|\hat{\boldsymbol{\beta}}_n - \boldsymbol{\beta}\| > \epsilon) = 0
\]

\textbf{Theorem 5.1 (Consistency of OLS):} Suppose the following conditions hold:
\begin{enumerate}
\item The model is linear: $y_i = \mathbf{x}_i^T \boldsymbol{\beta} + u_i$.
\item The data $\{(y_i, \mathbf{x}_i)\}$ are i.i.d.
\item \textbf{Predeterminedness:} $E[\mathbf{x}_i u_i] = \mathbf{0}$ (orthogonality condition).
\item $E[\mathbf{x}_i \mathbf{x}_i^T]$ is finite and non-singular.
\end{enumerate}

Then $\hat{\boldsymbol{\beta}}^{OLS} \xrightarrow{p} \boldsymbol{\beta}$.

\textbf{Proof:}

\[
\hat{\boldsymbol{\beta}} = \boldsymbol{\beta} + \left(\frac{1}{n}\mathbf{X}^T\mathbf{X}\right)^{-1} \left(\frac{1}{n}\mathbf{X}^T\mathbf{u}\right)
\]

By the LLN:
\[
\frac{1}{n}\mathbf{X}^T\mathbf{X} = \frac{1}{n}\sum_{i=1}^n \mathbf{x}_i \mathbf{x}_i^T \xrightarrow{p} E[\mathbf{x}_i \mathbf{x}_i^T] \equiv \mathbf{S}_{XX}
\]

\[
\frac{1}{n}\mathbf{X}^T\mathbf{u} = \frac{1}{n}\sum_{i=1}^n \mathbf{x}_i u_i \xrightarrow{p} E[\mathbf{x}_i u_i] = \mathbf{0}
\]

By Slutsky's theorem and the continuous mapping theorem:

\[
\hat{\boldsymbol{\beta}} \xrightarrow{p} \boldsymbol{\beta} + \mathbf{S}_{XX}^{-1} \mathbf{0} = \boldsymbol{\beta} \quad \square
\]

\subsection{5.2.1 Consistency vs. Unbiasedness}

\textbf{Key distinction:} Consistency only requires the weaker condition $E[\mathbf{x}_i u_i] = \mathbf{0}$ (predeterminedness / zero covariance), not strict exogeneity $E[u_i \mid \mathbf{X}] = \mathbf{0}$.

\begin{itemize}
\item Strict exogeneity $\Rightarrow$ predeterminedness, but NOT vice versa.
\item Predeterminedness is enough for consistency; strict exogeneity is needed for unbiasedness.
\end{itemize}

This means OLS can be \textbf{consistent but biased in small samples} --- the bias vanishes as $n \to \infty$.

\textbf{Practical implication:} In many panel data models (with lagged dependent variables), OLS is biased in small samples but consistent in large samples. For small panels, this bias can be substantial.

\medskip\hrule\medskip

\section{5.3 Asymptotic Normality and the Sandwich Variance}

\textbf{Theorem 5.2 (Asymptotic Normality of OLS):} Under the i.i.d. assumptions of Theorem 5.1, plus $E[\mathbf{x}_i \mathbf{x}_i^T u_i^2]$ finite:

\[
\sqrt{n}(\hat{\boldsymbol{\beta}} - \boldsymbol{\beta}) \xrightarrow{d} N\!\left(\mathbf{0}, \mathbf{S}_{XX}^{-1} \boldsymbol{\Omega} \mathbf{S}_{XX}^{-1}\right)
\]

where $\mathbf{S}_{XX} = E[\mathbf{x}_i \mathbf{x}_i^T]$ and $\boldsymbol{\Omega} = E[\mathbf{x}_i \mathbf{x}_i^T u_i^2] = E[u_i^2 \mathbf{x}_i \mathbf{x}_i^T]$.

\textbf{Proof:} The key step is applying the \textbf{Multivariate Central Limit Theorem (MCLT)} to $\frac{1}{\sqrt{n}}\mathbf{X}^T\mathbf{u}$:

\[
\frac{1}{\sqrt{n}}\mathbf{X}^T\mathbf{u} = \frac{1}{\sqrt{n}}\sum_{i=1}^n \mathbf{x}_i u_i \xrightarrow{d} N(\mathbf{0}, \boldsymbol{\Omega})
\]

By the LLN, $\frac{1}{n}\mathbf{X}^T\mathbf{X} \xrightarrow{p} \mathbf{S}_{XX}$. By Slutsky's theorem:

\[
\sqrt{n}(\hat{\boldsymbol{\beta}} - \boldsymbol{\beta}) = \left(\frac{1}{n}\mathbf{X}^T\mathbf{X}\right)^{-1} \frac{1}{\sqrt{n}}\mathbf{X}^T\mathbf{u} \xrightarrow{d} \mathbf{S}_{XX}^{-1} \cdot N(\mathbf{0}, \boldsymbol{\Omega}) = N(\mathbf{0}, \mathbf{S}_{XX}^{-1}\boldsymbol{\Omega}\mathbf{S}_{XX}^{-1}) \quad \square
\]

\begin{figure}[htbp]
  \centering
  \begin{tikzpicture}
    \begin{axis}[
      width=11cm, height=6.5cm,
      xlabel={$\sqrt{n}(\hat{\beta} - \beta_0)$},
      ylabel={Density},
      xmin=-4, xmax=4,
      ymin=0, ymax=0.5,
      axis lines=left,
      tick style={draw=none},
      xtick={0},
      xticklabels={0},
      ytick=\empty,
      legend style={at={(0.97,0.95)}, anchor=north east, font=\small},
    ]
      % Small n distribution (wider, less normal)
      \addplot[thick, gray!60, domain=-4:4, samples=100]
        {0.38*exp(-0.5*(x)^2/1.8^2)/sqrt(2*pi*1.8^2)*4};

      % Medium n
      \addplot[thick, econblue!60, domain=-4:4, samples=100]
        {0.4*exp(-0.5*(x)^2/1.2^2)/sqrt(2*pi*1.2^2)*4};

      % Large n (approaching normal)
      \addplot[very thick, econblue, domain=-4:4, samples=100]
        {0.42*exp(-0.5*(x)^2)/sqrt(2*pi)*4};

      % Standard normal (limiting)
      \addplot[thick, red!70!black, dashed, domain=-4:4, samples=100]
        {exp(-0.5*x^2)/sqrt(2*pi)};

      \addlegendentry{Small $n$}
      \addlegendentry{Moderate $n$}
      \addlegendentry{Large $n$}
      \addlegendentry{$\mathcal{N}(0,V_\beta)$ (limiting)}

      \node[font=\footnotesize, econblue] at (axis cs:-2.5, 0.38) {$\sqrt{n}(\hat{\beta}-\beta_0)\xrightarrow{d}\mathcal{N}(0,V_\beta)$};
    \end{axis}
  \end{tikzpicture}
  \caption{Asymptotic normality of the OLS estimator. As the sample size $n$ increases, the scaled and centered OLS estimator $\sqrt{n}(\hat{\beta} - \beta_0)$ converges in distribution to a normal random variable. The limiting distribution $\mathcal{N}(0, V_\beta)$ (dashed red) provides the basis for inference in large samples.}
  \label{fig:ch05_asymptotic_normality}
\end{figure}

\subsection{5.3.1 The Sandwich Variance Estimator}

The matrix $\mathbf{S}_{XX}^{-1}\boldsymbol{\Omega}\mathbf{S}_{XX}^{-1}$ is called the \textbf{sandwich variance} (or \textbf{Huber-White} variance). In finite samples:

\[
\text{Var}(\hat{\boldsymbol{\beta}}) \approx \frac{1}{n} \mathbf{S}_{XX}^{-1} \boldsymbol{\Omega} \mathbf{S}_{XX}^{-1} = (\mathbf{X}^T\mathbf{X})^{-1}\left[\mathbf{X}^T\boldsymbol{\Omega}\mathbf{X}\right](\mathbf{X}^T\mathbf{X})^{-1}
\]

\textbf{Under homoskedasticity:} $E[u_i^2 \mid \mathbf{x}_i] = \sigma^2$, so $\boldsymbol{\Omega} = \sigma^2 \mathbf{S}_{XX}$, and the sandwich variance simplifies to:

\[
\frac{\sigma^2}{n} \mathbf{S}_{XX}^{-1} \approx \sigma^2 (\mathbf{X}^T\mathbf{X})^{-1}
\]

This is the classical variance formula --- consistent with the Gauss-Markov result.

\textbf{Under heteroskedasticity:} $E[u_i^2 \mid \mathbf{x}_i] = \sigma_i^2 \neq \sigma^2$, the classical formula gives the wrong variance. We need to estimate $\boldsymbol{\Omega}$ consistently.

\medskip\hrule\medskip

\section{5.4 Heteroskedasticity-Robust Standard Errors}

\textbf{Definition 5.2 (Heteroskedasticity):} Errors are heteroskedastic if $V(u_i \mid \mathbf{x}_i)$ varies across observations.

\textbf{Example 5.1:} In a wage regression, it is common for high earners to have greater dispersion in wages than low earners --- the variance of $u_i$ may be an increasing function of predicted wages.

\subsection{5.4.1 Consequences of Ignoring Heteroskedasticity}

If we use the classical variance formula $\hat{\sigma}^2(\mathbf{X}^T\mathbf{X})^{-1}$ under heteroskedasticity:
\begin{itemize}
\item OLS is still \textbf{consistent} (consistency does not require spherical errors).
\item But standard errors are \textbf{wrong}: they may be too large or too small.
\item $t$-tests and $F$-tests are \textbf{invalid}: rejection rates under the null are not $\alpha$.
\end{itemize}

\subsection{5.4.2 The Heteroskedasticity-Robust (HC) Estimator}

\textbf{White (1980)} proposed a consistent estimator of the sandwich variance that does not require spherical errors:

\[
\widehat{\text{Var}}_{HC}(\hat{\boldsymbol{\beta}}) = (\mathbf{X}^T\mathbf{X})^{-1} \left(\sum_{i=1}^n \hat{u}_i^2 \mathbf{x}_i \mathbf{x}_i^T\right) (\mathbf{X}^T\mathbf{X})^{-1}
\]

This is a sample analogue of the sandwich formula, replacing $E[u_i^2 \mathbf{x}_i \mathbf{x}_i^T]$ with $\frac{1}{n}\sum_{i=1}^n \hat{u}_i^2 \mathbf{x}_i \mathbf{x}_i^T$.

\textbf{Why does this work?} By the consistency of OLS, $\hat{u}_i \approx u_i$ for large $n$. So $\hat{u}_i^2 \mathbf{x}_i \mathbf{x}_i^T$ is a consistent estimator of $E[u_i^2 \mathbf{x}_i \mathbf{x}_i^T]$.

\subsection{5.4.3 Small-Sample Corrections (HC1, HC2, HC3)}

In small samples, $\hat{u}_i^2$ is a biased downward estimate of $u_i^2$ (because OLS fits the residuals). Several corrections are used in practice:

\begin{center}
\begin{tabular}{@{} l l l @{}}
\toprule
\textbf{Version} & \textbf{Formula} & \textbf{Motivation} \\
\midrule
HC0 & $\hat{u}_i^2$ & Original White (1980) estimator \\
HC1 & $\frac{n}{n-k}\hat{u}_i^2$ & Degrees-of-freedom correction \\
HC2 & $\frac{\hat{u}_i^2}{1 - h_{ii}}$ & Corrects for leverage \\
HC3 & $\frac{\hat{u}_i^2}{(1 - h_{ii})^2}$ & Better in very small samples \\
\bottomrule
\end{tabular}
\end{center}

where $h_{ii} = [\mathbf{P_X}]_{ii}$ is the $i$-th diagonal of the hat matrix (leverage of observation $i$).

\textbf{In practice:} Use HC3 for small samples ($n < 250$), HC1 for larger samples. The \texttt{sandwich} and \texttt{lmtest} packages in R provide all variants.

\begin{keybox}{The Rule of Thumb:** In modern applied econometrics, always report heteroskedasticity-robust standard errors unless you have a strong theoretical reason to believe errors are homoskedastic.}
\end{keybox}

\medskip\hrule\medskip

\section{5.5 Cluster-Robust Standard Errors}

\subsection{5.5.1 When Clustering Matters}

\textbf{Clustering arises when observations are not independent.} Consider:

\[
y_{ig} = \mathbf{x}_{ig}^T \boldsymbol{\beta} + u_{ig}
\]

where $i$ indexes individuals within group (cluster) $g$. If individuals within the same group share common unobservables (e.g., students in the same school share a teacher quality shock), then $\text{Cov}(u_{ig}, u_{jg}) \neq 0$ for $i \neq j$.

\textbf{The variance-covariance structure is block-diagonal:}

\[
E[\mathbf{u}\mathbf{u}^T] = \begin{pmatrix} \boldsymbol{\Sigma}_1 & \mathbf{0} & \cdots & \mathbf{0} \\ \mathbf{0} & \boldsymbol{\Sigma}_2 & \cdots & \mathbf{0} \\ \vdots & & \ddots & \vdots \\ \mathbf{0} & \mathbf{0} & \cdots & \boldsymbol{\Sigma}_G \end{pmatrix}
\]

where $\boldsymbol{\Sigma}_g = E[\mathbf{u}_g \mathbf{u}_g^T]$ allows arbitrary within-cluster correlation.

\subsection{5.5.2 The Cluster-Robust Sandwich Estimator}

\textbf{Definition 5.3 (Cluster-Robust SE):}

\[
\widehat{\text{Var}}_{CR}(\hat{\boldsymbol{\beta}}) = (\mathbf{X}^T\mathbf{X})^{-1} \left(\sum_{g=1}^G \mathbf{X}_g^T \hat{\mathbf{u}}_g \hat{\mathbf{u}}_g^T \mathbf{X}_g\right) (\mathbf{X}^T\mathbf{X})^{-1}
\]

where $\mathbf{X}_g$ and $\hat{\mathbf{u}}_g$ are the design matrix and residuals for cluster $g$.

\textbf{Key difference from HC:} HC sums over individual observations; CR sums over clusters. Within each cluster, we allow arbitrary covariance in residuals.

\subsection{5.5.3 Practical Guidance on When to Cluster}

Clustering is needed when:

\begin{enumerate}
\item \textbf{Treatment is at a higher level than outcomes.} E.g., a minimum wage law applies to a state, but outcomes are individual workers. Cluster at the state level.
\end{enumerate}

\begin{enumerate}
\item \textbf{Panel data with repeated observations.} Each individual appears multiple times; errors are autocorrelated over time. Cluster at the individual level.
\end{enumerate}

\begin{enumerate}
\item \textbf{Assignment of treatment is clustered.} E.g., a program assigns villages to treatment; outcomes are individual household members. Cluster at the village level.
\end{enumerate}

\textbf{Rule of thumb:} Cluster at the level at which treatment varies (or at which observations are correlated).

\textbf{Warning --- few clusters:} Cluster-robust SEs rely on $G \to \infty$. With few clusters ($G < 20$), they are unreliable. In such cases, use the \textbf{wild cluster bootstrap} (see Cameron and Miller, 2015).

\textbf{Example 5.2 (Card and Krueger, 1994):} In their study of the effect of a New Jersey minimum wage increase on fast-food employment, Card and Krueger cluster at the restaurant level --- each restaurant is surveyed before and after the policy change, so two observations per restaurant are correlated.

\medskip\hrule\medskip

\section{5.6 When to Cluster? Practical Guidance}

A common question from students: "Should I cluster?" The answer depends on the data structure, not the size of the coefficient:

\begin{center}
\begin{tabular}{@{} l l l @{}}
\toprule
\textbf{Situation} & \textbf{Cluster?} & \textbf{Level} \\
\midrule
Random sample, i.i.d. observations & No & --- \\
Survey with multi-stage sampling & Yes & Primary sampling unit \\
Difference-in-differences & Yes & State/region/treatment unit \\
Panel data & Yes & Individual/firm \\
Cross-section with geographic variation & Often & Geographic unit \\
RCT with cluster randomization & Yes & Randomization unit \\
\bottomrule
\end{tabular}
\end{center}

\textbf{Never} choose whether to cluster based on whether it changes your results. Cluster based on the structure of the data.

\medskip\hrule\medskip

\section{5.7 Empirical Illustration: Wages and Education Revisited}

Consider the CPS wage regression:

\[
\log(\text{wage}_i) = \beta_1 + \beta_2 \cdot \text{educ}_i + \beta_3 \cdot \text{exp}_i + \beta_4 \cdot \text{exp}_i^2 + u_i
\]

\textbf{Table 5.1: OLS Estimates with Different Standard Errors}

\begin{center}
\begin{tabular}{@{} l l l l @{}}
\toprule
\textbf{} & \textbf{Classical SE} & \textbf{HC1 SE} & \textbf{HC3 SE} \\
\midrule
Education & 0.0709 (0.0010) & 0.0709 (0.0011) & 0.0709 (0.0011) \\
Experience & 0.0343 (0.0010) & 0.0343 (0.0013) & 0.0343 (0.0013) \\
Exp$^{2}$ & $-$0.0005 (0.000) & $-$0.0005 (0.000) & $-$0.0005 (0.000) \\
\bottomrule
\end{tabular}
\end{center}

Standard errors in parentheses.

In this large sample, classical and robust SEs are similar because $n$ is large. The key point: in datasets with clear heteroskedasticity patterns, or with clustering, the differences can be substantial.

\medskip\hrule\medskip

\section{Chapter Summary}

\begin{itemize}
\item Asymptotic theory delivers results about $\hat{\boldsymbol{\beta}}$ under weaker conditions than finite-sample theory.
\item OLS is \textbf{consistent} under predeterminedness ($E[\mathbf{x}_i u_i] = \mathbf{0}$) --- weaker than strict exogeneity.
\item The \textbf{asymptotic distribution} of $\sqrt{n}(\hat{\boldsymbol{\beta}} - \boldsymbol{\beta})$ is normal with sandwich variance $\mathbf{S}_{XX}^{-1}\boldsymbol{\Omega}\mathbf{S}_{XX}^{-1}$.
\item \textbf{Heteroskedasticity} means $V(u_i \mid \mathbf{x}_i)$ is not constant. It does not bias OLS but invalidates classical SEs.
\item \textbf{White (HC) standard errors} provide consistent inference under heteroskedasticity of unknown form.
\item \textbf{Cluster-robust standard errors} are needed when observations within groups are correlated; they cluster the "meat" of the sandwich at the group level.
\item In modern applied work, \textbf{always report robust standard errors} unless there is a strong reason to believe errors are spherical.
\end{itemize}

\medskip\hrule\medskip

\section{Key Terms}

\textbf{Consistency} --- $\hat{\boldsymbol{\beta}} \xrightarrow{p} \boldsymbol{\beta}$ as $n \to \infty$.

\textbf{Predeterminedness} --- $E[\mathbf{x}_i u_i] = \mathbf{0}$; weaker than strict exogeneity.

\textbf{Asymptotic normality} --- $\sqrt{n}(\hat{\boldsymbol{\beta}} - \boldsymbol{\beta}) \xrightarrow{d} N(\mathbf{0}, \mathbf{V})$ for some variance matrix $\mathbf{V}$.

\textbf{Sandwich variance} --- $\mathbf{S}_{XX}^{-1}\boldsymbol{\Omega}\mathbf{S}_{XX}^{-1}$; the asymptotic variance of OLS.

\textbf{Heteroskedasticity} --- non-constant error variance: $V(u_i \mid \mathbf{x}_i) \neq \sigma^2$.

\textbf{Heteroskedasticity-robust (HC) SE} --- consistent standard errors that allow for heteroskedasticity.

\textbf{Cluster-robust SE} --- consistent standard errors that allow for within-cluster correlation.

\textbf{Wild cluster bootstrap} --- a resampling method for inference with few clusters.

\medskip\hrule\medskip

\section{Exercises}

\subsection{Conceptual Questions}

\textbf{5.1} Explain the difference between \textbf{unbiasedness} and \textbf{consistency}. Give an example of an estimator that is consistent but biased in small samples.

\textbf{5.2} Why does heteroskedasticity not cause OLS to be biased or inconsistent? What does it affect?

\textbf{5.3} A researcher runs a difference-in-differences regression using state-level panel data (50 states, 10 years). A colleague says "You have 500 observations, so clustering doesn't matter." Evaluate this claim.

\textbf{5.4} What is the "sandwich" structure of the heteroskedasticity-robust variance estimator? Why is it called a sandwich?

\textbf{5.5} You are analyzing school data where 30 students are observed in each of 50 schools. Should you cluster by student, school, or not at all? Justify your answer.

\subsection{Analytical Questions}

\textbf{5.6} Prove that under predeterminedness and i.i.d. sampling, OLS is consistent.

\textbf{5.7} Suppose $V(u_i \mid x_i) = \sigma^2 x_i^2$. Derive the true variance of OLS in the bivariate regression of $y_i$ on $x_i$. Show it differs from the classical formula $\sigma^2/\text{SST}_x$.

\textbf{5.8} Show that when $V(u_i \mid \mathbf{x}_i) = \sigma^2$ (homoskedasticity), the HC0 sandwich estimator is NOT equal to the classical estimator, but converges to the same limit as $n \to \infty$.

\textbf{5.9} Consider panel data $y_{it} = \mathbf{x}_{it}^T\boldsymbol{\beta} + \alpha_i + \epsilon_{it}$ where $\alpha_i$ is an individual-specific shock and $\epsilon_{it}$ is i.i.d. If you run OLS pooling all observations without fixed effects, $\hat{\boldsymbol{\beta}}$ is generally biased. But even if you control for fixed effects (making $\hat{\boldsymbol{\beta}}$ consistent), why might you still need to cluster standard errors at the individual level?

\subsection{Applied Questions}

\textbf{5.10} Using the CPS data, estimate the wage regression with (a) classical SEs, (b) HC1 SEs, and (c) HC3 SEs. Report all three and compare. When do they differ?

\textbf{5.11} In R, demonstrate via Monte Carlo simulation that the classical $t$-test has the wrong size (rejection rate $\neq 5\%$ under $H_0$) when errors are heteroskedastic, but the HC-robust $t$-test maintains the correct size.

\medskip\hrule\medskip

\section{R Lab 5: Robust Standard Errors in Practice}

\subsection{Setup}

\begin{lstlisting}
library(tidyverse)
library(sandwich)   # for vcovHC
library(lmtest)    # for coeftest
library(AER)

data("CPS1988")
CPS1988 <- CPS1988 %>% mutate(log_wage = log(wage))

model <- lm(log_wage ~ education + experience + I(experience^2), data = CPS1988)
\end{lstlisting}

\subsection{Exercise 5.1: Compare Standard Errors}

\begin{lstlisting}
# Classical SEs
coeftest(model)

# HC0 (White 1980)
coeftest(model, vcov = vcovHC(model, type = "HC0"))

# HC1 (degrees-of-freedom correction)
coeftest(model, vcov = vcovHC(model, type = "HC1"))

# HC3 (best for small samples)
coeftest(model, vcov = vcovHC(model, type = "HC3"))
\end{lstlisting}

\subsection{Exercise 5.2: Detecting Heteroskedasticity}

\begin{lstlisting}
# Residuals vs fitted values plot
df_resid <- data.frame(
  fitted   = fitted(model),
  residuals = residuals(model)
)

ggplot(df_resid, aes(x = fitted, y = residuals)) +
  geom_point(alpha = 0.2, color = "steelblue") +
  geom_hline(yintercept = 0, color = "red") +
  geom_smooth(method = "loess", se = FALSE, color = "darkred") +
  labs(
    title = "Residuals vs. Fitted Values",
    subtitle = "Pattern indicates heteroskedasticity",
    x = "Fitted Values", y = "Residuals"
  ) +
  theme_minimal()

# Breusch-Pagan test for heteroskedasticity
library(lmtest)
bptest(model)
\end{lstlisting}

\subsection{Exercise 5.3: Cluster-Robust Standard Errors}

\begin{lstlisting}
library(sandwich)

# Simulate panel-like clustering by state (using ethnicity as a proxy cluster)
# In real applications, cluster = geographic unit or treatment unit

# Cluster by 'ethnicity' (as illustrative proxy)
vcov_cluster <- vcovCL(model, cluster = ~ CPS1988$ethnicity)
coeftest(model, vcov = vcov_cluster)
\end{lstlisting}

\subsection{Exercise 5.4: Monte Carlo --- Size of Classical vs. Robust t-tests under Heteroskedasticity}

\begin{lstlisting}
set.seed(42)
n_sims <- 5000
n      <- 200
reject_classical <- numeric(n_sims)
reject_robust    <- numeric(n_sims)

for (s in 1:n_sims) {
  x <- rnorm(n)
  u <- rnorm(n, sd = abs(x))   # heteroskedastic: V(u|x) = x^2
  y <- 1 + 0 * x + u            # H0: beta = 0 is TRUE

  fit <- lm(y ~ x)
  
  # Classical t-stat
  t_classic <- summary(fit)$coefficients["x", "t value"]
  reject_classical[s] <- abs(t_classic) > 1.96
  
  # Robust t-stat
  t_robust <- coeftest(fit, vcov = vcovHC(fit, "HC1"))["x", "t value"]
  reject_robust[s] <- abs(t_robust) > 1.96
}

cat("Classical test rejection rate:", round(mean(reject_classical), 3),
    "(should be ~0.05)\n")
cat("Robust test rejection rate:   ", round(mean(reject_robust), 3),
    "(should be ~0.05)\n")
\end{lstlisting}

\textbf{Expected finding:} The classical test rejects too often (size > 5\%) under heteroskedasticity; the robust test maintains approximately the correct size.

\medskip\hrule\medskip

\textit{End of Chapter 5}

\medskip\hrule\medskip

\textbf{References}

White, H. (1980). A heteroskedasticity-consistent covariance matrix estimator and a direct test for heteroskedasticity. \textit{Econometrica}, 48(4), 817--838.

Davidson, R., and MacKinnon, J. G. (2004). \textit{Econometric Theory and Methods}. Chapter 3.

Angrist, J. D., and Pischke, J. (2009). \textit{Mostly Harmless Econometrics}. Chapter 8.

Cameron, A. C., and Miller, D. L. (2015). A practitioner's guide to cluster-robust inference. \textit{Journal of Human Resources}, 50(2), 317--372.

Cunningham, S. (2021). \textit{Causal Inference: The Mixtape}. Chapter on regression.


%% ===========================================================
%%  PART II: THE REGRESSION TOOLKIT
%% ===========================================================
\part{The Regression Toolkit}

\chapter{Chapter 6: The Frisch-Waugh-Lovell Theorem and Omitted Variable Bias}

\medskip\hrule\medskip

\section{Chapter Overview}

Chapter 3 introduced the Frisch-Waugh-Lovell (FWL) Theorem as a result about the algebra of OLS. This chapter makes the theorem earn its keep: it is the key tool for understanding \textbf{omitted variable bias (OVB)} --- the most important source of endogeneity in applied work. We derive the OVB formula, discuss its direction and magnitude, and explain the difference between bad and good controls. We also cover proxy variables and the concept of "bad controls" --- variables that, when included, can actually increase bias rather than reduce it.

\medskip\hrule\medskip

\section{Learning Objectives}

After completing this chapter, students will be able to:

\begin{enumerate}
\item Formally state and prove the \textbf{Frisch-Waugh-Lovell (FWL) Theorem}.
\item Derive the \textbf{omitted variable bias (OVB) formula} and use it to sign the direction of bias.
\item Explain when a control variable reduces bias vs. when it introduces bias ("bad controls").
\item Distinguish between "bad controls" and "good controls."
\item Explain the concept of a \textbf{proxy variable} and when it is useful.
\item Apply these concepts to empirical examples in wages, health, and education.
\end{enumerate}

\medskip\hrule\medskip

\section{6.1 Partitioned Regression and the FWL Theorem}

We proved the FWL Theorem in Chapter 3. Here we develop its implications more carefully.

\textbf{Setup:} Partition the regression:

\[
\mathbf{y} = \mathbf{X}_1 \boldsymbol{\beta}_1 + \mathbf{X}_2 \boldsymbol{\beta}_2 + \mathbf{u}
\]

where $\mathbf{X}_1$ is $n \times k_1$ ("control variables") and $\mathbf{X}_2$ is $n \times k_2$ ("variables of interest").

\textbf{Theorem 6.1 (Frisch-Waugh-Lovell):} The OLS estimate $\hat{\boldsymbol{\beta}}_2$ from the full regression equals the OLS estimate from regressing $\tilde{\mathbf{y}} = \mathbf{M}_{X_1}\mathbf{y}$ on $\tilde{\mathbf{X}}_2 = \mathbf{M}_{X_1}\mathbf{X}_2$:

\[
\hat{\boldsymbol{\beta}}_2 = (\tilde{\mathbf{X}}_2^T \tilde{\mathbf{X}}_2)^{-1} \tilde{\mathbf{X}}_2^T \tilde{\mathbf{y}}
\]

\textbf{Economic interpretation of "residualizing":}

\begin{itemize}
\item $\tilde{\mathbf{y}} = \mathbf{M}_{X_1}\mathbf{y}$: the part of $y$ that is NOT linearly predictable from $\mathbf{X}_1$.
\item $\tilde{\mathbf{X}}_2 = \mathbf{M}_{X_1}\mathbf{X}_2$: the part of $\mathbf{X}_2$ that is NOT linearly predictable from $\mathbf{X}_1$.
\end{itemize}

The FWL regression asks: \textit{among the remaining variation in $\mathbf{X}_2$ after removing what $\mathbf{X}_1$ can predict, how much of the remaining variation in $y$ does it explain?}

This is exactly what we mean by "controlling for $\mathbf{X}_1$" --- we remove the linear influence of $\mathbf{X}_1$ on both $y$ and $\mathbf{X}_2$ before asking about the relationship between $\mathbf{X}_2$ and $y$.

\begin{figure}[htbp]
  \centering
  \begin{tikzpicture}[
    box/.style={draw=econblue, fill=econblue!8, rounded corners, thick, minimum width=4cm, minimum height=1.1cm, align=center, font=\small},
    arr/.style={-{Stealth}, thick, econblue},
    scale=0.9
  ]
    % Step 1
    \node[box, minimum width=5cm] (step1a) at (0, 0) {Regress $Y_i$ on $\mathbf{W}_i$\\Get residuals $\tilde{Y}_i$};
    \node[box, minimum width=5cm] (step1b) at (6.5, 0) {Regress $D_i$ on $\mathbf{W}_i$\\Get residuals $\tilde{D}_i$};

    % Step 2
    \node[box, minimum width=5cm] (step2) at (3.25, -2) {Regress $\tilde{Y}_i$ on $\tilde{D}_i$};

    % Arrows
    \draw[arr] (step1a) -- (step2) node[midway, left, font=\footnotesize] {};
    \draw[arr] (step1b) -- (step2);

    % Result
    \node[draw=econblue, fill=econblue!15, thick, rounded corners, minimum width=5cm, minimum height=0.9cm, align=center, font=\small\bfseries] (result) at (3.25, -3.8) {Coefficient on $\tilde{D}_i$ $=$ OLS coefficient on $D_i$ in full regression};
    \draw[arr] (step2) -- (result);

    \node[font=\footnotesize, gray, align=center] at (3.25, -5.0) {FWL: $\hat{\beta}_D^{\text{full}} = (\tilde{\mathbf{D}}^T\tilde{\mathbf{D}})^{-1}\tilde{\mathbf{D}}^T\tilde{\mathbf{Y}}$};
    \node[font=\footnotesize, gray, align=center] at (3.25, -5.5) {where tildes denote residuals after projecting out $\mathbf{W}$};
  \end{tikzpicture}
  \caption{The Frisch-Waugh-Lovell (FWL) Theorem as a two-step procedure. The coefficient on $D_i$ in the full regression (controlling for $\mathbf{W}_i$) equals the coefficient from regressing the ``residualized'' outcome $\tilde{Y}_i$ on the ``residualized'' treatment $\tilde{D}_i$. This demonstrates that OLS ``partials out'' the controls.}
  \label{fig:ch06_fwl}
\end{figure}

\medskip\hrule\medskip

\section{6.2 Omitted Variable Bias}

The most important implication of the FWL Theorem is for \textbf{omitted variable bias (OVB)}: what happens to estimated coefficients when we leave out a variable that belongs in the model?

\subsection{6.2.1 The Setup}

Suppose the true model is:

\[
y_i = \beta_1 + \beta_2 x_i + \beta_3 z_i + u_i \tag{Long}
\]

but we instead estimate the \textbf{short regression}:

\[
y_i = \tilde{\beta}_1 + \tilde{\beta}_2 x_i + v_i \tag{Short}
\]

What is the relationship between $\hat{\tilde{\beta}}_2$ (from the short regression) and $\hat{\beta}_2$ (from the long regression)?

\subsection{6.2.2 The OVB Formula}

\textbf{Theorem 6.2 (Omitted Variable Bias):} The short-regression coefficient satisfies:

\[
\hat{\tilde{\beta}}_2 = \hat{\beta}_2 + \hat{\beta}_3 \cdot \hat{\delta}_{zx}
\]

where $\hat{\delta}_{zx}$ is the \textbf{auxiliary regression coefficient} from regressing $z_i$ on $x_i$ (the slope from OLS of $z$ on $x$).

\textbf{In population (probability limits):}

\[
\text{plim}(\hat{\tilde{\beta}}_2) = \beta_2 + \beta_3 \cdot \delta_{zx}
\]

where $\delta_{zx} = \text{Cov}(x_i, z_i)/V(x_i)$.

The \textbf{omitted variable bias} is:

\[
\text{OVB} = \underbrace{\beta_3}_{\text{effect of } z \text{ on } y} \cdot \underbrace{\delta_{zx}}_{\text{relationship between } z \text{ and } x}
\]

\textbf{Proof:} By OLS algebra and the FWL Theorem. The short regression estimates:

\[
\hat{\tilde{\beta}}_2 = \frac{\text{Cov}_s(x_i, y_i)}{V_s(x_i)}
\]

Substituting $y_i = \beta_1 + \beta_2 x_i + \beta_3 z_i + u_i$ and using properties of covariance:

\[
= \beta_2 + \beta_3 \frac{\text{Cov}(x_i, z_i)}{V(x_i)} + \frac{\text{Cov}(x_i, u_i)}{V(x_i)} = \beta_2 + \beta_3 \delta_{zx} \quad \square
\]

(where $\text{Cov}(x_i, u_i) = 0$ by the long regression assumption.)

\subsection{6.2.3 Signing the Bias}

The OVB formula has two factors. To determine the sign of the bias:

\[
\text{sign(OVB)} = \text{sign}(\beta_3) \times \text{sign}(\delta_{zx})
\]

\begin{center}
\begin{tabular}{@{} l l l l @{}}
\toprule
\textbf{Effect of $z$ on $y$ ($\beta_3$)} & \textbf{Correlation of $x$ and $z$ ($\delta_{zx}$)} & \textbf{OVB} & \textbf{Short regression} \\
\midrule
Positive & Positive & Positive (upward) & Overestimates $\beta_2$ \\
Positive & Negative & Negative (downward) & Underestimates $\beta_2$ \\
Negative & Positive & Negative (downward) & Underestimates $\beta_2$ \\
Negative & Negative & Positive (upward) & Overestimates $\beta_2$ \\
\bottomrule
\end{tabular}
\end{center}

\begin{figure}[htbp]
  \centering
  \begin{tikzpicture}[scale=1.0,
    cell/.style={draw=econblue!40, minimum width=3.5cm, minimum height=1.1cm, align=center, font=\small},
    header/.style={draw=econblue, fill=econblue!20, minimum width=3.5cm, minimum height=0.9cm, align=center, font=\small\bfseries}
  ]
    % Header row
    \node[header, minimum width=3.2cm] at (0, 0) {};
    \node[header] at (3.2, 0) {$\rho_{D\tilde{W}} > 0$};
    \node[header] at (6.4, 0) {$\rho_{D\tilde{W}} < 0$};

    % Left header column
    \node[header, minimum width=3.2cm] at (0, -1.2) {$\delta_W > 0$};
    \node[header, minimum width=3.2cm] at (0, -2.4) {$\delta_W < 0$};

    % Inner cells
    \node[cell, fill=red!15] at (3.2, -1.2) {$\hat{\delta} > \delta$\\(upward bias)};
    \node[cell, fill=econblue!12] at (6.4, -1.2) {$\hat{\delta} < \delta$\\(downward bias)};
    \node[cell, fill=econblue!12] at (3.2, -2.4) {$\hat{\delta} < \delta$\\(downward bias)};
    \node[cell, fill=red!15] at (6.4, -2.4) {$\hat{\delta} > \delta$\\(upward bias)};

    % Labels
    \node[font=\footnotesize, gray] at (3.2, -3.3) {$\text{OVB} = \delta_W \cdot \rho_{D\tilde{W}}$};
    \node[font=\footnotesize, align=center, gray] at (-1.8, -1.8) {Effect of\\omitted\\variable};
    \node[font=\footnotesize, align=center, gray] at (4.8, 0.9) {Correlation of omitted variable with $D$};
  \end{tikzpicture}
  \caption{Omitted Variable Bias (OVB) direction table. The sign of the bias in $\hat{\delta}$ depends on (1) whether the omitted variable $W$ has a positive or negative effect on $Y$ (holding $D$ fixed), and (2) whether $W$ is positively or negatively correlated with the treatment $D$. Bias is upward (red) when these signs agree, downward (blue) when they disagree.}
  \label{fig:ch06_ovb_direction}
\end{figure}

\textbf{Example 6.1 (Returns to Education):} True model: log wages = $\beta_1 + \beta_2 \cdot \text{educ} + \beta_3 \cdot \text{ability} + u$.

If we omit ability:
\begin{itemize}
\item $\beta_3 > 0$: ability increases wages.
\item $\delta_{zx} > 0$: ability is positively correlated with education.
\item OVB $> 0$: the short regression \textbf{overestimates} the return to education.
\end{itemize}

This is the classical "ability bias" problem in the returns to schooling literature. Simply regressing wages on years of schooling produces an upward-biased estimate because more able people both get more education and earn more.

\textbf{Example 6.2 (Class Size and Test Scores):} True model: test scores = $\beta_1 + \beta_2 \cdot \text{classsize} + \beta_3 \cdot \text{income} + u$.

If we omit family income:
\begin{itemize}
\item $\beta_3 > 0$: higher income families have children with higher scores.
\item $\delta_{zx} < 0$: higher income districts have smaller classes.
\item OVB $< 0$: the short regression \textbf{overestimates} the negative effect of class size (makes classes seem more harmful than they are).
\end{itemize}

\subsection{6.2.4 The Multivariate OVB Formula}

With multiple omitted variables $\mathbf{z}_i = (z_{1i}, \ldots, z_{mi})^T$:

\[
\text{plim}(\hat{\tilde{\beta}}_2) = \beta_2 + \boldsymbol{\beta}_z^T \boldsymbol{\delta}_{zx}
\]

where $\boldsymbol{\beta}_z$ is the vector of coefficients on $\mathbf{z}$ in the long regression and $\boldsymbol{\delta}_{zx}$ is the vector of auxiliary regression coefficients. OVB can be positive or negative depending on the signs of all $\beta_{z_j} \delta_{z_j x}$ terms.

\medskip\hrule\medskip

\section{6.3 Bad Controls}

A natural response to OVB is to "include more controls." But including too many controls --- or the wrong controls --- can introduce new biases. Variables that seem like good controls can actually be \textbf{bad controls}.

\subsection{6.3.1 What Is a Bad Control?}

\textbf{Definition 6.1 (Bad Control):} A variable $z_i$ is a \textbf{bad control} if it is itself affected by the treatment variable $x_i$ (i.e., $z_i$ is a "mediator" on the causal pathway from $x_i$ to $y_i$, or if it shares unobserved determinants with $y_i$ beyond those that caused $z_i$ to be correlated with $x_i$).

\textbf{Canonical example (Angrist and Pischke, 2009):} Suppose we want to estimate the effect of college education ($x$) on wages ($y$). Occupation ($z$) is a potential control variable. But:
\begin{itemize}
\item College education affects the occupation a worker enters.
\item Occupation affects wages.
\item Occupation is thus a \textbf{mediator} on the causal path: Education $\to$ Occupation $\to$ Wages.
\end{itemize}

If we control for occupation, we are "blocking" part of the causal pathway through which education affects wages. The coefficient on education will \textbf{understate} the total effect of education on wages.

\subsection{6.3.2 The Bias from Controlling for a Mediator}

When $z_i$ is a mediator:

\[
y_i = \beta_1 + \beta_2^{direct} x_i + \beta_3 z_i + u_i
\]

where $\beta_2^{direct}$ is the \textit{direct} effect of $x$ on $y$ (not through $z$). But:

\[
z_i = \gamma_1 + \gamma_2 x_i + \epsilon_i
\]

The \textbf{total} effect of $x$ on $y$ is $\beta_2^{direct} + \beta_3 \gamma_2$ (direct + indirect through $z$). Controlling for $z$ recovers only the direct effect.

\textbf{Is this wrong?} It depends on the question. If you want the \textbf{total} causal effect of education on wages, you should NOT control for occupation. If you want the \textbf{direct} effect (not mediated by occupation), you should control for it.

\subsection{6.3.3 Selection into the Control Variable}

A more subtle problem arises when the control variable $z_i$ has its own unobserved determinants $w_i$ that also affect $y_i$:

\begin{itemize}
\item $x$ $\to$ $y$ (the causal path we want)
\item $x$ $\to$ $z$ (treatment affects the control)
\item $w$ $\to$ $z$ (unobserved factor affects control)
\item $w$ $\to$ $y$ (same unobserved factor affects outcome)
\end{itemize}

Controlling for $z$ can open up a "backdoor" correlation between $x$ and $y$ through $w$. Including $z$ introduces bias from $w$ that wasn't there before.

\textbf{Example:} Suppose we want the effect of gender ($x$) on wages ($y$). A naive control might be job title ($z$). But:
\begin{itemize}
\item Gender affects job title (discrimination in promotion).
\item $w$ = unobserved productivity affects both job title and wages.
\item Conditioning on job title opens a spurious path: Gender $\to$ Job Title $\leftarrow$ Productivity $\to$ Wages.
\item This introduces bias rather than removing it.
\end{itemize}

\subsection{6.3.4 Rules of Thumb for Control Variables}

\textbf{Good controls:} Pre-treatment variables that affect both $x$ and $y$ but are not affected by $x$. These reduce bias from confounding.

\textbf{Bad controls:} Variables that are:
\begin{enumerate}
\item Caused by the treatment $x$ (mediators)
\item Share unobserved determinants with $y$ beyond their relationship with $x$
\item Selected into by the same process that creates the endogeneity problem
\end{enumerate}

\begin{keybox}{Practical Advice:\textbf{ Think carefully about the causal structure before choosing controls. A useful framework is the }Directed Acyclic Graph (DAG)**, which visually represents causal relationships and helps identify good and bad controls.}
\end{keybox}

\medskip\hrule\medskip

\section{6.4 Proxy Variables}

When the true omitted variable $z_i$ is unobserved, we may have access to a \textbf{proxy} --- a variable $p_i$ that is correlated with $z_i$ and satisfies certain conditions.

\textbf{Definition 6.2 (Proxy Variable):} $p_i$ is a valid proxy for $z_i$ if:
\begin{enumerate}
\item $p_i$ is correlated with $z_i$: $\text{Cov}(p_i, z_i) \neq 0$.
\item $p_i$ is uncorrelated with the structural error $u_i$ conditional on $z_i$: $E[u_i \mid x_i, z_i, p_i] = E[u_i \mid x_i, z_i]$.
\end{enumerate}

Condition (2) says $p_i$ affects $y_i$ only through $z_i$ --- once we control for the underlying unobservable $z_i$, the proxy $p_i$ adds no additional information.

\textbf{Example:} In the wage-education model, IQ score ($p_i$) may serve as a proxy for cognitive ability ($z_i$). If IQ only affects wages through its relationship with ability, and doesn't directly affect wages, it satisfies the proxy conditions.

\textbf{Limitation:} The OLS estimate using a proxy will generally not recover the true structural coefficient on $x_i$, but it reduces OVB compared to omitting the proxy entirely.

\medskip\hrule\medskip

\section{Chapter Summary}

\begin{itemize}
\item The \textbf{FWL Theorem} shows that OLS slope coefficients measure partial effects after "residualizing" on all other controls.
\item \textbf{Omitted variable bias} arises when an important determinant of $y$ is both correlated with included regressors and omitted from the model.
\item The \textbf{OVB formula} is $\text{OVB} = \beta_3 \cdot \delta_{zx}$, the product of the omitted variable's effect on $y$ and its correlation with the included regressor.
\item "More controls are better" is NOT always true. \textbf{Bad controls} --- mediators or variables that share unobserved determinants with $y$ --- can introduce new biases.
\item \textbf{Proxy variables} can reduce OVB when the true omitted variable is unobservable but an imperfect proxy is available.
\end{itemize}

\medskip\hrule\medskip

\section{Key Terms}

\textbf{Omitted Variable Bias (OVB)} --- bias in an estimated coefficient when a variable that belongs in the model is omitted.

\textbf{Long regression} --- the correctly specified model including all relevant variables.

\textbf{Short regression} --- the misspecified model omitting one or more relevant variables.

\textbf{Auxiliary regression} --- the regression of the omitted variable on the included regressors.

\textbf{OVB formula} --- $\text{plim}(\hat{\tilde{\beta}}_2) = \beta_2 + \beta_3 \delta_{zx}$.

\textbf{Bad control} --- a variable that, when included, introduces bias rather than reducing it.

\textbf{Mediator} --- a variable on the causal pathway from the treatment to the outcome.

\textbf{Proxy variable} --- an observable variable correlated with an unobservable omitted variable.

\textbf{Directed Acyclic Graph (DAG)} --- a visual representation of assumed causal relationships.

\medskip\hrule\medskip

\section{Exercises}

\subsection{Conceptual Questions}

\textbf{6.1} State the OVB formula. What are its two components? Explain the direction of OVB for the ability bias in the returns-to-education literature.

\textbf{6.2} Explain why "including more controls" does not always reduce bias. Give an example of a bad control.

\textbf{6.3} You are studying the effect of a job training program ($x$) on earnings ($y$). A colleague suggests controlling for whether the participant found a new job after the program ($z$). Is this a good or bad control? Explain.

\textbf{6.4} Define a proxy variable. What conditions must it satisfy? Give an example from economics.

\textbf{6.5} In a regression of wages on education and experience, a researcher finds that adding "number of previous employers" reduces the education coefficient from 0.10 to 0.06. What does this suggest about the omitted variable bias in the original regression?

\subsection{Analytical Questions}

\textbf{6.6} True model: $y_i = 2 + 3x_i + (-1)z_i + u_i$, where $\text{Cov}(x,z)/V(x) = 0.5$.
(a) What is the short-regression coefficient on $x$?
(b) Is the estimate biased upward or downward relative to the true $\beta_2 = 3$?

\textbf{6.7} Suppose we have the true model $y_i = \beta_0 + \beta_1 x_{1i} + \beta_2 x_{2i} + \beta_3 x_{3i} + u_i$, but we estimate $y_i = \gamma_0 + \gamma_1 x_{1i} + \gamma_2 x_{2i} + v_i$ (omitting $x_{3i}$). Derive the formula for $\text{plim}(\hat{\gamma}_1)$ and $\text{plim}(\hat{\gamma}_2)$.

\textbf{6.8} Consider a regression of health status on income and region. Suppose region is correlated with both income and health (richer regions have better health care). Describe what happens to the income coefficient if you (a) include region as a control; (b) omit region.

\subsection{Applied Questions}

\textbf{6.9} Using the CPS data:
(a) Estimate the return to education with no controls.
(b) Add experience as a control. Does the education coefficient change? In which direction? Is this consistent with OVB?
(c) Add ethnicity as a control. Does the education coefficient change? What does this imply?

\textbf{6.10} Conduct the following experiment in R:
\begin{itemize}
\item Generate data: $y_i = 1 + 2x_i + 3z_i + u_i$, $z_i = 0.5 x_i + \epsilon_i$.
\item Estimate the short regression (omitting $z$) and the long regression (including $z$).
\item Verify the OVB formula numerically.
\end{itemize}

\medskip\hrule\medskip

\section{R Lab 6: Omitted Variable Bias and the FWL Theorem}

\subsection{Exercise 6.1: Demonstrating OVB}

\begin{lstlisting}
library(tidyverse)
set.seed(2024)
n <- 1000

# True parameters
beta_0 <- 1; beta_x <- 2; beta_z <- 3

# Simulate omitted variable correlated with x
x <- rnorm(n)
z <- 0.5 * x + rnorm(n)  # delta_zx ~= 0.5
u <- rnorm(n)
y <- beta_0 + beta_x * x + beta_z * z + u

# Short regression (omitting z)
short_reg <- lm(y ~ x)
cat("Short regression coefficient on x:", round(coef(short_reg)["x"], 3), "\n")

# Long regression (including z)
long_reg <- lm(y ~ x + z)
cat("Long regression coefficient on x:", round(coef(long_reg)["x"], 3), "\n")

# Auxiliary regression
aux_reg <- lm(z ~ x)
delta_zx <- coef(aux_reg)["x"]
cat("delta_zx (z on x):", round(delta_zx, 3), "\n")

# OVB formula
ovb_formula <- beta_z * delta_zx
cat("OVB formula prediction:", round(ovb_formula, 3), "\n")
cat("Actual OVB:", round(coef(short_reg)["x"] - coef(long_reg)["x"], 3), "\n")
\end{lstlisting}

\subsection{Exercise 6.2: Visualizing Bad Controls}

\begin{lstlisting}
# Simulate: education -> occupation -> wages; ability -> wages
set.seed(42)
n <- 2000

ability   <- rnorm(n)
education <- 0.5 * ability + rnorm(n)
# occupation is a mediator: caused by education AND ability
occupation_index <- 0.4 * education + 0.3 * ability + rnorm(n)
wages <- 0.5 * education + 0.3 * ability + 0.2 * occupation_index + rnorm(n)

# True total effect of education on wages (should be approx 0.5 + 0.2*0.4 = 0.58)
cat("True total effect:", 0.5 + 0.2 * 0.4, "\n")

# Without controlling for occupation
reg1 <- lm(wages ~ education)
cat("OLS without occupation:", round(coef(reg1)["education"], 3), "\n")

# Controlling for occupation (BAD: occupation is a mediator)
reg2 <- lm(wages ~ education + occupation_index)
cat("OLS with occupation (bad control):", round(coef(reg2)["education"], 3), "\n")

# Controlling for ability (GOOD: pre-treatment confounder)
reg3 <- lm(wages ~ education + ability)
cat("OLS with ability (good control):", round(coef(reg3)["education"], 3), "\n")
\end{lstlisting}

\subsection{Exercise 6.3: Verifying FWL Numerically}

\begin{lstlisting}
library(AER)
data("CPS1988")
CPS1988 <- CPS1988 %>% mutate(log_wage = log(wage))

# Full regression
full_model <- lm(log_wage ~ education + experience + I(experience^2) + ethnicity,
                  data = CPS1988)

# FWL: residualize log_wage and education on other controls
M_controls_y    <- residuals(lm(log_wage   ~ experience + I(experience^2) + ethnicity, data = CPS1988))
M_controls_educ <- residuals(lm(education  ~ experience + I(experience^2) + ethnicity, data = CPS1988))

fwl_model <- lm(M_controls_y ~ M_controls_educ)

cat("Full OLS education coefficient:", round(coef(full_model)["education"], 6), "\n")
cat("FWL education coefficient:     ", round(coef(fwl_model)["M_controls_educ"], 6), "\n")
\end{lstlisting}

\medskip\hrule\medskip

\textit{End of Chapter 6}

\medskip\hrule\medskip

\textbf{References}

Angrist, J. D., and Pischke, J. (2009). \textit{Mostly Harmless Econometrics}. Chapters 3 and 6.

Davidson, R., and MacKinnon, J. G. (2004). \textit{Econometric Theory and Methods}. Chapter 2.

Pearl, J. (2009). \textit{Causality: Models, Reasoning, and Inference} (2nd ed.). Cambridge University Press.

Wooldridge, J. M. (2019). \textit{Introductory Econometrics} (7th ed.). Chapter 3.

\chapter{Chapter 7: Regression with Special Regressors: Dummies, Interactions, and Non-linearities}

\medskip\hrule\medskip

\section{Chapter Overview}

Real economic relationships rarely look like straight lines. This chapter extends the linear regression toolkit to handle discrete and categorical variables, interactions between variables, polynomial terms, and logarithmic transformations. We also discuss heterogeneous treatment effects and weighted regression. These tools are essential for applied work and appear in virtually every empirical paper.

\medskip\hrule\medskip

\section{Learning Objectives}

After completing this chapter, students will be able to:

\begin{enumerate}
\item Use \textbf{binary dummy variables} to capture discrete categories.
\item Avoid the \textbf{dummy variable trap} (perfect multicollinearity with the intercept).
\item Include \textbf{interaction terms} to allow effects to vary across groups.
\item Interpret \textbf{log-linear}, \textbf{log-log}, and \textbf{linear-log} regression specifications.
\item Use \textbf{polynomial regression} to capture non-linear relationships.
\item Interpret \textbf{regression with saturated specifications} as estimating the CEF exactly.
\item Understand how heterogeneous treatment effects are reflected in regression.
\item Apply \textbf{weighted least squares (WLS)} when error variances are known.
\end{enumerate}

\medskip\hrule\medskip

\section{7.1 Binary and Categorical Variables}

\subsection{7.1.1 The Binary Dummy Variable}

A \textbf{binary (dummy) variable} $D_i \in \{0, 1\}$ captures membership in a category:

\[
D_i = \begin{cases} 1 & \text{if condition holds} \\ 0 & \text{otherwise} \end{cases}
\]

Examples: employed/unemployed, male/female, urban/rural, treated/not treated.

\textbf{Regression with a single dummy:}

\[
y_i = \alpha + \rho D_i + u_i
\]

Interpretation:
\begin{itemize}
\item $\alpha = E[y_i \mid D_i = 0]$ = mean outcome for the reference group.
\item $\rho = E[y_i \mid D_i = 1] - E[y_i \mid D_i = 0]$ = difference in mean outcomes between the two groups.
\end{itemize}

The OLS estimate $\hat{\rho}$ is simply the difference in sample means: $\bar{y}_1 - \bar{y}_0$.

\begin{figure}[htbp]
  \centering
  \begin{tikzpicture}
    \begin{axis}[
      width=9cm, height=6.5cm,
      xlabel={$X$ (continuous variable)},
      ylabel={$Y$},
      xmin=0, xmax=6, ymin=0.5, ymax=5.5,
      axis lines=left,
      tick style={draw=none},
      xtick=\empty, ytick=\empty,
      legend style={at={(0.05,0.95)}, anchor=north west, font=\small},
    ]
      % Group 0 line (D=0)
      \addplot[very thick, econblue, domain=0.5:5.5] {0.6*x + 0.8};
      % Group 1 line (D=1), shifted up
      \addplot[very thick, red!70!black, domain=0.5:5.5] {0.6*x + 2.0};

      % Brace for alpha
      \draw[<->, thick, gray] (axis cs:0.7, 1.22) -- (axis cs:0.7, 2.42)
        node[midway, left, font=\small, gray] {$\alpha$};

      % Data points group 0
      \addplot[only marks, mark=o, mark size=2pt, color=econblue]
        coordinates {(1,1.4)(2,2.0)(3,2.6)(4,3.5)(5,3.8)};
      % Data points group 1
      \addplot[only marks, mark=square*, mark size=2pt, color=red!70!black]
        coordinates {(1,2.7)(2,3.4)(3,3.8)(4,4.5)(5,5.1)};

      \legend{$D_i=0$ (control), $D_i=1$ (treated)};
      \node[font=\footnotesize, gray] at (axis cs:4, 1.2) {$Y_i = \beta_0 + \beta_1 X_i + \alpha D_i + u_i$};
    \end{axis}
  \end{tikzpicture}
  \caption{A dummy variable shifts the intercept by $\alpha$ while keeping the slope constant. The two regression lines for treated ($D_i=1$, red) and untreated ($D_i=0$, blue) are parallel, separated by the treatment effect $\alpha$.}
  \label{fig:ch07_dummy_parallel}
\end{figure}

\subsection{7.1.2 Categorical Variables with $J$ Categories}

For a variable with $J$ categories (e.g., education levels: less than high school, high school, some college, college, graduate), create $J-1$ dummy variables, omitting one \textbf{reference category}:

\[
y_i = \alpha + \sum_{j=2}^{J} \rho_j D_{ij} + u_i
\]

where $D_{ij} = 1$ if individual $i$ is in category $j$.

\textbf{Interpretation:} Each $\rho_j$ is the mean difference in $y$ between category $j$ and the reference category (category 1), conditional on other controls.

\textbf{The Dummy Variable Trap:} Including $J$ dummies plus an intercept creates perfect multicollinearity: $\sum_{j=1}^J D_{ij} = 1 = $ intercept. Always omit one category.

\textbf{Example 7.1 (Education and Wages):} Create dummies for education groups:
\begin{itemize}
\item $D_{\text{HS},i}$ = 1 if high school graduate (reference category: less than HS)
\item $D_{\text{SC},i}$ = 1 if some college
\item $D_{\text{BA},i}$ = 1 if bachelor's degree
\item $D_{\text{MA},i}$ = 1 if graduate degree
\end{itemize}

Coefficients $\hat{\rho}_{\text{HS}}, \hat{\rho}_{\text{SC}}, \hat{\rho}_{\text{BA}}, \hat{\rho}_{\text{MA}}$ are the wage premia for each group relative to those with less than a high school diploma.

\medskip\hrule\medskip

\section{7.2 Interaction Terms}

\subsection{7.2.1 Why Interactions?}

An \textbf{interaction term} allows the effect of one variable to depend on the value of another. Without interactions, we impose that the slope coefficient on $x$ is the same for all individuals. This is often too restrictive.

\subsection{7.2.2 Interaction Between Two Dummy Variables}

\[
y_i = \alpha + \beta_1 D_{1i} + \beta_2 D_{2i} + \beta_3 (D_{1i} \times D_{2i}) + u_i
\]

The expected value of $y_i$ for each combination:

\begin{center}
\begin{tabular}{@{} l l l @{}}
\toprule
\textbf{$D_1$} & \textbf{$D_2$} & \textbf{$E[y_i]$} \\
\midrule
0 & 0 & $\alpha$ \\
1 & 0 & $\alpha + \beta_1$ \\
0 & 1 & $\alpha + \beta_2$ \\
1 & 1 & $\alpha + \beta_1 + \beta_2 + \beta_3$ \\
\bottomrule
\end{tabular}
\end{center}

The interaction coefficient $\beta_3$ is the \textbf{additional effect} of having both $D_1 = 1$ and $D_2 = 1$ beyond the sum of their individual effects.

\textbf{Example 7.2:} Regress wages on a female dummy ($D_1$), a college dummy ($D_2$), and their interaction. The interaction coefficient captures whether the college wage premium is different for women than for men.

\subsection{7.2.3 Interaction Between a Dummy and a Continuous Variable}

\[
y_i = \alpha + \beta_1 D_i + \beta_2 x_i + \beta_3 (D_i \times x_i) + u_i
\]

Interpretation:
\begin{itemize}
\item When $D_i = 0$: $E[y_i \mid x_i, D_i=0] = \alpha + \beta_2 x_i$ (slope = $\beta_2$).
\item When $D_i = 1$: $E[y_i \mid x_i, D_i=1] = (\alpha + \beta_1) + (\beta_2 + \beta_3) x_i$ (slope = $\beta_2 + \beta_3$).
\end{itemize}

The interaction allows for different slopes in the two groups. $\beta_3$ is the \textbf{differential slope} for the group with $D_i = 1$.

\textbf{Example 7.3:} Allow the return to education to differ by gender. Regress log wages on education, a female dummy, and the education $\times$ female interaction. The interaction coefficient gives the female-male difference in returns to education.

\begin{figure}[htbp]
  \centering
  \begin{tikzpicture}
    \begin{axis}[
      width=9cm, height=6.5cm,
      xlabel={$X$},
      ylabel={$Y$},
      xmin=0, xmax=6, ymin=0, ymax=6,
      axis lines=left,
      tick style={draw=none},
      xtick=\empty, ytick=\empty,
      legend style={at={(0.05,0.95)}, anchor=north west, font=\small},
    ]
      % Group 0 (D=0): slope beta1
      \addplot[very thick, econblue, domain=0.5:5.5] {0.4*x + 0.5};
      % Group 1 (D=1): slope beta1 + gamma
      \addplot[very thick, red!70!black, domain=0.5:5.5] {0.9*x + 0.8};

      % Interaction label
      \node[font=\footnotesize, gray, align=center] at (axis cs:3.5, 1.0) {$Y_i = \beta_0 + \beta_1 X_i + \alpha D_i + \gamma (D_i \cdot X_i) + u_i$};

      % Slope difference annotation
      \draw[->, thick, gray] (axis cs:4.5, 4.85) arc (280:240:1.2)
        node[left, font=\footnotesize, gray] {slope $= \beta_1+\gamma$};
      \draw[->, thick, gray] (axis cs:4.2, 2.18) arc (100:130:1.2)
        node[left, font=\footnotesize, gray] {slope $= \beta_1$};

      \legend{$D_i=0$, $D_i=1$}
    \end{axis}
  \end{tikzpicture}
  \caption{An interaction term $D_i \cdot X_i$ allows the slope of $X$ to differ between groups. When $\gamma \neq 0$, the two regression lines are \emph{not} parallel: the effect of $X$ on $Y$ depends on treatment status.}
  \label{fig:ch07_interaction}
\end{figure}

\subsection{7.2.4 The Saturated Regression Model}

A \textbf{saturated} model includes a separate indicator for every possible combination of the categorical variables. For two binary variables $D_1$ and $D_2$:

\[
y_i = \alpha + \beta_1 D_{1i} + \beta_2 D_{2i} + \beta_3 (D_{1i} \times D_{2i}) + u_i
\]

This is saturated for two binary variables. The OLS estimates of $E[y \mid D_1, D_2]$ for each cell are exactly the sample cell means. \textbf{The saturated model estimates the CEF exactly} (for discrete regressors).

\medskip\hrule\medskip

\section{7.3 Polynomial and Logarithmic Transformations}

\subsection{7.3.1 Polynomial Regression}

To capture non-linearities in a continuous variable, include higher-order polynomial terms:

\[
y_i = \beta_0 + \beta_1 x_i + \beta_2 x_i^2 + \cdots + \beta_p x_i^p + u_i
\]

This is still a \textbf{linear} regression (linear in the parameters), so OLS applies directly.

\textbf{Application: Experience-Earnings Profile}

In the labour economics literature, the relationship between experience and wages is typically hump-shaped --- wages rise rapidly at first, then more slowly, eventually declining near retirement. A quadratic specification captures this:

\[
\log(\text{wage}_i) = \beta_0 + \beta_1 \text{exp}_i + \beta_2 \text{exp}_i^2 + \ldots
\]

The peak of the experience-earnings profile occurs at:

\[
\text{exp}^* = -\frac{\hat{\beta}_1}{2\hat{\beta}_2}
\]

\textbf{Example 7.4:} Typical estimates: $\hat{\beta}_1 \approx 0.04$, $\hat{\beta}_2 \approx -0.0007$. Peak at $\text{exp}^* \approx 0.04/(2 \times 0.0007) \approx 28$ years of experience.

\subsection{7.3.2 Logarithmic Transformations}

Logarithmic transformations are ubiquitous in economics. They:
\begin{itemize}
\item Reduce skewness in right-tailed distributions (wages, income, firm size).
\item Allow interpretation of coefficients as \textbf{percentage changes}.
\item Provide a natural way to estimate \textbf{elasticities}.
\end{itemize}

\textbf{Three standard specifications:}

\begin{center}
\begin{tabular}{@{} l l l @{}}
\toprule
\textbf{Model} & \textbf{Equation} & \textbf{Interpretation of $\hat{\beta}$} \\
\midrule
\textbf{Log-linear} & $\log(y_i) = \beta_0 + \beta_1 x_i + u_i$ & 1-unit $\uparrow$ in $x$ $\to$ $100\hat{\beta}_1$\% change in $y$ \\
\textbf{Linear-log} & $y_i = \beta_0 + \beta_1 \log(x_i) + u_i$ & 1\% $\uparrow$ in $x$ $\to$ $\hat{\beta}_1/100$ unit change in $y$ \\
\textbf{Log-log} & $\log(y_i) = \beta_0 + \beta_1 \log(x_i) + u_i$ & 1\% $\uparrow$ in $x$ $\to$ $\hat{\beta}_1$\% change in $y$ (elasticity) \\
\bottomrule
\end{tabular}
\end{center}

\textbf{Example 7.5 (Mincer Earnings Equation):} The classic Mincer (1974) equation uses a log-linear specification:

\[
\log(\text{wage}_i) = \beta_0 + \beta_1 \cdot \text{schooling}_i + \beta_2 \cdot \text{exp}_i + \beta_3 \cdot \text{exp}_i^2 + u_i
\]

The coefficient $\hat{\beta}_1 \approx 0.10$ implies each additional year of schooling raises wages by approximately 10\%.

\textbf{Approximation note:} $\log(y_1) - \log(y_0) \approx (y_1 - y_0)/y_0$ for small changes. For large changes (>20\%), the approximation becomes less accurate and we should use $\exp(\hat{\beta}_1) - 1$ for the exact percentage change.

\medskip\hrule\medskip

\section{7.4 Heterogeneous Treatment Effects and Regression}

\subsection{7.4.1 What Regression Recovers with Heterogeneous Effects}

In Chapter 1, we assumed constant treatment effects: $Y_{1i} - Y_{0i} = \rho$ for all $i$. In reality, treatment effects vary across individuals. What does OLS estimate in this case?

Consider:

\[
Y_i = \alpha + \rho_i D_i + \eta_i
\]

where $\rho_i = Y_{1i} - Y_{0i}$ is the individual treatment effect and $\eta_i = Y_{0i} - E[Y_{0i}]$.

If $D_i$ is randomly assigned, one can show that OLS estimates:

\[
\hat{\rho}^{OLS} \xrightarrow{p} E[\rho_i] + \frac{\text{Cov}(\rho_i, D_i(1-D_i))}{E[D_i(1-D_i)]}
\]

Under random assignment ($D_i$ independent of $\rho_i$), this simplifies to:

\[
\hat{\rho}^{OLS} \xrightarrow{p} E[\rho_i] = \text{ATE}
\]

OLS recovers the ATE under random assignment, even with heterogeneous effects.

\subsection{7.4.2 Regression Weighted Average of Effects}

More generally, Angrist (1998) shows that with heterogeneous effects and a continuous treatment, OLS computes a \textbf{variance-weighted average} of causal effects:

\[
\hat{\beta}^{OLS} \xrightarrow{p} \frac{E[\omega_i \cdot \beta_i]}{E[\omega_i]}
\]

where $\omega_i = (x_i - E[x_i])^2$ and $\beta_i$ is the individual-specific treatment effect. Observations with more variation in $x_i$ receive higher weight.

\textbf{Implication:} OLS gives greater weight to the treatment effect for individuals whose $x_i$ is far from the mean. This is not necessarily the same as the average effect for the typical individual.

\medskip\hrule\medskip

\section{7.5 Weighting and Weighted Least Squares}

\subsection{7.5.1 When to Weight?}

Weighted regression assigns observation $i$ weight $w_i > 0$ and minimizes:

\[
\sum_{i=1}^n w_i (y_i - \mathbf{x}_i^T \boldsymbol{\beta})^2
\]

The WLS estimator is:

\[
\hat{\boldsymbol{\beta}}^{WLS} = (\mathbf{X}^T \mathbf{W} \mathbf{X})^{-1} \mathbf{X}^T \mathbf{W} \mathbf{y}
\]

where $\mathbf{W} = \text{diag}(w_1, \ldots, w_n)$.

\textbf{Two main justifications for weighting:}

\begin{enumerate}
\item \textbf{Known heteroskedasticity:} If $V(u_i) = \sigma^2/w_i$, then WLS with weights $w_i$ is equivalent to OLS on a transformed model with homoskedastic errors, making WLS the efficient estimator (BLUE under heteroskedasticity).
\end{enumerate}

\begin{enumerate}
\item \textbf{Survey weights:} Complex surveys oversample certain groups. Survey weights adjust for differential sampling probabilities to ensure estimates are representative of the population.
\end{enumerate}

\textbf{Example 7.6:} Suppose $V(u_i) = \sigma^2 x_i$ (variance proportional to $x_i$). Use WLS with $w_i = 1/x_i$. This makes the transformed errors homoskedastic: $u_i^* = u_i/\sqrt{x_i}$ has constant variance $\sigma^2$.

\subsection{7.5.2 Debate on Weighting in Applied Work}

A subtle issue (Solon, Haider, and Wooldridge, 2015): survey weights adjust for sampling design. But should we use them to estimate conditional means?

\begin{itemize}
\item If the model is \textbf{correctly specified} (CEF is truly linear with the same coefficients for all groups), weighting is unnecessary --- unweighted OLS is more efficient.
\item If the model is \textbf{misspecified} (the CEF is non-linear or slopes vary across groups), weighted estimates may better describe the population-representative effect.
\end{itemize}

There is no universal answer. Practical guidance: report both weighted and unweighted estimates; if they differ substantially, investigate why.

\medskip\hrule\medskip

\section{Chapter Summary}

\begin{itemize}
\item \textbf{Dummy variables} capture categorical distinctions; always omit one reference category to avoid the dummy variable trap.
\item \textbf{Interaction terms} allow effects to vary across groups or with the level of another variable.
\item The \textbf{saturated model} (full set of interactions for discrete regressors) estimates the CEF exactly at each cell.
\item \textbf{Polynomial terms} capture non-linearities; logarithmic transformations facilitate elasticity interpretation.
\item With heterogeneous treatment effects, OLS estimates a \textbf{variance-weighted average} of effects.
\item \textbf{Weighted least squares (WLS)} is appropriate under known heteroskedasticity or with survey weights.
\end{itemize}

\medskip\hrule\medskip

\section{Key Terms}

\textbf{Dummy variable} --- a binary $\{0,1\}$ variable indicating group membership.

\textbf{Reference category} --- the omitted category in a set of dummy variables.

\textbf{Dummy variable trap} --- perfect multicollinearity from including all category dummies with an intercept.

\textbf{Interaction term} --- the product of two regressors, allowing the effect of one to vary with the other.

\textbf{Saturated model} --- a model that includes all possible combinations of categorical variables, estimating the CEF exactly.

\textbf{Log-linear model} --- $\log(y) = \mathbf{x}^T\boldsymbol{\beta} + u$; coefficients are approximate percentage changes.

\textbf{Elasticity} --- coefficient from a log-log model; measures percentage change in $y$ per 1\% change in $x$.

\textbf{Mincer earnings equation} --- the classic log-linear wage regression with schooling and a quadratic in experience.

\textbf{Weighted Least Squares (WLS)} --- minimizes a weighted sum of squared residuals.

\medskip\hrule\medskip

\section{Exercises}

\subsection{Conceptual Questions}

\textbf{7.1} You include a gender dummy (1 = female) in a wage regression. The coefficient is $-$0.15. Interpret this coefficient. What does it mean? What doesn't it tell you?

\textbf{7.2} Explain the dummy variable trap. Why is it a problem? How do you avoid it?

\textbf{7.3} A researcher runs: $\log(\text{wage}) = \beta_0 + \beta_1 \text{educ} + \beta_2 \text{female} + \beta_3 (\text{educ} \times \text{female}) + u$. Interpret each coefficient.

\textbf{7.4} In the quadratic experience-earnings model, derive the formula for the experience level at which predicted wages are maximized.

\textbf{7.5} "A log-log regression coefficient is the elasticity of $y$ with respect to $x$." Is this exactly correct? When is it an approximation?

\subsection{Analytical Questions}

\textbf{7.6} Suppose $y_i = \alpha + \beta_1 D_{1i} + \beta_2 D_{2i} + \beta_3 (D_{1i} \times D_{2i}) + u_i$ where $D_1$ = male (0=female), $D_2$ = college (0=no college). 

Write the expected value of $y$ for each of the four groups: (female, no college), (female, college), (male, no college), (male, college). Express each in terms of $\alpha, \beta_1, \beta_2, \beta_3$.

\textbf{7.7} Show that the OLS estimator for WLS with weights $w_i$ solves $\hat{\boldsymbol{\beta}}^{WLS} = (\mathbf{X}^T\mathbf{W}\mathbf{X})^{-1}\mathbf{X}^T\mathbf{W}\mathbf{y}$.

\textbf{7.8} Prove that the saturated model with two binary regressors produces fitted values equal to the sample cell means.

\subsection{Applied Questions}

\textbf{7.9} Using CPS data, estimate the Mincer earnings equation. Report: (a) the return to education; (b) the return to experience at 5, 15, and 25 years; (c) the peak experience level.

\textbf{7.10} Add an interaction between education and ethnicity. Does the return to education differ significantly by ethnicity?

\medskip\hrule\medskip

\section{R Lab 7: Dummies, Interactions, and Non-linearities}

\begin{lstlisting}
library(tidyverse)
library(AER)
data("CPS1988")
CPS1988 <- CPS1988 %>%
  mutate(log_wage = log(wage),
         ethnicity_dummy = as.integer(ethnicity == "afam"))

# 7.1: Mincer earnings equation
mincer <- lm(log_wage ~ education + experience + I(experience^2), data = CPS1988)
summary(mincer)

# Peak experience
beta_exp  <- coef(mincer)["experience"]
beta_exp2 <- coef(mincer)["I(experience^2)"]
peak_exp  <- -beta_exp / (2 * beta_exp2)
cat("Peak experience:", round(peak_exp, 1), "years\n")

# 7.2: Interaction with ethnicity dummy
mincer_int <- lm(log_wage ~ education * ethnicity + experience + I(experience^2),
                  data = CPS1988)
summary(mincer_int)

# 7.3: Plot experience-earnings profile
exp_grid <- seq(0, 45, by = 1)
# Predicted log wages at mean education, varying experience
newdata <- data.frame(
  education = mean(CPS1988$education),
  experience = exp_grid,
  ethnicity = "cauc"
)
newdata$pred <- predict(mincer_int, newdata = newdata)

ggplot(newdata, aes(x = experience, y = pred)) +
  geom_line(color = "steelblue", linewidth = 1.2) +
  labs(
    title = "Experience-Earnings Profile",
    subtitle = "Quadratic specification at mean education",
    x = "Years of Experience",
    y = "Predicted Log Wage"
  ) +
  theme_minimal()

# 7.4: Log-log regression for wage elasticity w.r.t. education
log_log <- lm(log_wage ~ log(education) + experience + I(experience^2), data = CPS1988)
summary(log_log)
cat("Elasticity of wage w.r.t. education:", round(coef(log_log)["log(education)"], 3), "\n")
\end{lstlisting}

\medskip\hrule\medskip

\textit{End of Chapter 7}

\medskip\hrule\medskip

\textbf{References}

Angrist, J. D. (1998). Estimating the labor market impact of voluntary military service using social security data on military applicants. \textit{Econometrica}, 66(2), 249--288.

Mincer, J. (1974). \textit{Schooling, Experience, and Earnings}. NBER.

Solon, G., Haider, S. J., and Wooldridge, J. M. (2015). What are we weighting for? \textit{Journal of Human Resources}, 50(2), 301--316.

Wooldridge, J. M. (2019). \textit{Introductory Econometrics}. Chapter 6.


%% ===========================================================
%%  PART III: IDENTIFICATION STRATEGIES
%% ===========================================================
\part{Identification Strategies}

\chapter{Chapter 8: Instrumental Variables and Two-Stage Least Squares}

\medskip\hrule\medskip

\section{Chapter Overview}

Chapters 2--7 developed regression analysis as a tool for estimating conditional means and partial effects. But we repeatedly noted that OLS identifies causal effects only under strict exogeneity --- an assumption that fails whenever there is omitted variable bias, reverse causality, or measurement error. This chapter introduces \textbf{Instrumental Variables (IV)} --- arguably the most important identification strategy in applied economics --- as a way to estimate causal effects when OLS is endogenous.

The key insight: if we can find a variable $Z_i$ (the "instrument") that is correlated with the endogenous regressor $X_i$ but affects the outcome $Y_i$ only through $X_i$, then we can use the exogenous variation in $Z_i$ to "purge" the endogeneity from $X_i$.

\medskip\hrule\medskip

\section{Learning Objectives}

After completing this chapter, students will be able to:

\begin{enumerate}
\item Explain the \textbf{endogeneity problem} and why OLS fails.
\item Define the two conditions for a valid instrument: \textbf{relevance} and \textbf{exclusion}.
\item Derive the \textbf{IV estimator} (Wald estimator) in the bivariate case.
\item Derive and interpret the \textbf{Two-Stage Least Squares (2SLS)} estimator in the general case.
\item Derive the asymptotic distribution of 2SLS.
\item Explain the \textbf{weak instruments} problem and its consequences.
\item Use the \textbf{first-stage $F$-statistic} to test for weak instruments.
\item Apply IV methods to canonical empirical examples (Angrist and Krueger, 1991).
\end{enumerate}

\medskip\hrule\medskip

\section{8.1 The Endogeneity Problem}

\textbf{Definition 8.1 (Endogeneity):} The regressor $X_i$ is \textbf{endogenous} in the regression $Y_i = \alpha + \rho X_i + u_i$ if $\text{Cov}(X_i, u_i) \neq 0$.

When this holds, OLS is biased and inconsistent:

\[
\text{plim}(\hat{\rho}^{OLS}) = \rho + \frac{\text{Cov}(X_i, u_i)}{V(X_i)} \neq \rho
\]

The three main sources of endogeneity are:

\begin{enumerate}
\item \textbf{Omitted variables:} $u_i$ contains an omitted variable $A_i$ correlated with $X_i$ (e.g., ability in the wage regression).
\item \textbf{Reverse causality:} $Y_i$ causes $X_i$ as well as vice versa (e.g., income and health).
\item \textbf{Measurement error:} $X_i$ is measured with error correlated with the true $X_i^*$ (attenuation bias).
\end{enumerate}

\medskip\hrule\medskip

\section{8.2 The Instrumental Variables (IV) Estimator}

\subsection{8.2.1 What Is an Instrument?}

An \textbf{instrument} $Z_i$ must satisfy two conditions:

\textbf{Condition 1 --- Relevance (First Stage):} $\text{Cov}(Z_i, X_i) \neq 0$.

The instrument must be correlated with the endogenous regressor.

\textbf{Condition 2 --- Exclusion Restriction (Exogeneity):} $\text{Cov}(Z_i, u_i) = 0$.

The instrument must be uncorrelated with the structural error --- it affects $Y_i$ only through $X_i$.

Graphically: $Z_i \to X_i \to Y_i$ (with $Z_i$ having no direct path to $Y_i$).

\begin{figure}[htbp]
  \centering
  \begin{tikzpicture}[
    mynode/.style={circle, draw=econblue, fill=econblue!10, thick, minimum size=12mm, font=\small\bfseries},
    confnode/.style={circle, draw=gray!60, fill=gray!10, thick, minimum size=10mm, font=\small\bfseries},
    ivnode/.style={circle, draw=green!50!black, fill=green!10, thick, minimum size=12mm, font=\small\bfseries},
    myarrow/.style={-{Stealth[scale=1.2]}, thick, econblue},
    confarrow/.style={-{Stealth[scale=1.2]}, thick, dashed, gray},
    scale=1.2
  ]
    \node[ivnode] (Z) at (0, 0) {$Z_i$};
    \node[mynode] (D) at (3, 0) {$D_i$};
    \node[mynode] (Y) at (6, 0) {$Y_i$};
    \node[confnode] (U) at (4.5, 2) {$U_i$};

    % First stage
    \draw[myarrow, green!50!black] (Z) -- (D)
      node[midway, below, font=\footnotesize] {1st stage};
    % Second stage
    \draw[myarrow] (D) -- (Y)
      node[midway, below, font=\footnotesize] {causal effect};
    % Confounder
    \draw[confarrow] (U) -- (D);
    \draw[confarrow] (U) -- (Y);
    % No direct Z->Y (exclusion restriction)
    \draw[red!70!black, thick, dashed] (Z) to[bend left=30] node[above, font=\footnotesize, red!70!black] {$\times$ excluded} (Y);

    % Labels
    \node[font=\footnotesize, green!40!black, align=center] at (0, -0.9) {Instrument\\(exogenous)};
    \node[font=\footnotesize, gray, align=center] at (4.5, 2.7) {Confounder\\(unobserved)};

    % IV conditions
    \node[draw=econblue!30, fill=econblue!5, rounded corners, font=\footnotesize, align=left] at (8.5, 1)
      {\textbf{IV Conditions:}\\(1) Relevance: $Z_i \not\!\perp D_i$\\(2) Exclusion: $Z_i \perp Y_i | D_i, U_i$\\(3) Independence: $Z_i \perp U_i$};
  \end{tikzpicture}
  \caption{DAG for an instrumental variables (IV) strategy. The instrument $Z_i$ (green) affects the outcome $Y_i$ \emph{only} through the treatment $D_i$ (exclusion restriction, red dashed), is correlated with $D_i$ (relevance), and is independent of the unobserved confounder $U_i$.}
  \label{fig:ch08_iv_dag}
\end{figure}

\subsection{8.2.2 The Wald Estimator (Bivariate Case)}

In the bivariate model $Y_i = \alpha + \rho X_i + u_i$ with a single binary instrument $Z_i \in \{0,1\}$:

\[
\hat{\rho}^{IV} = \frac{E[Y_i \mid Z_i = 1] - E[Y_i \mid Z_i = 0]}{E[X_i \mid Z_i = 1] - E[X_i \mid Z_i = 0]} = \frac{\text{Reduced Form}}{\text{First Stage}}
\]

This is the \textbf{Wald estimator}: the ratio of the effect of the instrument on the outcome to the effect of the instrument on the endogenous regressor.

\textbf{Why this works:} If $Z_i$ is randomly assigned (or as good as random), it affects $Y_i$ only through $X_i$. Dividing the total effect of $Z_i$ on $Y_i$ by the effect of $Z_i$ on $X_i$ "scales" the outcome effect to be per unit of $X_i$ --- yielding the causal effect of $X_i$ on $Y_i$.

\subsection{8.2.3 Consistency of the IV Estimator}

\textbf{Proof:} Using sample moments:

\[
\hat{\rho}^{IV} = \frac{\text{Cov}_s(Z_i, Y_i)}{\text{Cov}_s(Z_i, X_i)}
\]

Substituting $Y_i = \alpha + \rho X_i + u_i$:

\[
\frac{\text{Cov}(Z_i, Y_i)}{\text{Cov}(Z_i, X_i)} = \frac{\text{Cov}(Z_i, \alpha + \rho X_i + u_i)}{\text{Cov}(Z_i, X_i)} = \frac{\rho \text{Cov}(Z_i, X_i) + \text{Cov}(Z_i, u_i)}{\text{Cov}(Z_i, X_i)}
\]

Under the exclusion restriction $\text{Cov}(Z_i, u_i) = 0$: $\hat{\rho}^{IV} \xrightarrow{p} \rho$. $\square$

\medskip\hrule\medskip

\section{8.3 Two-Stage Least Squares (2SLS): The General Case}

With multiple endogenous variables and multiple instruments, we use \textbf{Two-Stage Least Squares (2SLS)}.

\subsection{8.3.1 The General Setup}

The \textbf{structural equation} is:

\[
\mathbf{y} = \mathbf{X}\boldsymbol{\beta} + \mathbf{u}
\]

where $\mathbf{X} = [\mathbf{X}_1, \mathbf{X}_2]$: $\mathbf{X}_1$ contains the exogenous regressors (controls) and $\mathbf{X}_2$ contains the $k_2$ endogenous variables.

The \textbf{instrument set} is $\mathbf{Z} = [\mathbf{X}_1, \mathbf{Z}_2]$: the same exogenous controls $\mathbf{X}_1$ plus $\ell \geq k_2$ excluded instruments $\mathbf{Z}_2$.

For \textbf{identification}: $\ell \geq k_2$ (order condition --- at least as many instruments as endogenous variables).

\subsection{8.3.2 The Two Stages}

\textbf{Stage 1 (First Stage):} Regress each endogenous variable on all instruments:

\[
\mathbf{X}_2 = \mathbf{Z}\boldsymbol{\Pi} + \mathbf{V}
\]

Obtain fitted values: $\hat{\mathbf{X}}_2 = \mathbf{Z}(\mathbf{Z}^T\mathbf{Z})^{-1}\mathbf{Z}^T\mathbf{X}_2 = \mathbf{P_Z}\mathbf{X}_2$

This decomposes $\mathbf{X}_2$ into a part predicted by the instruments ($\hat{\mathbf{X}}_2$, exogenous variation) and a residual ($\hat{\mathbf{V}}$, endogenous variation).

\textbf{Stage 2 (Second Stage):} Regress $\mathbf{y}$ on $\hat{\mathbf{X}} = [\mathbf{X}_1, \hat{\mathbf{X}}_2]$:

\[
\hat{\boldsymbol{\beta}}^{2SLS} = (\hat{\mathbf{X}}^T\hat{\mathbf{X}})^{-1}\hat{\mathbf{X}}^T\mathbf{y}
\]

\subsection{8.3.3 The 2SLS Estimator in Matrix Form}

The 2SLS estimator can be written more compactly:

\[
\hat{\boldsymbol{\beta}}^{2SLS} = (\mathbf{X}^T\mathbf{P_Z}\mathbf{X})^{-1}\mathbf{X}^T\mathbf{P_Z}\mathbf{y}
\]

\textbf{Key insight:} 2SLS projects both $\mathbf{X}$ and $\mathbf{y}$ onto the column space of $\mathbf{Z}$ before running the regression. This replaces the endogenous variation in $\mathbf{X}$ with the exogenous variation from the instruments.

\textbf{Special case:} When $\mathbf{Z} = \mathbf{X}$ (all regressors are exogenous), $\mathbf{P_Z} = \mathbf{P_X}$ and 2SLS = OLS.

\begin{figure}[htbp]
  \centering
  \begin{tikzpicture}[
    box/.style={draw=econblue, fill=econblue!8, rounded corners, thick, minimum width=4cm, minimum height=1.1cm, align=center, font=\small},
    arr/.style={-{Stealth}, thick, econblue},
    scale=0.9
  ]
    % Stage 1
    \node[box, minimum width=5.5cm] (s1reg) at (0, 0) {Regress $D_i$ on $Z_i$ (and controls)\\\small First Stage};
    \node[box, minimum width=5.5cm, fill=green!10, draw=green!50!black] (s1out) at (0, -1.8) {$\hat{D}_i = \hat{\pi}_0 + \hat{\pi}_1 Z_i$\\\small Predicted (exogenous) treatment};
    \draw[arr, green!50!black] (s1reg) -- (s1out);

    % Stage 2
    \node[box, minimum width=5.5cm] (s2reg) at (7, 0) {Regress $Y_i$ on $\hat{D}_i$ (and controls)\\\small Second Stage};
    \node[box, minimum width=5.5cm, fill=econblue!20] (s2out) at (7, -1.8) {$\hat{\beta}^{2SLS}$: IV estimate of\\causal effect of $D$ on $Y$};
    \draw[arr] (s2reg) -- (s2out);

    % Arrow from stage 1 to stage 2
    \draw[arr, thick] (s1out) -- (s2reg)
      node[midway, above, font=\footnotesize] {use $\hat{D}_i$ in place of $D_i$};

    % Notes
    \node[font=\footnotesize, gray, align=center] at (3.5, -3.2)
      {Key: $\hat{D}_i$ captures only the exogenous variation in $D_i$ induced by $Z_i$,\\purging the endogeneity caused by $U_i$.};
  \end{tikzpicture}
  \caption{Two-Stage Least Squares (2SLS). In the first stage, we project $D_i$ onto the instrument $Z_i$ to obtain the ``clean'' predicted values $\hat{D}_i$. In the second stage, we regress $Y_i$ on $\hat{D}_i$, using only the exogenous variation in treatment to identify the causal effect.}
  \label{fig:ch08_2sls}
\end{figure}

\medskip\hrule\medskip

\section{8.4 Asymptotic Properties of 2SLS}

\textbf{Theorem 8.1 (Consistency of 2SLS):} Under the conditions:
\begin{enumerate}
\item Linear model $\mathbf{y} = \mathbf{X}\boldsymbol{\beta} + \mathbf{u}$
\item Instrument relevance: $E[\mathbf{z}_i \mathbf{x}_i^T]$ has rank $k$
\item Exclusion restriction: $E[\mathbf{z}_i u_i] = \mathbf{0}$
\end{enumerate}

Then $\hat{\boldsymbol{\beta}}^{2SLS} \xrightarrow{p} \boldsymbol{\beta}$.

\textbf{Theorem 8.2 (Asymptotic Distribution of 2SLS):}

\[
\sqrt{n}(\hat{\boldsymbol{\beta}}^{2SLS} - \boldsymbol{\beta}) \xrightarrow{d} N\!\left(\mathbf{0}, \mathbf{V}^{2SLS}\right)
\]

where under homoskedasticity:

\[
\mathbf{V}^{2SLS} = \sigma^2 \left(E[\mathbf{x}_i \mathbf{z}_i^T]\left(E[\mathbf{z}_i\mathbf{z}_i^T]\right)^{-1}E[\mathbf{z}_i\mathbf{x}_i^T]\right)^{-1}
\]

In practice, use the 2SLS residuals $\hat{\mathbf{u}}^{2SLS} = \mathbf{y} - \mathbf{X}\hat{\boldsymbol{\beta}}^{2SLS}$ (note: NOT the residuals from Stage 2) to estimate $\sigma^2$ and use robust/cluster SEs.

\textbf{Important warning:} Never use software's default Stage-2 SEs --- they use Stage-2 residuals and are wrong. Always use the correct 2SLS standard errors (or use dedicated IV commands: \texttt{ivreg()} in R, \texttt{ivregress} in Stata).

\medskip\hrule\medskip

\section{8.5 Instrument Validity: Relevance and Exclusion}

\subsection{8.5.1 Testing Relevance: The First-Stage $F$-Statistic}

The relevance assumption can be \textbf{tested}: run the first stage and report the $F$-statistic for the joint significance of the excluded instruments.

\textbf{Rule of thumb (Staiger and Stock, 1997):} First-stage $F > 10$ is often cited as a threshold for "strong" instruments.

\textbf{The weak instruments problem (preview):} When instruments are only weakly correlated with the endogenous variable ($F$ close to 0), 2SLS has poor finite-sample properties:
\begin{itemize}
\item The 2SLS estimator is biased toward OLS.
\item Standard errors are misleading; $t$-tests are severely size-distorted.
\end{itemize}

\subsection{8.5.2 The Exclusion Restriction Cannot Be Tested}

The exclusion restriction $E[\mathbf{z}_i u_i] = \mathbf{0}$ is \textbf{fundamentally untestable} (in the sense that $u_i$ is unobserved). Researchers must argue for it on substantive grounds.

\textbf{Testing over-identifying restrictions (when $\ell > k_2$):} When we have more instruments than needed, we can test whether different instruments give similar estimates --- the \textbf{Sargan-Hansen test}. Rejection suggests at least one instrument is invalid.

\textbf{However:} The Sargan-Hansen test only checks consistency across instruments, not absolute validity. If all instruments are invalid in the same way, it will not detect the problem.

\medskip\hrule\medskip

\section{8.6 Weak Instruments: Diagnosis and Remedies}

\subsection{8.6.1 The Weak Instruments Problem}

\textbf{Definition 8.2 (Weak Instruments):} Instruments are weak if the first-stage relationship is weak --- the instruments explain little variation in the endogenous variable beyond controls.

With weak instruments, even modest violations of the exclusion restriction cause large bias. The bias of 2SLS relative to OLS is approximately:

\[
\text{Bias}(2SLS) \approx \frac{1}{F}\text{Bias}(OLS)
\]

where $F$ is the first-stage $F$-statistic. With $F = 5$, 2SLS is biased about 20\% of the way toward OLS.

\subsection{8.6.2 Detection}

\begin{itemize}
\item \textbf{First-stage $F$-statistic:} Always report and examine. $F < 10$ is cause for concern.
\item \textbf{Partial $R^2$ of instruments:} The fraction of variation in the endogenous variable explained by the excluded instruments alone.
\item \textbf{Anderson-Rubin confidence sets:} Inference that is robust to weak instruments.
\end{itemize}

\subsection{8.6.3 Remedies}

\begin{itemize}
\item Find stronger instruments.
\item Use \textbf{Limited Information Maximum Likelihood (LIML)}, which is median-unbiased under many-instrument asymptotics.
\item Use \textbf{Anderson-Rubin tests} for inference robust to weak instruments.
\end{itemize}

\medskip\hrule\medskip

\section{8.7 Empirical Illustration: Angrist and Krueger (1991)}

\subsection{8.7.1 The Research Question}

Angrist and Krueger (1991) study the return to education using compulsory schooling laws as an instrument. Their key insight: US students born early in the year (Q1) start school at an older age and can legally drop out having completed less education than students born later in the year (Q4). Therefore, \textbf{quarter of birth} creates quasi-random variation in education levels.

\textbf{Instrument:} Quarter of birth (Q1 vs. Q4).

\textbf{Endogenous variable:} Years of schooling.

\textbf{Outcome:} Log weekly earnings.

\subsection{8.7.2 Instrument Validity}

\textbf{Relevance:} Q1 births have on average 0.1--0.2 fewer years of schooling than Q4 births. This is the "season of birth" effect documented in the data. The first-stage relationship, while small, is statistically significant.

\textbf{Exclusion restriction:} Quarter of birth should affect earnings \textit{only} through its effect on schooling. (This is plausible if Q1 vs. Q4 birth is essentially random and birth quarter has no direct effect on labor market outcomes.)

\subsection{8.7.3 Results}

\textbf{OLS estimate:} Return to schooling $\approx$ 7--8\% per year.

\textbf{IV/2SLS estimate (using Q1 birth as instrument):} Return to schooling $\approx$ 8--10\% per year.

The OLS and IV estimates are remarkably similar. This suggests (though doesn't prove) that ability bias in OLS may be smaller than often assumed, or that the instrument is strong enough to identify a similar population parameter.

\begin{keybox}{Box 8.1 --- The Angrist-Krueger (1991) Paper}

Angrist and Krueger use quarter of birth interacted with year of birth and state of birth as instruments --- generating many instruments. Card (1995) later pointed out that many weak instruments can lead to biased 2SLS estimates toward OLS. This prompted important work on weak instruments and LIML. The original paper remains a canonical example of creative instrument design.
\end{keybox}

\subsection{8.7.4 Reduced Form and First Stage}

A useful way to present IV results is to show:

\begin{enumerate}
\item \textbf{First Stage:} Quarter of birth $\to$ Education (must be significant).
\item \textbf{Reduced Form:} Quarter of birth $\to$ Log wages (directly).
\item \textbf{2SLS:} The ratio, rescaling reduced form to "per year of education."
\end{enumerate}

If the reduced form is insignificant, the instrument has no explanatory power --- the IV estimate has no credibility.

\medskip\hrule\medskip

\section{Chapter Summary}

\begin{itemize}
\item \textbf{Endogeneity} ($\text{Cov}(X_i, u_i) \neq 0$) makes OLS biased and inconsistent.
\item An instrument $Z_i$ must satisfy \textbf{relevance} ($\text{Cov}(Z, X) \neq 0$) and \textbf{exclusion} ($\text{Cov}(Z, u) = 0$).
\item The \textbf{Wald estimator} is the ratio of the reduced form to the first stage.
\item \textbf{2SLS} generalizes IV to multiple endogenous variables and instruments by projecting regressors onto the instrument space.
\item \textbf{Weak instruments} cause large bias toward OLS; always report the first-stage $F$-statistic.
\item The exclusion restriction is untestable; defend it on substantive grounds.
\item Angrist and Krueger (1991) is a canonical application using quarter of birth as an instrument for education.
\end{itemize}

\medskip\hrule\medskip

\section{Key Terms}

\textbf{Endogeneity} --- $\text{Cov}(X_i, u_i) \neq 0$; OLS is biased and inconsistent.

\textbf{Instrument} --- variable $Z_i$ correlated with $X_i$ but uncorrelated with $u_i$.

\textbf{Relevance} --- $\text{Cov}(Z_i, X_i) \neq 0$.

\textbf{Exclusion restriction} --- $\text{Cov}(Z_i, u_i) = 0$.

\textbf{Wald estimator} --- $\hat{\rho}^{IV} = \frac{\text{Cov}(Z,Y)}{\text{Cov}(Z,X)}$.

\textbf{Two-Stage Least Squares (2SLS)} --- IV estimator for the general case with multiple instruments and endogenous variables.

\textbf{First stage} --- regression of endogenous variable(s) on all instruments.

\textbf{Reduced form} --- regression of outcome on instruments directly.

\textbf{Weak instruments} --- instruments with small first-stage $F$-statistic; cause bias and size distortion.

\textbf{Sargan-Hansen test} --- test of over-identifying restrictions; checks consistency across instruments.

\medskip\hrule\medskip

\section{Exercises}

\subsection{Conceptual Questions}

\textbf{8.1} Explain why OLS is biased when $\text{Cov}(X_i, u_i) \neq 0$. Give three economic examples of endogeneity.

\textbf{8.2} State the two conditions for a valid instrument. Can we test both in practice? Explain.

\textbf{8.3} Describe the two stages of 2SLS in plain language. What is "purged" in the first stage?

\textbf{8.4} Why is the first-stage $F$-statistic important? What is the consequence of ignoring weak instruments?

\textbf{8.5} "The exclusion restriction is satisfied if the instrument is randomly assigned." Evaluate this claim.

\subsection{Analytical Questions}

\textbf{8.6} In the bivariate model $Y_i = \alpha + \rho X_i + u_i$, with a single instrument $Z_i$:
(a) Derive the Wald estimator.
(b) Show it is consistent when $\text{Cov}(Z_i, u_i) = 0$ and $\text{Cov}(Z_i, X_i) \neq 0$.
(c) What is the probability limit when $\text{Cov}(Z_i, u_i) \neq 0$?

\textbf{8.7} Show that the 2SLS estimator equals $\hat{\boldsymbol{\beta}}^{2SLS} = (\mathbf{X}^T\mathbf{P_Z}\mathbf{X})^{-1}\mathbf{X}^T\mathbf{P_Z}\mathbf{y}$.

\textbf{8.8} In the Angrist-Krueger example with $Z_i$ = 1(Q1 birth) and $X_i$ = years of schooling:
\begin{itemize}
\item $E[Y_i \mid Z_i = 1] = 5.80$, $E[Y_i \mid Z_i = 0] = 5.90$
\item $E[X_i \mid Z_i = 1] = 12.68$, $E[X_i \mid Z_i = 0] = 12.79$
\end{itemize}

Compute the Wald estimate of the return to schooling.

\subsection{Applied Questions}

\textbf{8.9} Using the \texttt{card.dta} data (Card, 1995): college proximity as an instrument for education. In R:
(a) Run OLS of log wages on education and controls.
(b) Report the first-stage $F$-statistic.
(c) Run 2SLS using college proximity as the instrument.
(d) Compare OLS and 2SLS estimates. Does the direction of change make sense?

\medskip\hrule\medskip

\section{R Lab 8: IV and 2SLS}

\subsection{Setup}

\begin{lstlisting}
library(tidyverse)
library(AER)    # ivreg() for 2SLS
library(lmtest)
library(sandwich)

# Card (1995) data: effect of education on wages using college proximity as IV
# Available in the ivreg package or can be read from the PS3 data
# install.packages("ivreg")
library(ivreg)
data("SchoolingReturns", package = "ivreg")
df <- SchoolingReturns
\end{lstlisting}

\subsection{Exercise 8.1: OLS Baseline}

\begin{lstlisting}
# OLS regression: log wage on education and controls
ols_model <- lm(wage ~ education + experience + I(experience^2) + ethnicity + 
                  smsa + south, data = df)
summary(ols_model)
cat("OLS coefficient on education:", round(coef(ols_model)["education"], 4), "\n")
\end{lstlisting}

\subsection{Exercise 8.2: First Stage}

\begin{lstlisting}
# First stage: education on nearcollege (proximity to 4-year college) + controls
first_stage <- lm(education ~ nearcollege + experience + I(experience^2) + 
                    ethnicity + smsa + south, data = df)
summary(first_stage)
cat("First-stage coefficient on nearcollege:", round(coef(first_stage)["nearcollege"], 4), "\n")
cat("First-stage F-statistic:", round(summary(first_stage)$fstatistic[1], 2), "\n")

# Partial F-test for just the instrument
ols_no_iv <- lm(education ~ experience + I(experience^2) + ethnicity + smsa + south, data = df)
F_partial <- anova(ols_no_iv, first_stage)[["F"]][2]
cat("Partial F-statistic for instrument:", round(F_partial, 2), "\n")
\end{lstlisting}

\subsection{Exercise 8.3: Reduced Form}

\begin{lstlisting}
# Reduced form: log wage on nearcollege directly
reduced_form <- lm(wage ~ nearcollege + experience + I(experience^2) + 
                     ethnicity + smsa + south, data = df)
summary(reduced_form)
cat("Reduced form coefficient:", round(coef(reduced_form)["nearcollege"], 4), "\n")
\end{lstlisting}

\subsection{Exercise 8.4: 2SLS Estimation}

\begin{lstlisting}
# 2SLS using ivreg()
iv_model <- ivreg(wage ~ education + experience + I(experience^2) + ethnicity + smsa + south |
                  nearcollege + experience + I(experience^2) + ethnicity + smsa + south,
                  data = df)
summary(iv_model, diagnostics = TRUE)

cat("2SLS coefficient on education:", round(coef(iv_model)["education"], 4), "\n")

# Wald estimator (bivariate comparison)
wald_est <- coef(reduced_form)["nearcollege"] / coef(first_stage)["nearcollege"]
cat("Wald estimate:", round(wald_est, 4), "\n")
cat("2SLS estimate:", round(coef(iv_model)["education"], 4), "\n")
\end{lstlisting}

\subsection{Exercise 8.5: Visualizing First Stage}

\begin{lstlisting}
ggplot(df, aes(x = factor(nearcollege), y = education)) +
  geom_boxplot(fill = "steelblue", alpha = 0.7) +
  labs(
    title = "First Stage: College Proximity and Years of Education",
    subtitle = "Living near a college is associated with more schooling",
    x = "Near College (0 = No, 1 = Yes)",
    y = "Years of Education"
  ) +
  theme_minimal()
\end{lstlisting}

\medskip\hrule\medskip

\textit{End of Chapter 8}

\medskip\hrule\medskip

\textbf{References}

Angrist, J. D., and Krueger, A. B. (1991). Does compulsory school attendance affect schooling and earnings? \textit{Quarterly Journal of Economics}, 106(4), 979--1014.

Card, D. (1995). Using geographic variation in college proximity to estimate the return to schooling. In L. Christofides et al. (eds.), \textit{Aspects of Labour Market Behaviour}. University of Toronto Press.

Staiger, D., and Stock, J. H. (1997). Instrumental variables regression with weak instruments. \textit{Econometrica}, 65(3), 557--586.

Davidson, R., and MacKinnon, J. G. (2004). \textit{Econometric Theory and Methods}. Chapter 8.

Angrist, J. D., and Pischke, J. (2009). \textit{Mostly Harmless Econometrics}. Chapter 4.

\chapter{Chapter 9: Heterogeneous Treatment Effects and the Local Average Treatment Effect}

\medskip\hrule\medskip

\section{Chapter Overview}

Chapter 8 presented IV as a solution to endogeneity when the exclusion restriction holds. But what does IV actually estimate when treatment effects are heterogeneous --- when the effect of the treatment varies across individuals? This chapter answers that question using the \textbf{potential outcomes framework} developed in Chapter 1.

The central result --- the \textbf{LATE Theorem} of Imbens and Angrist (1994) --- shows that IV with a binary instrument estimates the \textbf{Local Average Treatment Effect}: the average causal effect for a specific subpopulation called \textbf{compliers} --- individuals whose treatment status is affected by the instrument. This has profound implications for interpreting IV estimates and designing empirical research.

\medskip\hrule\medskip

\section{Learning Objectives}

After completing this chapter, students will be able to:

\begin{enumerate}
\item Explain the limitations of the \textbf{constant-effects model} and why heterogeneous effects matter.
\item Define \textbf{potential outcomes} for both the outcome and the treatment.
\item Define the four \textbf{compliance types} (always-takers, never-takers, compliers, defiers).
\item State the four \textbf{LATE assumptions} (independence, exclusion, first stage, monotonicity).
\item State and prove the \textbf{LATE Theorem} (Imbens and Angrist, 1994).
\item Interpret IV estimates as LATE, and understand when LATE differs from ATE.
\item Explain the relationship between the \textbf{Wald estimator} and LATE.
\item Apply the LATE framework to canonical examples (Vietnam draft lottery, compulsory schooling).
\end{enumerate}

\medskip\hrule\medskip

\section{9.1 The Constant-Effects Model and Its Limitations}

In Chapter 8 (and Chapter 1), we wrote the structural equation as:

\[
Y_i = \alpha + \rho D_i + u_i
\]

where $\rho = Y_{1i} - Y_{0i}$ is the same for everyone --- the \textbf{constant-effects model}.

This assumption is convenient but unrealistic. In labour economics, the return to education varies by individual ability, occupation, and timing. In health economics, the effect of a treatment varies by disease severity. In program evaluation, training program effects vary by background and motivation.

\textbf{The question this chapter answers:} When effects are heterogeneous and we use an IV estimator, what causal parameter does it recover?

\medskip\hrule\medskip

\section{9.2 The Potential Outcomes Framework Revisited}

We now have potential outcomes for \textbf{both} the treatment and the outcome.

\textbf{Potential treatments:} For each individual $i$ and each possible value of the instrument $Z_i \in \{0,1\}$:

\[
D_{0i} = \text{treatment status individual } i \text{ would take if } Z_i = 0
\]
\[
D_{1i} = \text{treatment status individual } i \text{ would take if } Z_i = 1
\]

\textbf{Potential outcomes:} Given treatment $d \in \{0,1\}$:

\[
Y_{di} = \text{outcome individual } i \text{ would experience if } D_i = d
\]

\textbf{Observed outcomes} (via the switching equations):

\[
D_i = D_{0i} + (D_{1i} - D_{0i})Z_i
\]
\[
Y_i = Y_{0i} + (Y_{1i} - Y_{0i})D_i
\]

\medskip\hrule\medskip

\section{9.3 SUTVA}

A key background assumption is \textbf{Stable Unit Treatment Value Assumption (SUTVA)}:

\textbf{Assumption 9.1 (SUTVA):} 
\begin{enumerate}
\item \textbf{No interference:} One unit's treatment does not affect another's outcomes.
\item \textbf{No hidden versions of treatment:} There is only one version of each treatment level.
\end{enumerate}

SUTVA allows us to write individual potential outcomes without conditioning on others' treatments. It fails in settings with spillovers (e.g., vaccination programs that create herd immunity) or when there are multiple treatment variants (e.g., different types of job training).

\medskip\hrule\medskip

\section{9.4 Compliance Types: Always-Takers, Never-Takers, Compliers, Defiers}

Given the instrument $Z_i$ and potential treatments $D_{0i}, D_{1i}$, we can classify all individuals into four compliance types:

\begin{center}
\begin{tabular}{@{} l l l l @{}}
\toprule
\textbf{Type} & \textbf{$D_{0i}$} & \textbf{$D_{1i}$} & \textbf{Behavior} \\
\midrule
\textbf{Always-takers} & 1 & 1 & Take treatment regardless of $Z$ \\
\textbf{Never-takers} & 0 & 0 & Never take treatment regardless of $Z$ \\
\textbf{Compliers} & 0 & 1 & Take treatment when $Z=1$, not when $Z=0$ \\
\textbf{Defiers} & 1 & 0 & Take treatment when $Z=0$, not when $Z=1$ \\
\bottomrule
\end{tabular}
\end{center}

\begin{figure}[htbp]
  \centering
  \begin{tikzpicture}[
    cell/.style={draw=econblue!30, minimum width=3.5cm, minimum height=1.3cm, align=center, font=\small},
    header/.style={draw=econblue, fill=econblue!20, minimum width=3.5cm, minimum height=0.9cm, align=center, font=\small\bfseries},
    scale=1.0
  ]
    % Top-left corner label
    \node[font=\footnotesize, gray] at (-0.5, 0.5) {};

    % Column headers
    \node[header] at (2.0, 0) {$D_i(1) = 1$\\(take if offered)};
    \node[header] at (5.5, 0) {$D_i(1) = 0$\\(refuse if offered)};

    % Row headers
    \node[header, minimum width=2.5cm, minimum height=1.3cm] at (-0.8, -1.5) {$D_i(0)=0$\\(don't take\\if not offered)};
    \node[header, minimum width=2.5cm, minimum height=1.3cm] at (-0.8, -3.2) {$D_i(0)=1$\\(take even\\if not offered)};

    % Inner cells
    \node[cell, fill=green!15] at (2.0, -1.5) {\textbf{Complier}\\follows instrument};
    \node[cell, fill=red!10] at (5.5, -1.5) {\textbf{Never-taker}\\never treated};
    \node[cell, fill=orange!12] at (2.0, -3.2) {\textbf{Always-taker}\\always treated};
    \node[cell, fill=gray!15] at (5.5, -3.2) {\textbf{Defier}\\(monotonicity rules out)};

    % LATE annotation
    \node[draw=green!50!black, fill=green!8, rounded corners, font=\footnotesize, align=center] at (3.75, -4.5)
      {IV/2SLS identifies LATE = causal effect for \textbf{Compliers} only};
  \end{tikzpicture}
  \caption{The four compliance types, defined by potential treatment status under each instrument value. IV estimates the Local Average Treatment Effect (LATE): the average causal effect for \emph{compliers} --- units whose treatment status is switched by the instrument. Under monotonicity (no defiers), LATE is point identified.}
  \label{fig:ch09_compliance_types}
\end{figure}

\textbf{Key insight:}
\begin{itemize}
\item We can \textbf{never directly observe} which type an individual is, because we observe only one of $D_{0i}$ or $D_{1i}$.
\item Always-takers and never-takers are "uninformative" about the causal effect --- their treatment status doesn't change with the instrument.
\item The instrument only affects the treatment of \textbf{compliers} (and defiers).
\end{itemize}

\textbf{Example 9.1 (Vietnam Draft Lottery):} Consider a lottery that drafts men into military service. 
\begin{itemize}
\item Always-takers: men who volunteer regardless of whether they're drafted.
\item Never-takers: men who avoid service regardless (flee the country, obtain exemptions).
\item Compliers: men who serve if drafted and don't serve if not drafted.
\item Defiers: men who serve only when NOT drafted (unlikely to exist).
\end{itemize}

\medskip\hrule\medskip

\section{9.5 The Four LATE Assumptions}

\textbf{Assumption 9.2 (Independence / Random Assignment):}

\[
(Y_{0i}, Y_{1i}, D_{0i}, D_{1i}) \perp\!\!\!\perp Z_i
\]

The instrument $Z_i$ is independent of all potential outcomes and potential treatments. This is satisfied by random assignment.

\textbf{Assumption 9.3 (Exclusion Restriction):}

\[
Y_{di} = Y_{di}(Z_i) \text{ for } d \in \{0,1\}
\]

Or more precisely: the instrument affects outcomes only through the treatment channel. $Z_i$ has no direct effect on $Y_i$ given $D_i$.

\textbf{Assumption 9.4 (First Stage / Relevance):}

\[
E[D_{1i} - D_{0i}] \neq 0
\]

The instrument has a non-zero (on average) effect on the treatment. 

\textbf{Assumption 9.5 (Monotonicity):}

\[
D_{1i} \geq D_{0i} \text{ for all } i \quad \text{(or } D_{1i} \leq D_{0i} \text{ for all } i)
\]

The instrument shifts treatment in the same direction for everyone. This rules out defiers: if $Z = 1$ increases treatment for some, it does not decrease treatment for others.

The combination of the first-stage and monotonicity assumptions implies that the complier fraction is positive: $P(D_{1i} > D_{0i}) > 0$.

\medskip\hrule\medskip

\section{9.6 The LATE Theorem}

\textbf{Theorem 9.1 (Imbens and Angrist, 1994):} Under Assumptions 9.2--9.5, the IV/Wald estimator identifies the \textbf{Local Average Treatment Effect}:

\[
\text{LATE} \equiv E[Y_{1i} - Y_{0i} \mid D_{1i} > D_{0i}] = \frac{E[Y_i \mid Z_i = 1] - E[Y_i \mid Z_i = 0]}{E[D_i \mid Z_i = 1] - E[D_i \mid Z_i = 0]}
\]

The LATE is the average treatment effect for the subpopulation of \textbf{compliers} --- those individuals for whom $D_{1i} > D_{0i}$ (who are induced into treatment by the instrument).

\textbf{Proof:}

\textbf{Denominator (First Stage):}

\[
E[D_i \mid Z_i = 1] - E[D_i \mid Z_i = 0]
\]

By independence (Assumption 9.2): $E[D_i \mid Z_i = z] = E[D_{zi}]$.

\[
= E[D_{1i}] - E[D_{0i}] = E[D_{1i} - D_{0i}]
\]

By monotonicity: $D_{1i} - D_{0i} \in \{0, 1\}$ (no defiers means the difference is either 0 or 1). So:

\[
= P(D_{1i} > D_{0i}) = P(\text{Complier})
\]

\textbf{Numerator (Reduced Form):}

\[
E[Y_i \mid Z_i = 1] - E[Y_i \mid Z_i = 0]
\]

By independence and the switching equation:

\[
= E[Y_{D_{1i}, i}] - E[Y_{D_{0i}, i}]
\]

\[
= E[Y_{1i} D_{1i} + Y_{0i}(1-D_{1i})] - E[Y_{1i} D_{0i} + Y_{0i}(1-D_{0i})]
\]

\[
= E[(Y_{1i} - Y_{0i})(D_{1i} - D_{0i})]
\]

By monotonicity, $D_{1i} - D_{0i} \geq 0$, so:

\[
= E[(Y_{1i} - Y_{0i})(D_{1i} - D_{0i})]
\]

\[
= E[Y_{1i} - Y_{0i} \mid D_{1i} > D_{0i}] \cdot P(D_{1i} > D_{0i})
\]

(the always-takers and never-takers have $D_{1i} = D_{0i}$, so they contribute zero to the sum).

\textbf{Combining:}

\[
\frac{\text{Numerator}}{\text{Denominator}} = \frac{E[Y_{1i} - Y_{0i} \mid D_{1i} > D_{0i}] \cdot P(\text{Complier})}{P(\text{Complier})} = E[Y_{1i} - Y_{0i} \mid D_{1i} > D_{0i}] = \text{LATE} \quad \square
\]

\medskip\hrule\medskip

\section{9.7 Interpreting LATE}

\subsection{9.7.1 Who Are Compliers?}

Compliers are the people "on the margin" between treatment and control who are moved by the instrument. Their identity depends on what the instrument is:

\begin{itemize}
\item \textbf{Vietnam draft lottery (Angrist, 1990):} Compliers are men who served in the military because they were drafted, but would not have served voluntarily.
\item \textbf{Compulsory schooling (Angrist-Krueger, 1991):} Compliers are individuals who stayed in school because they weren't old enough to drop out by the end of the term, and would have dropped out otherwise.
\item \textbf{College proximity (Card, 1995):} Compliers are individuals who attended college because a college was nearby, and would not have attended otherwise.
\end{itemize}

\subsection{9.7.2 LATE vs. ATE vs. ATT}

In general, \textbf{LATE $\neq$ ATE $\neq$ ATT} when:

\begin{enumerate}
\item Treatment effects are heterogeneous (vary across individuals).
\item The complier group differs systematically from the full population.
\end{enumerate}

\begin{center}
\begin{tabular}{@{} l l l @{}}
\toprule
\textbf{Parameter} & \textbf{Population} & \textbf{Method of identification} \\
\midrule
ATE & Full population & Random assignment to all \\
ATT & Treated units & Selection on observables (CIA) \\
LATE & Compliers & Instrumental variable \\
\bottomrule
\end{tabular}
\end{center}

\textbf{When does LATE $\approx$ ATE?}
\begin{itemize}
\item When treatment effects are homogeneous.
\item When compliers are representative of the full population.
\end{itemize}

\textbf{When does LATE $\neq$ ATE?}
\begin{itemize}
\item When the instrument selects a non-representative subset. The Vietnam draft lottery compliers are men who didn't volunteer --- possibly those with lower "returns" to military service.
\end{itemize}

\subsection{9.7.3 LATE and External Validity}

The LATE from one instrument is not necessarily the same as the LATE from another instrument --- even when both instruments are valid. Different instruments define different complier populations.

\textbf{Example:} Two valid instruments for education:
\begin{enumerate}
\item Quarter of birth (Angrist-Krueger): compliers are people who wanted to drop out as soon as legally allowed.
\item College proximity (Card): compliers are people who would have attended college if it were nearby.
\end{enumerate}

These two groups likely have very different returns to education. This means the two IV estimates can legitimately differ --- not because one is wrong, but because they estimate different parameters for different populations.

This has profound implications for \textbf{external validity}: an IV estimate tells you the causal effect for the compliers defined by that specific instrument, not necessarily the effect for the average person or for a new policy context.

\medskip\hrule\medskip

\section{9.8 2SLS as LATE with Multiple Instruments}

When there are multiple instruments $Z_1, Z_2, \ldots, Z_\ell$ for a single endogenous variable, 2SLS combines them. Under the LATE framework:

\[
\hat{\rho}^{2SLS} \xrightarrow{p} \sum_{j=1}^\ell \omega_j \cdot \text{LATE}_j
\]

where $\text{LATE}_j$ is the LATE defined by instrument $Z_j$ (the effect for compliers of instrument $j$) and $\omega_j$ are weights that sum to one. The weights are proportional to the variance of the first-stage predictions from each instrument.

This means 2SLS with multiple instruments estimates a \textbf{weighted average of LATEs} --- a complex estimand whose interpretation requires care.

\medskip\hrule\medskip

\section{9.9 Empirical Applications}

\subsection{9.9.1 Application 1: Vietnam Draft Lottery (Angrist, 1990)}

\textbf{Research question:} What is the effect of military service on lifetime earnings?

\textbf{Identification problem:} Veterans and non-veterans differ systematically in many ways (selection into service). OLS estimates of the effect of veteran status are confounded by these differences.

\textbf{Instrument:} Vietnam-era draft lottery. Men born in years 1950--1953 were assigned draft lottery numbers 1--365 based on birthday. Low lottery numbers were called first; men with numbers below a threshold were drafted.

\textbf{Validity:}
\begin{itemize}
\item \textbf{Independence:} Lottery numbers were randomly assigned (approximately).
\item \textbf{Exclusion:} Lottery number affects earnings only through military service (no direct effect of knowing your lottery number on future earnings).
\item \textbf{First stage:} Being assigned a low lottery number significantly increases the probability of military service.
\item \textbf{Monotonicity:} Low numbers increase service (no defiers expected).
\end{itemize}

\textbf{Results:} IV (Wald) estimates suggest military service reduced civilian earnings by approximately 15\% for compliers (white veterans), with effects persisting for at least 15 years. The OLS estimate was biased upward because veterans tend to come from lower-income backgrounds (attenuation by selection).

\textbf{Who are the compliers?} Men who were called up by the draft lottery and served. These are NOT representative of all veterans (who also include volunteers) and NOT representative of all men (since many draft-eligible men found other ways to avoid service).

\subsection{9.9.2 Application 2: Compulsory Schooling (Angrist and Krueger, 1991)}

As discussed in Chapter 8:

\textbf{Compliers:} Men who would have dropped out of school as soon as legally permitted (when they turned 16 or after completing the grade they were enrolled in at age 16), but who in fact had to stay in school longer because they started school at a younger age (and were thus younger than the dropout cutoff when they reached certain grades).

\textbf{LATE interpretation:} The Angrist-Krueger estimates measure the return to education for this specific group of marginal students --- those constrained by compulsory schooling laws. The returns for this group may be higher or lower than for the average student.

\medskip\hrule\medskip

\section{Chapter Summary}

\begin{itemize}
\item With heterogeneous treatment effects, the question "what does IV estimate?" is non-trivial.
\item The \textbf{LATE Theorem} shows that IV identifies the average treatment effect for \textbf{compliers} --- those whose treatment is affected by the instrument.
\item The four LATE assumptions are: independence, exclusion restriction, first stage, and monotonicity.
\item \textbf{Always-takers} and \textbf{never-takers} contribute no information to the IV estimate; their effects are "washed out."
\item \textbf{LATE $\neq$ ATE} when compliers are not representative of the full population.
\item Different instruments define different complier populations and may give different (but all valid) estimates.
\item The LATE framework forces us to be explicit about \textit{which population} our IV estimate describes --- crucial for policy relevance.
\end{itemize}

\medskip\hrule\medskip

\section{Key Terms}

\textbf{SUTVA} --- Stable Unit Treatment Value Assumption; no interference between units, no hidden treatment versions.

\textbf{Potential treatments} --- $D_{0i}$ and $D_{1i}$: treatment under $Z_i = 0$ and $Z_i = 1$.

\textbf{Compliance types} --- always-takers, never-takers, compliers, defiers.

\textbf{Compliers} --- individuals with $D_{1i} > D_{0i}$; treatment is affected by the instrument.

\textbf{Monotonicity} --- no defiers: $D_{1i} \geq D_{0i}$ for all $i$.

\textbf{LATE (Local Average Treatment Effect)} --- $E[Y_{1i} - Y_{0i} \mid D_{1i} > D_{0i}]$; average effect for compliers.

\textbf{ATE} --- $E[Y_{1i} - Y_{0i}]$; average effect for the full population.

\textbf{External validity} --- whether an estimate from one context generalizes to another.

\medskip\hrule\medskip

\section{Exercises}

\subsection{Conceptual Questions}

\textbf{9.1} Explain in your own words why IV estimates LATE rather than ATE.

\textbf{9.2} A researcher uses random assignment of a health insurance subsidy as an instrument for health insurance take-up, to estimate the effect of insurance on health outcomes. Who are the compliers in this setting? Who are the always-takers and never-takers? Does LATE have a meaningful policy interpretation?

\textbf{9.3} Why is monotonicity an important assumption for the LATE theorem? What would happen if defiers existed?

\textbf{9.4} Two researchers use different instruments to estimate the return to education. Researcher A uses college proximity; Researcher B uses compulsory schooling laws. Both produce valid IV estimates. Explain why the two estimates might differ even though both are valid.

\textbf{9.5} Evaluate this statement: "Since IV estimates LATE for compliers, and compliers are a small fraction of the population, IV estimates are not useful for policy."

\subsection{Analytical Questions}

\textbf{9.6} Prove the LATE Theorem carefully. Specifically, show that:
(a) The denominator of the Wald estimator equals $P(\text{Complier})$.
(b) The numerator equals $E[Y_{1i} - Y_{0i} \mid \text{Complier}] \times P(\text{Complier})$.

\textbf{9.7} Suppose the population is 20\% always-takers (TE = $\tau_{AT}$), 50\% never-takers (TE = $\tau_{NT}$), and 30\% compliers (TE = $\tau_C$). Express the ATE in terms of $\tau_{AT}$, $\tau_{NT}$, $\tau_C$. Under what conditions does LATE = ATE?

\textbf{9.8} In the draft lottery example, suppose:
\begin{itemize}
\item $P(\text{Complier}) = 0.16$ (16\% of eligible men were induced to serve by the lottery).
\item $E[Y_i \mid Z_i = 1] - E[Y_i \mid Z_i = 0] = -0.024$ (2.4\% reduction in earnings from being in a low-lottery cohort).
\item $E[D_i \mid Z_i = 1] - E[D_i \mid Z_i = 0] = 0.16$.
\end{itemize}

Compute the LATE. Interpret it.

\subsection{Applied Questions}

\textbf{9.9} In R, simulate the LATE framework:
\begin{itemize}
\item Create a population of 10,000 with 4 compliance types.
\item Assign a binary instrument $Z_i$ randomly.
\item Compute the Wald estimator and show it equals the average treatment effect for compliers.
\end{itemize}

\medskip\hrule\medskip

\section{R Lab 9: Simulating the LATE Framework}

\subsection{Exercise 9.1: Simulating Compliance Types}

\begin{lstlisting}
library(tidyverse)
set.seed(2024)
n <- 10000

# Define compliance types (known only to God!)
type <- sample(c("always_taker", "never_taker", "complier", "defier"),
               n, replace = TRUE,
               prob = c(0.20, 0.30, 0.45, 0.05))

# Potential treatment under Z=0 and Z=1
D0 <- ifelse(type %in% c("always_taker", "defier"), 1, 0)
D1 <- ifelse(type %in% c("always_taker", "complier"), 1, 0)

# True heterogeneous treatment effects by type
te <- case_when(
  type == "always_taker" ~ rnorm(n, 5, 2),
  type == "never_taker"  ~ rnorm(n, 3, 2),
  type == "complier"     ~ rnorm(n, 8, 2),  # compliers have higher returns
  type == "defier"       ~ rnorm(n, 1, 2)
)

# Potential outcomes
Y0 <- rnorm(n, 10, 3)         # baseline untreated outcome
Y1 <- Y0 + te                  # treated outcome

# Randomly assign instrument
Z  <- rbinom(n, 1, 0.5)

# Observed treatment and outcome
D  <- ifelse(Z == 1, D1, D0)
Y  <- ifelse(D == 1, Y1, Y0)

df_late <- data.frame(Z, D, Y, type, te, D0, D1)
\end{lstlisting}

\subsection{Exercise 9.2: Computing the LATE}

\begin{lstlisting}
# True LATE (average treatment effect for compliers)
true_LATE <- mean(te[type == "complier"])
cat("True LATE (compliers only):", round(true_LATE, 4), "\n")

# True ATE
true_ATE <- mean(te)
cat("True ATE:", round(true_ATE, 4), "\n")

# Wald estimator (IV)
Y_Z1 <- mean(Y[Z == 1])
Y_Z0 <- mean(Y[Z == 0])
D_Z1 <- mean(D[Z == 1])
D_Z0 <- mean(D[Z == 0])

wald_est <- (Y_Z1 - Y_Z0) / (D_Z1 - D_Z0)
cat("Wald (IV) estimate:", round(wald_est, 4), "\n")
cat("Complier share:", round(D_Z1 - D_Z0, 4), "\n")
\end{lstlisting}

\subsection{Exercise 9.3: IV Regression via ivreg}

\begin{lstlisting}
library(ivreg)
iv_model <- ivreg(Y ~ D | Z, data = df_late)
summary(iv_model)
cat("ivreg IV estimate:", round(coef(iv_model)["D"], 4), "\n")
\end{lstlisting}

\subsection{Exercise 9.4: Distribution of Treatment Effects by Type}

\begin{lstlisting}
df_late %>%
  ggplot(aes(x = te, fill = type)) +
  geom_density(alpha = 0.5) +
  geom_vline(xintercept = true_LATE, linetype = "dashed", color = "red") +
  geom_vline(xintercept = true_ATE, linetype = "dotted", color = "blue") +
  labs(
    title = "Heterogeneous Treatment Effects by Compliance Type",
    subtitle = "Red: LATE (complier mean); Blue: ATE (population mean)",
    x = "Individual Treatment Effect",
    y = "Density",
    fill = "Compliance Type"
  ) +
  theme_minimal()
\end{lstlisting}

\medskip\hrule\medskip

\textit{End of Chapter 9}

\medskip\hrule\medskip

\textbf{References}

Imbens, G. W., and Angrist, J. D. (1994). Identification and estimation of local average treatment effects. \textit{Econometrica}, 62(2), 467--475.

Angrist, J. D. (1990). Lifetime earnings and the Vietnam era draft lottery: Evidence from social security administrative records. \textit{American Economic Review}, 80(3), 313--336.

Angrist, J. D., and Krueger, A. B. (1991). Does compulsory school attendance affect schooling and earnings? \textit{Quarterly Journal of Economics}, 106(4), 979--1014.

Card, D. (1995). Using geographic variation in college proximity to estimate the return to schooling. University of Toronto Press.

Angrist, J. D., and Pischke, J. (2009). \textit{Mostly Harmless Econometrics}. Chapter 4.

Cunningham, S. (2021). \textit{Causal Inference: The Mixtape}. Chapter 7.

\chapter{Chapter 10: Panel Data and Fixed Effects}

\medskip\hrule\medskip

\section{Chapter Overview}

Panel data --- repeated observations on the same units over time --- are among the most powerful data structures in applied economics. They allow researchers to control for \textbf{unobserved individual heterogeneity}: time-invariant characteristics of individuals, firms, or regions that are correlated with both the treatment variable and the outcome, but that we can never directly observe.

This chapter develops the \textbf{fixed effects (FE)} estimator, shows how it is related to OLS via the Frisch-Waugh-Lovell theorem, and discusses its assumptions, limitations, and extensions.

\medskip\hrule\medskip

\section{Learning Objectives}

After completing this chapter, students will be able to:

\begin{enumerate}
\item Explain the \textbf{panel data structure} and define key terminology (balanced/unbalanced panels, $N$, $T$).
\item Explain the \textbf{unobserved heterogeneity problem} and why pooled OLS fails when individual effects are correlated with regressors.
\item Derive the \textbf{within (fixed effects) estimator} via demeaning.
\item Show that FE is equivalent to \textbf{LSDV (Least Squares Dummy Variables)}.
\item State the identifying assumption of FE and explain what it rules out.
\item Contrast the \textbf{first-differences estimator} with the within estimator.
\item Discuss the \textbf{random effects estimator} and the \textbf{Hausman test}.
\item Apply fixed effects to the Card (1995) earnings study.
\end{enumerate}

\medskip\hrule\medskip

\section{10.1 Why Panel Data? The Unobserved Heterogeneity Problem}

\subsection{10.1.1 The Setup}

Suppose we observe individual $i$ in periods $t = 1, 2, \ldots, T$. The outcome $Y_{it}$ depends on an observed regressor $X_{it}$ and an \textbf{unobserved individual effect} $\alpha_i$:

\[
Y_{it} = \alpha_i + \beta X_{it} + u_{it} \tag{10.1}
\]

where:
\begin{itemize}
\item $\alpha_i$ is a \textbf{time-invariant} unobserved characteristic of individual $i$ (e.g., innate ability, firm quality, state-specific resources).
\item $u_{it}$ is an idiosyncratic error term.
\end{itemize}

\subsection{10.1.2 The Pooled OLS Problem}

If we ignore $\alpha_i$ and run OLS on the stacked dataset:

\[
Y_{it} = \alpha + \beta X_{it} + v_{it}
\]

where $v_{it} = \alpha_i - \alpha + u_{it}$, then OLS is biased if $\text{Cov}(X_{it}, \alpha_i) \neq 0$.

\textbf{Example 10.1:} Estimating the return to union membership ($X_{it}$) on wages ($Y_{it}$). $\alpha_i$ might be unobserved worker productivity --- high-productivity workers are both more likely to be unionized and have higher wages independently of union membership. So $\text{Cov}(X_{it}, \alpha_i) > 0$, and pooled OLS overstates the union wage premium.

\begin{figure}[htbp]
  \centering
  \begin{tikzpicture}[scale=0.9,
    unit/.style={draw=econblue, fill=econblue!10, rounded corners, minimum width=1.5cm, minimum height=0.8cm, align=center, font=\footnotesize},
    time/.style={draw=gray!50, fill=gray!10, rounded corners, minimum width=1.5cm, minimum height=0.6cm, align=center, font=\footnotesize\bfseries},
    obs/.style={draw=gray!30, fill=white, minimum width=1.5cm, minimum height=0.8cm, align=center, font=\footnotesize}
  ]
    % Time headers
    \foreach \t/\tname in {1/$t=1$, 2/$t=2$, 3/$t=3$, 4/$t=T$} {
      \node[time] at (\t*1.8 + 0.6, 0) {\tname};
    }
    \node[font=\footnotesize\itshape, gray] at (8.0, 0) {$\cdots$};

    % Unit rows
    \foreach \i/\iname in {1/$i=1$, 2/$i=2$, 3/$i=N$} {
      \node[unit] at (0, -\i*1.2) {\iname};
      \foreach \t in {1,2,3,4} {
        \node[obs] at (\t*1.8 + 0.6, -\i*1.2) {$(Y_{it},\,X_{it})$};
      }
      \node[font=\footnotesize\itshape, gray] at (8.0, -\i*1.2) {$\cdots$};
    }
    \node[font=\footnotesize\itshape, gray] at (0, -3.6) {$\vdots$};
    \foreach \t in {1,2,3,4} {
      \node[font=\footnotesize\itshape, gray] at (\t*1.8+0.6, -3.6) {$\vdots$};
    }

    \node[font=\footnotesize, gray, align=center] at (4.5, -5.2) {Balanced panel: $N$ units observed at $T$ time periods $\Rightarrow$ $NT$ observations total};
  \end{tikzpicture}
  \caption{Structure of a balanced panel dataset. Each cell $(i,t)$ contains observations on unit $i$ at time $t$. With $N$ units and $T$ periods, there are $NT$ total observations. Panel data allow researchers to control for time-invariant unobserved heterogeneity by exploiting within-unit variation over time.}
  \label{fig:ch10_panel_structure}
\end{figure}

\medskip\hrule\medskip

\section{10.2 The Panel Data Model}

\textbf{Definition 10.1 (One-Way Fixed Effects Model):}

\[
Y_{it} = \alpha_i + \beta X_{it} + u_{it}
\]

where $i = 1, \ldots, N$ and $t = 1, \ldots, T$.

\textbf{Key assumption (strict exogeneity within the panel):}

\[
E[u_{it} \mid X_{i1}, X_{i2}, \ldots, X_{iT}, \alpha_i] = 0 \quad \text{for all } t
\]

This says that, conditional on the individual effect $\alpha_i$, the idiosyncratic error is uncorrelated with regressors in all periods. This is stronger than requiring $E[u_{it} \mid X_{it}, \alpha_i] = 0$.

\textbf{No assumption on $\text{Cov}(X_{it}, \alpha_i)$:} Unlike pooled OLS, we allow the individual effect to be correlated with the regressors.

\medskip\hrule\medskip

\section{10.3 The Within (Fixed Effects) Estimator}

\subsection{10.3.1 Demeaning}

The key insight: \textbf{subtract the individual time-mean} to remove $\alpha_i$.

For individual $i$, the time-mean of equation (10.1) is:

\[
\bar{Y}_i = \alpha_i + \beta \bar{X}_i + \bar{u}_i
\]

where $\bar{Y}_i = \frac{1}{T}\sum_{t=1}^T Y_{it}$, etc.

Subtracting from equation (10.1):

\[
\underbrace{(Y_{it} - \bar{Y}_i)}_{\tilde{Y}_{it}} = \beta \underbrace{(X_{it} - \bar{X}_i)}_{\tilde{X}_{it}} + \underbrace{(u_{it} - \bar{u}_i)}_{\tilde{u}_{it}}
\]

The individual effect $\alpha_i$ is eliminated. This is called the \textbf{within transformation} or \textbf{demeaning}.

\subsection{10.3.2 The Within (FE) Estimator}

\textbf{Definition 10.2 (Within Estimator):}

\[
\hat{\beta}^{FE} = \left(\sum_{i=1}^N \sum_{t=1}^T \tilde{X}_{it}^2\right)^{-1} \sum_{i=1}^N \sum_{t=1}^T \tilde{X}_{it} \tilde{Y}_{it}
\]

In matrix notation: $\hat{\boldsymbol{\beta}}^{FE} = (\tilde{\mathbf{X}}^T\tilde{\mathbf{X}})^{-1}\tilde{\mathbf{X}}^T\tilde{\mathbf{Y}}$, where $\tilde{\mathbf{Y}}$ and $\tilde{\mathbf{X}}$ are the demeaned outcome and regressor matrices.

This is \textbf{OLS on the demeaned (within-individual) data}.

\subsection{10.3.3 Connection to the FWL Theorem}

The within estimator is a special case of the FWL Theorem. Let $\mathbf{X}_1$ be the matrix of individual dummies (one per individual) and $\mathbf{X}_2 = X_{it}$. Then:

\[
\hat{\beta}^{FE} = \hat{\beta}_2^{FWL} = (\tilde{\mathbf{X}}^T\tilde{\mathbf{X}})^{-1}\tilde{\mathbf{X}}^T\tilde{\mathbf{Y}}
\]

where $\tilde{\mathbf{X}} = \mathbf{M}_{X_1}\mathbf{X}$ and $\tilde{\mathbf{Y}} = \mathbf{M}_{X_1}\mathbf{Y}$ are the residuals from projecting onto the individual dummy space --- which is exactly demeaning.

\begin{figure}[htbp]
  \centering
  \begin{tikzpicture}
    \begin{axis}[
      width=10cm, height=6.5cm,
      xlabel={$X_{it}$ (e.g., hours worked)},
      ylabel={$Y_{it}$ (e.g., output)},
      xmin=0, xmax=7, ymin=0, ymax=6,
      axis lines=left,
      tick style={draw=none},
      xtick=\empty, ytick=\empty,
      legend style={at={(0.05,0.95)}, anchor=north west, font=\small},
    ]
      % Unit 1 observations (low-output firm, low hours)
      \addplot[only marks, mark=*, mark size=3pt, color=econblue]
        coordinates {(1.5,1.2)(2.0,1.8)(2.5,2.2)};
      % Unit 2 observations (mid-output firm)
      \addplot[only marks, mark=square*, mark size=3pt, color=red!70!black]
        coordinates {(3.0,3.1)(3.5,3.6)(4.0,4.0)};
      % Unit 3 observations (high-output firm)
      \addplot[only marks, mark=triangle*, mark size=3pt, color=green!50!black]
        coordinates {(4.5,4.5)(5.0,4.9)(5.5,5.4)};

      % OLS (between variation -- misleading slope)
      \addplot[thick, gray!60, dashed, domain=0.5:6.5] {0.77*x + 0.05};

      % Within-unit slopes (same slope for all units, shifted)
      \addplot[very thick, econblue, domain=1.0:3.0] {0.5*x + 0.45};
      \addplot[very thick, red!70!black, domain=2.5:4.5] {0.5*x + 1.55};
      \addplot[very thick, green!50!black, domain=4.0:6.0] {0.5*x + 2.25};

      \legend{Unit 1, Unit 2, Unit 3, OLS (pooled), Within-unit slope}
    \end{axis}
  \end{tikzpicture}
  \caption{Within vs.\ between variation in panel data. OLS on pooled data (dashed gray) uses both between-unit and within-unit variation, and is biased when unobserved heterogeneity (unit fixed effects) correlates with $X$. Fixed effects estimation uses only within-unit variation (solid colored lines, common slope) to identify the causal effect.}
  \label{fig:ch10_within_between}
\end{figure}

\subsection{10.3.4 Identifying Assumption of FE}

\textbf{Assumption 10.1 (No Time-Varying Omitted Variables):}

The FE estimator is consistent if, conditional on the individual effect $\alpha_i$, there are no time-varying confounders correlated with $X_{it}$ in $u_{it}$.

\textbf{What FE controls for:} Any individual-specific, time-invariant unobservable.

\textbf{What FE cannot control for:} Time-varying confounders (variables that vary both across individuals and over time, and are correlated with $X_{it}$).

\textbf{Example of what FE handles:} Studying the effect of union membership on wages. FE removes the effect of permanent individual ability (time-invariant). What remains is the association between changes in union status and changes in wages, within the same individual over time.

\textbf{Example of what FE cannot handle:} Studying the effect of a new employment policy on wages, when individuals self-select into regions that adopt the policy at different times based on their expectations of local economic conditions (time-varying, endogenous selection).

\medskip\hrule\medskip

\section{10.4 The LSDV Representation}

\textbf{Proposition 10.1 (LSDV Equivalence):} The within estimator is algebraically identical to the OLS estimator from a regression of $Y_{it}$ on $X_{it}$ and $N$ individual dummy variables (LSDV = Least Squares Dummy Variables):

\[
Y_{it} = \sum_{j=1}^N \alpha_j D_{ij} + \beta X_{it} + u_{it}
\]

where $D_{ij} = 1$ if $j = i$ (individual $i$'s dummy variable).

\textbf{Proof:} By the FWL theorem, the coefficient on $X_{it}$ from the LSDV regression equals the coefficient from regressing $\mathbf{M}_{X_1}\mathbf{Y}$ on $\mathbf{M}_{X_1}\mathbf{X}$, where $\mathbf{X}_1$ is the matrix of individual dummies. Projecting onto individual dummies means subtracting individual means --- which is demeaning. $\square$

\textbf{Practical note:} With large $N$, LSDV is computationally expensive (requires inverting a large matrix). The within estimator is computationally equivalent but much faster.

\medskip\hrule\medskip

\section{10.5 The First-Differences Estimator}

An alternative to demeaning is \textbf{first differencing}: subtract period $t-1$ from period $t$.

\[
\Delta Y_{it} = \beta \Delta X_{it} + \Delta u_{it}
\]

where $\Delta Y_{it} = Y_{it} - Y_{i,t-1}$, etc. Individual effects are again eliminated.

\textbf{FD estimator:} OLS on the first-differenced data.

\textbf{When FD $\equiv$ FE:} When $T = 2$.

\textbf{When FD $\neq$ FE (T > 2):}
\begin{itemize}
\item Under strict exogeneity, both are consistent.
\item If $u_{it}$ is serially uncorrelated, FE is more efficient.
\item If $u_{it}$ follows a random walk ($u_{it} = u_{i,t-1} + \epsilon_{it}$), FD is more efficient.
\end{itemize}

\medskip\hrule\medskip

\section{10.6 Two-Way Fixed Effects}

In many panel settings, we want to control for both individual effects AND time effects --- for example, macro shocks that affect all individuals simultaneously.

\textbf{Two-way FE model:}

\[
Y_{it} = \alpha_i + \gamma_t + \beta X_{it} + u_{it}
\]

where $\gamma_t$ is a \textbf{time fixed effect} (common to all individuals in period $t$).

\textbf{Estimation:} Include individual and time dummies (LSDV) or demean by subtracting both individual means and time means (within transformation in two dimensions).

\textbf{Application:} The two-way FE model is the backbone of the \textbf{Differences-in-Differences} estimator studied in Chapter 11.

\medskip\hrule\medskip

\section{10.7 Random Effects and the Hausman Test}

\subsection{10.7.1 The Random Effects (RE) Model}

\textbf{Definition 10.3 (Random Effects):} The RE model treats $\alpha_i$ as a random variable drawn from a distribution, uncorrelated with $X_{it}$:

\[
\alpha_i \perp X_{it}
\]

Under this assumption, OLS on the pooled data is consistent but inefficient (errors within individuals are correlated). The RE estimator uses a GLS transformation to achieve efficiency.

\subsection{10.7.2 FE vs. RE: The Hausman Test}

\textbf{Hausman (1978) Test:}

\begin{itemize}
\item Under $H_0$: $\text{Cov}(X_{it}, \alpha_i) = 0$ (RE is consistent and efficient; FE is consistent but inefficient).
\item Under $H_1$: $\text{Cov}(X_{it}, \alpha_i) \neq 0$ (RE is inconsistent; FE is consistent).
\end{itemize}

The test statistic is:

\[
H = (\hat{\boldsymbol{\beta}}^{FE} - \hat{\boldsymbol{\beta}}^{RE})^T \left[\text{Var}(\hat{\boldsymbol{\beta}}^{FE}) - \text{Var}(\hat{\boldsymbol{\beta}}^{RE})\right]^{-1} (\hat{\boldsymbol{\beta}}^{FE} - \hat{\boldsymbol{\beta}}^{RE}) \sim \chi^2_k
\]

If $H$ is large (reject $H_0$), use FE. If not, RE may be preferred.

\textbf{Practical advice:} In applied economics, FE is used far more commonly than RE. The assumption that individual effects are uncorrelated with regressors is strong and often implausible. Use RE only when you have a strong theoretical reason to believe the assumption holds.

\medskip\hrule\medskip

\section{10.8 Empirical Illustration: Returns to Education (Card, 1995)}

Card (1995) uses panel data to study whether the returns to education are different from OLS estimates once unobserved individual ability is controlled for.

\textbf{Data:} National Longitudinal Survey (NLS) following the same workers over time.

\textbf{Key finding:} FE estimates of the return to education are \textit{higher} than OLS estimates --- not lower, as ability bias would predict. This suggests that:

\begin{enumerate}
\item Men who complete more education may have had \textit{lower} ability (conditional on measured characteristics) --- a negative $\text{Cov}(X_{it}, \alpha_i)$.
\item Or: OLS ability bias is smaller than typically assumed, while measurement error in education attenuates OLS estimates downward.
\end{enumerate}

This is a sobering finding: the "obvious" direction of ability bias may not be correct once we look within individuals over time.

\medskip\hrule\medskip

\section{Chapter Summary}

\begin{itemize}
\item \textbf{Panel data} allow us to control for time-invariant unobserved heterogeneity.
\item \textbf{Pooled OLS} is biased when individual effects are correlated with regressors.
\item The \textbf{within (FE) estimator} eliminates individual effects by demeaning; it is equivalent to OLS with individual dummies (LSDV).
\item FE's identifying assumption is that there are no time-varying confounders; it controls for everything time-invariant but nothing that varies over time.
\item The \textbf{first-differences estimator} is an alternative to FE; they are equivalent when $T = 2$.
\item \textbf{Two-way FE} controls for both individual and time effects.
\item The \textbf{Hausman test} helps choose between FE and RE, but FE is generally preferred in applied economics.
\item Card (1995) shows FE can produce higher (not lower) estimates of returns to education --- reversing the naive prediction about ability bias.
\end{itemize}

\medskip\hrule\medskip

\section{Key Terms}

\textbf{Panel data} --- data with multiple observations per unit over time.

\textbf{Individual fixed effect} ($\alpha_i$) --- time-invariant unobserved characteristic of unit $i$.

\textbf{Within transformation (demeaning)} --- subtracting individual time-means.

\textbf{Within (FE) estimator} --- OLS on demeaned data; eliminates individual effects.

\textbf{LSDV} --- Least Squares Dummy Variables; OLS with individual dummies; algebraically equivalent to FE.

\textbf{First-differences (FD) estimator} --- OLS on first-differenced data; also eliminates individual effects.

\textbf{Two-way FE} --- controls for both individual and time effects.

\textbf{Random effects (RE)} --- treats individual effects as random draws uncorrelated with regressors; more efficient but requires stronger assumptions.

\textbf{Hausman test} --- tests whether $\text{Cov}(X_{it}, \alpha_i) = 0$; guides FE vs. RE choice.

\medskip\hrule\medskip

\section{Exercises}

\subsection{Conceptual Questions}

\textbf{10.1} Explain why pooled OLS is biased in the presence of unobserved individual heterogeneity. What is the sign of the bias in the returns-to-union-membership example?

\textbf{10.2} What does FE "control for"? What does it NOT control for? Give an example of a confounder that FE cannot address.

\textbf{10.3} Why is FE equivalent to LSDV? Explain using the FWL theorem.

\textbf{10.4} Compare the FE and FD estimators. When are they equivalent? When do they differ?

\textbf{10.5} Explain the Hausman test. What is its null hypothesis? What should you do if the test rejects?

\subsection{Analytical Questions}

\textbf{10.6} Show formally that the within estimator eliminates the individual fixed effect $\alpha_i$.

\textbf{10.7} In the two-period panel ($T = 2$), show that FD and FE give identical estimates.

\textbf{10.8} Consider the model $Y_{it} = \alpha_i + \beta X_{it} + u_{it}$. Suppose $X_{it} = \alpha_i + \epsilon_{it}$ (the regressor depends on the individual effect). Show that pooled OLS is biased and derive the direction of bias.

\textbf{10.9} Show that including individual dummy variables in a regression and applying OLS gives the same coefficient as the within estimator (LSDV equivalence).

\subsection{Applied Questions}

\textbf{10.10} Using a wages panel dataset in R, estimate: (a) pooled OLS; (b) FE (within estimator); (c) first differences. Compare and interpret the results.

\textbf{10.11} Apply the Hausman test to your estimates. What do you conclude about the appropriate model?

\medskip\hrule\medskip

\section{R Lab 10: Fixed Effects Estimation}

\subsection{Setup}

\begin{lstlisting}
library(tidyverse)
library(plm)     # panel data estimation
library(lmtest)
library(sandwich)

# Install and load data
# install.packages("plm")
data("Wages", package = "plm")
head(Wages)
# Variables: id, time, lwage (log wage), ed (education), exp (experience),
#            union (union status), black, etc.
\end{lstlisting}

\subsection{Exercise 10.1: Pooled OLS}

\begin{lstlisting}
# Convert to panel data format
wages_panel <- pdata.frame(Wages, index = c("id", "time"))

# Pooled OLS
pool_model <- plm(lwage ~ ed + exp + I(exp^2) + union + black,
                   data = wages_panel, model = "pooling")
summary(pool_model)
\end{lstlisting}

\subsection{Exercise 10.2: Fixed Effects (Within Estimator)}

\begin{lstlisting}
# Fixed Effects (Within) Model
fe_model <- plm(lwage ~ ed + exp + I(exp^2) + union + black,
                 data = wages_panel, model = "within")
summary(fe_model)
cat("FE coefficient on union:", round(coef(fe_model)["union"], 4), "\n")
cat("Pool OLS coefficient on union:", round(coef(pool_model)["union"], 4), "\n")
\end{lstlisting}

\subsection{Exercise 10.3: First Differences}

\begin{lstlisting}
# First Differences
fd_model <- plm(lwage ~ ed + exp + I(exp^2) + union + black,
                 data = wages_panel, model = "fd")
summary(fd_model)
\end{lstlisting}

\subsection{Exercise 10.4: LSDV (verify FE equivalence)}

\begin{lstlisting}
# LSDV using lm() -- add factor(id) as dummy variables
# WARNING: this is slow for large N!
# Use a small subset for illustration
wages_small <- Wages[Wages$id <= 50, ]

lsdv_model <- lm(lwage ~ ed + exp + I(exp^2) + union + factor(id),
                   data = wages_small)
fe_small <- plm(lwage ~ ed + exp + I(exp^2) + union,
                 data = pdata.frame(wages_small, index = c("id", "time")),
                 model = "within")

cat("LSDV union coefficient:", round(coef(lsdv_model)["union"], 6), "\n")
cat("FE union coefficient:  ", round(coef(fe_small)["union"], 6), "\n")
\end{lstlisting}

\subsection{Exercise 10.5: Random Effects and Hausman Test}

\begin{lstlisting}
# Random Effects
re_model <- plm(lwage ~ ed + exp + I(exp^2) + union + black,
                 data = wages_panel, model = "random")
summary(re_model)

# Hausman Test
phtest(fe_model, re_model)
\end{lstlisting}

\subsection{Exercise 10.6: Two-Way Fixed Effects}

\begin{lstlisting}
# Two-Way FE (individual + time effects)
twofe_model <- plm(lwage ~ ed + exp + I(exp^2) + union + black,
                    data = wages_panel, model = "within", effect = "twoways")
summary(twofe_model)
\end{lstlisting}

\medskip\hrule\medskip

\textit{End of Chapter 10}

\medskip\hrule\medskip

\textbf{References}

Card, D. (1995). Using geographic variation in college proximity to estimate the return to schooling. University of Toronto Press.

Wooldridge, J. M. (2019). \textit{Introductory Econometrics}. Chapter 14.

Angrist, J. D., and Pischke, J. (2009). \textit{Mostly Harmless Econometrics}. Chapter 5.

Cunningham, S. (2021). \textit{Causal Inference: The Mixtape}. Chapter on panel data.

Hausman, J. A. (1978). Specification tests in econometrics. \textit{Econometrica}, 46(6), 1251--1271.

\chapter{Chapter 11: Differences-in-Differences}

\medskip\hrule\medskip

\section{Chapter Overview}

\textbf{Differences-in-Differences (DiD)} is the workhorse identification strategy of modern empirical economics. It exploits variation in the timing and geography of policy changes to estimate causal effects: we compare the change in outcomes over time for units that were exposed to a treatment (the "treatment group") to the change in outcomes over time for units that were not (the "control group"). The difference of these two differences is the DiD estimator.

The key identifying assumption is \textbf{parallel trends}: in the absence of treatment, the treatment and control groups would have followed the same trend over time. When this assumption holds, the DiD estimator is unbiased. Testing it --- and understanding when it is plausible --- is the central challenge of DiD research.

\medskip\hrule\medskip

\section{Learning Objectives}

After completing this chapter, students will be able to:

\begin{enumerate}
\item Set up the \textbf{2$\times$2 DiD framework} (two groups, two periods).
\item State the \textbf{parallel trends assumption} and explain its meaning.
\item Derive the \textbf{DiD estimator} from simple mean comparisons.
\item Formulate DiD as a \textbf{regression with interaction terms}.
\item Explain the role of \textbf{clustering} standard errors in DiD.
\item Use \textbf{event study plots} to assess parallel pre-trends.
\item Identify key threats to DiD validity: anticipation, policy endogeneity, compositional changes.
\item Describe the complications of \textbf{staggered DiD} and recent methodological advances.
\item Apply DiD to the Card and Krueger (1994) minimum wage study.
\end{enumerate}

\medskip\hrule\medskip

\section{11.1 The 2$\times$2 Differences-in-Differences Setup}

\subsection{11.1.1 Notation}

Consider two groups and two time periods:
\begin{itemize}
\item \textbf{Groups:} $g \in \{0, 1\}$ where $g = 1$ is the treatment group, $g = 0$ is the control group.
\item \textbf{Periods:} $t \in \{0, 1\}$ where $t = 0$ is the pre-treatment period, $t = 1$ is the post-treatment period.
\item \textbf{Treatment:} Group 1 receives the treatment in period 1; group 0 never receives treatment.
\end{itemize}

We define potential outcomes $Y_{gt}(d)$: the mean outcome in group $g$, period $t$, if treatment status were $d$.

\subsection{11.1.2 The Estimand}

The object of interest is the \textbf{Average Treatment Effect on the Treated (ATT)} in the post-period:

\[
\text{ATT} = E[Y_{i1}(1) - Y_{i1}(0) \mid g_i = 1]
\]

This is the average effect of treatment on the treated units in the post-treatment period. The key challenge: we do not observe $Y_{i1}(0)$ for treated units --- we need a counterfactual for what the treatment group's outcome would have been had they not been treated.

\subsection{11.1.3 The Four Observed Cell Means}

\begin{center}
\begin{tabular}{@{} l l l @{}}
\toprule
\textbf{Period} & \textbf{Control Group ($g=0$)} & \textbf{Treatment Group ($g=1$)} \\
\midrule
Pre ($t=0$) & $\bar{Y}_{00}$ & $\bar{Y}_{10}$ \\
Post ($t=1$) & $\bar{Y}_{01}$ & $\bar{Y}_{11}$ \\
\bottomrule
\end{tabular}
\end{center}

\begin{figure}[htbp]
  \centering
  \begin{tikzpicture}
    \begin{axis}[
      width=10cm, height=7cm,
      xlabel={Time period},
      ylabel={Outcome $Y$},
      xmin=0.3, xmax=2.7,
      ymin=0.5, ymax=5.5,
      axis lines=left,
      tick style={draw=none},
      xtick={1,2},
      xticklabels={Pre ($t=0$), Post ($t=1$)},
      ytick=\empty,
      legend style={at={(0.05,0.95)}, anchor=north west, font=\small},
    ]
      % Control group: goes from 1.5 to 2.5
      \addplot[very thick, gray!60, mark=*, mark size=3pt]
        coordinates {(1,1.5)(2,2.5)};
      % Treatment group (actual): goes from 2.5 to 4.5
      \addplot[very thick, econblue, mark=square*, mark size=3pt]
        coordinates {(1,2.5)(2,4.5)};
      % Counterfactual (parallel trend): goes from 2.5 to 3.5
      \addplot[thick, econblue, dashed, mark=square, mark size=3pt]
        coordinates {(1,2.5)(2,3.5)};

      % DiD bracket
      \draw[<->, very thick, red!70!black] (axis cs:2.15, 3.5) -- (axis cs:2.15, 4.5)
        node[midway, right, font=\small\bfseries, red!70!black] {DiD $= \hat{\delta}$};

      % Parallel trend arrow
      \draw[->, gray, dashed, thick] (axis cs:1.5,2.2) to[bend right=10] (axis cs:1.85, 3.2);
      \node[font=\footnotesize, gray, align=center] at (axis cs:1.2, 2.0) {parallel\\trend};

      % Treatment timing
      \draw[gray!50, dashed, thick] (axis cs:1.5, 0.5) -- (axis cs:1.5, 5.5)
        node[above, font=\footnotesize, gray] {Treatment};

      \legend{Control group, Treatment group (actual), Treatment group (counterfactual)}
    \end{axis}
  \end{tikzpicture}
  \caption{The 2$\times$2 Differences-in-Differences design. Both groups trend upward by 1 unit from pre to post under parallel trends. The treatment group (blue) actually rises by 2 units, so the DiD estimate $\hat{\delta} = (4.5-2.5)-(2.5-1.5) = 1$ captures the treatment effect. The dashed blue line shows the counterfactual trend the treatment group would have followed absent treatment.}
  \label{fig:ch11_did_diagram}
\end{figure}

\medskip\hrule\medskip

\section{11.2 The Parallel Trends Assumption}

\textbf{Assumption 11.1 (Parallel Trends):}

\[
E[Y_{t}(0) \mid g=1, t=1] - E[Y_{t}(0) \mid g=1, t=0] = E[Y_{t}(0) \mid g=0, t=1] - E[Y_{t}(0) \mid g=0, t=0]
\]

In words: \textbf{in the absence of treatment, the treatment group would have experienced the same change in outcomes as the control group.}

This does NOT require the two groups to start at the same level. It only requires them to share the same time trend in untreated potential outcomes.

\textbf{Visual test:} Plot outcome trends for both groups in the pre-treatment periods. If the two groups have similar pre-trends (parallel), the assumption is more credible.

\medskip\hrule\medskip

\section{11.3 The DiD Estimator}

\subsection{11.3.1 Deriving the Estimator}

The observed outcome for a treated unit in the post-period is:

\[
E[Y_{i1} \mid g_i = 1, t=1] = E[Y_{i1}(1) \mid g_i = 1]
\]

We need $E[Y_{i1}(0) \mid g_i = 1]$ --- the counterfactual. Under parallel trends:

\[
E[Y_{i1}(0) \mid g_i = 1] = E[Y_{i0}(0) \mid g_i = 1] + \underbrace{(E[Y_{t}(0) \mid g=0, t=1] - E[Y_{t}(0) \mid g=0, t=0])}_{\text{control group trend}}
\]

\[
= \bar{Y}_{10} + (\bar{Y}_{01} - \bar{Y}_{00})
\]

Therefore:

\[
\text{ATT} = E[Y_{i1}(1) \mid g_i=1] - E[Y_{i1}(0) \mid g_i=1]
\]

\[
= \bar{Y}_{11} - \left[\bar{Y}_{10} + (\bar{Y}_{01} - \bar{Y}_{00})\right]
\]

\[
= \boxed{(\bar{Y}_{11} - \bar{Y}_{10}) - (\bar{Y}_{01} - \bar{Y}_{00})}
\]

\textbf{This is the DiD estimator:} the difference in before-after changes between the treatment and control groups.

\subsection{11.3.2 Decomposing the DiD}

\[
\hat{\delta}^{DiD} = \underbrace{(\bar{Y}_{11} - \bar{Y}_{10})}_{\text{Change in treatment group}} - \underbrace{(\bar{Y}_{01} - \bar{Y}_{00})}_{\text{Change in control group}}
\]

\begin{center}
\begin{tabular}{@{} l l @{}}
\toprule
\textbf{Comparison} & \textbf{What it measures} \\
\midrule
$(\bar{Y}_{11} - \bar{Y}_{10})$ & Pre-post change for treated group = Treatment effect + Time trend \\
$(\bar{Y}_{01} - \bar{Y}_{00})$ & Pre-post change for control group = Time trend \\
DiD & Treatment effect (time trend differenced out) \\
\bottomrule
\end{tabular}
\end{center}

The control group provides the counterfactual trend. Subtracting it removes the common time trend, leaving only the causal effect.

\subsection{11.3.3 Identification Under Parallel Trends}

\textbf{Proposition 11.1:} Under Assumption 11.1 (Parallel Trends) and the overlap condition ($P(g=1) > 0$, $P(g=0) > 0$):

\[
\hat{\delta}^{DiD} \xrightarrow{p} \text{ATT}
\]

\medskip\hrule\medskip

\section{11.4 DiD as a Regression}

The DiD estimator has a natural regression formulation.

\textbf{Regression DiD:}

\[
Y_{igt} = \alpha + \beta_g G_i + \beta_t T_t + \delta (G_i \times T_t) + u_{igt}
\]

where:
\begin{itemize}
\item $G_i = 1$ if unit $i$ is in the treatment group
\item $T_t = 1$ if period $t$ is the post-treatment period
\item $G_i \times T_t$ is the interaction term
\end{itemize}

\textbf{Interpreting coefficients:}

\begin{center}
\begin{tabular}{@{} l l l @{}}
\toprule
\textbf{Group} & \textbf{Period} & \textbf{$E[Y]$} \\
\midrule
Control ($G=0$) & Pre ($T=0$) & $\alpha$ \\
Control ($G=0$) & Post ($T=1$) & $\alpha + \beta_t$ \\
Treatment ($G=1$) & Pre ($T=0$) & $\alpha + \beta_g$ \\
Treatment ($G=1$) & Post ($T=1$) & $\alpha + \beta_g + \beta_t + \delta$ \\
\bottomrule
\end{tabular}
\end{center}

Therefore:
\begin{itemize}
\item $\beta_g$: pre-existing level difference between treatment and control.
\item $\beta_t$: common time trend (post-pre change for the control group).
\item $\delta$: the \textbf{DiD estimator} --- the additional change for the treatment group in the post period = \textbf{ATT under parallel trends}.
\end{itemize}

\textbf{Verification:} $\delta = (\alpha + \beta_g + \beta_t + \delta) - (\alpha + \beta_t) - (\alpha + \beta_g) + \alpha = \hat{\delta}^{DiD}$. $\checkmark$

\subsection{11.4.1 With Additional Controls}

To increase precision and relax the parallel trends assumption:

\[
Y_{igt} = \alpha_g + \gamma_t + \delta (G_i \times T_t) + \mathbf{X}_{igt}^T \boldsymbol{\theta} + u_{igt}
\]

where $\alpha_g$ are group fixed effects, $\gamma_t$ are time fixed effects, and $\mathbf{X}_{igt}$ are time-varying controls.

\textbf{Note:} With multiple time periods and groups, this becomes the \textbf{two-way fixed effects DiD} model of Chapter 10.

\medskip\hrule\medskip

\section{11.5 Standard Errors in DiD: Clustering}

Standard errors in DiD settings need careful treatment. The treatment is typically assigned at the group level (e.g., a state adopts a policy), but outcomes are measured at the individual level. This creates \textbf{within-group correlation} in residuals that, if ignored, leads to severely underestimated standard errors.

\textbf{Bertrand, Duflo, and Mullainathan (2004)} showed that ignoring clustering can lead to rejection rates of 45\% at the nominal 5\% level when the true effect is zero.

\textbf{Solution:} Always cluster standard errors at the level of treatment assignment (e.g., state, school, firm).

\textbf{Rule:} Cluster at the group $\times$ time level that the policy variation comes from. If treatment varies at the state level, cluster at the state level even if outcomes are at the individual level.

\medskip\hrule\medskip

\section{11.6 Testing Parallel Trends: Event Study Plots}

The parallel trends assumption is the key identifying assumption of DiD. While it cannot be directly tested (we never observe the counterfactual trend for the treatment group), we can assess its plausibility using \textbf{pre-treatment trends}.

\subsection{11.6.1 The Event Study Specification}

Replace the single post-period interaction with a set of interactions for each time period:

\[
Y_{igt} = \alpha_g + \gamma_t + \sum_{k \neq -1} \delta_k (G_i \times \mathbb{1}[t=k]) + \mathbf{X}_{igt}^T\boldsymbol{\theta} + u_{igt}
\]

where $k$ indexes periods relative to the treatment date, and $k = -1$ is the omitted (reference) period just before treatment.

\textbf{The $\delta_k$ coefficients:}
\begin{itemize}
\item For $k < 0$ (pre-treatment): should be zero if parallel trends holds.
\item For $k \geq 0$ (post-treatment): these are the dynamic treatment effects.
\end{itemize}

\textbf{Event study plot:} Plot $\hat{\delta}_k$ with confidence intervals against time $k$. If pre-treatment coefficients are near zero (statistically and practically), parallel trends is more plausible.

\begin{figure}[htbp]
  \centering
  \begin{tikzpicture}
    \begin{axis}[
      width=11cm, height=6cm,
      xlabel={Period relative to treatment ($k$)},
      ylabel={Coefficient $\hat{\delta}_k$ (relative to $k=-1$)},
      xmin=-4.5, xmax=4.5,
      ymin=-1.5, ymax=3.5,
      axis lines=left,
      tick style={draw=none},
      xtick={-4,-3,-2,-1,0,1,2,3,4},
      xticklabels={$-4$,$-3$,$-2$,$-1$,$0$,$1$,$2$,$3$,$4$},
      ytick={0,1,2,3},
      yticklabels={0,1,2,3},
      grid=none,
    ]
      % Zero reference line
      \addplot[gray!50, thick, dashed] coordinates {(-4.5,0)(4.5,0)};
      % Treatment timing
      \draw[gray!60, dashed, thick] (axis cs:-0.5,-1.5) -- (axis cs:-0.5,3.5)
        node[above right, font=\footnotesize, gray] {treatment};

      % Pre-trends (flat, near zero) with CI
      \addplot[only marks, mark=*, mark size=3pt, econblue]
        coordinates {(-4,-0.1)(-3,0.15)(-2,-0.05)(-1,0)};
      % Post-treatment effects (rising)
      \addplot[only marks, mark=*, mark size=3pt, red!70!black]
        coordinates {(0,1.1)(1,1.8)(2,2.3)(3,2.5)(4,2.4)};

      % CI bars pre (hardcoded endpoints = point +/- 0.35)
      \draw[econblue!60, thick] (axis cs:-4, -0.45) -- (axis cs:-4,  0.25);
      \draw[econblue!60, thick] (axis cs:-3, -0.20) -- (axis cs:-3,  0.50);
      \draw[econblue!60, thick] (axis cs:-2, -0.40) -- (axis cs:-2,  0.30);
      \draw[econblue!60, thick] (axis cs:-1, -0.35) -- (axis cs:-1,  0.35);
      % CI bars post (hardcoded endpoints = point +/- 0.40)
      \draw[red!70!black!60, thick] (axis cs:0,  0.70) -- (axis cs:0,  1.50);
      \draw[red!70!black!60, thick] (axis cs:1,  1.40) -- (axis cs:1,  2.20);
      \draw[red!70!black!60, thick] (axis cs:2,  1.90) -- (axis cs:2,  2.70);
      \draw[red!70!black!60, thick] (axis cs:3,  2.10) -- (axis cs:3,  2.90);
      \draw[red!70!black!60, thick] (axis cs:4,  2.00) -- (axis cs:4,  2.80);

      \node[font=\footnotesize, econblue] at (axis cs:-2.5, -1.1) {Pre-trend: flat $\checkmark$};
      \node[font=\footnotesize, red!70!black] at (axis cs:2.0, 0.5) {Post-treatment: rising};
    \end{axis}
  \end{tikzpicture}
  \caption{An event study (dynamic DiD) plot. Coefficients $\hat{\delta}_k$ measure the treatment group's outcome relative to the control group, normalized to zero at $k=-1$ (one period before treatment). Flat pre-treatment coefficients (blue) support the parallel trends assumption. Rising post-treatment coefficients (red) reveal the treatment effect growing over time.}
  \label{fig:ch11_event_study}
\end{figure}

\subsection{11.6.2 Caution: Pre-trends Tests Are Noisy}

Even if the parallel trends assumption holds, pre-trends tests have low power with few pre-treatment periods. Conversely, a non-significant pre-trend test does not guarantee that parallel trends holds. The test is necessary but not sufficient evidence.

\medskip\hrule\medskip

\section{11.7 Threats to DiD Validity}

\subsection{11.7.1 Violation of Parallel Trends}

The most common threat. If treatment group would have experienced a different trend even without treatment, the DiD estimate is biased.

\textbf{Example:} If firms that adopt a new technology (treatment) were on a faster growth trajectory before adoption, the DiD estimator confounds the treatment effect with the pre-existing growth differential.

\textbf{Mitigation:} Use multiple pre-treatment periods to document parallel pre-trends; use control groups that are more similar to the treatment group.

\subsection{11.7.2 Anticipation Effects}

If treated units change behavior in anticipation of the treatment before it officially begins, the pre-treatment outcomes are already affected by the treatment. This biases the DiD downward (understates the total effect).

\textbf{Example:} A minimum wage increase announced 6 months before it takes effect. Firms may start adjusting employment before the law takes effect.

\textbf{Mitigation:} Extend the pre-treatment window to see if anticipation effects appear.

\subsection{11.7.3 Policy Endogeneity}

The treatment is often adopted in response to current or expected outcomes. If states with declining fast-food employment adopt a minimum wage increase, the control group (states with stable employment) will show different trends.

\textbf{Mitigation:} Find settings where treatment assignment is exogenous --- natural experiments, sharp rules, randomized rollouts.

\subsection{11.7.4 Compositional Changes}

If the composition of the treatment and control groups changes over time (e.g., selective migration in response to a policy), the pre-post comparison may reflect composition changes rather than treatment effects.

\medskip\hrule\medskip

\section{11.8 Staggered DiD and Recent Advances}

\subsection{11.8.1 Staggered Treatment Adoption}

In many real settings, different units adopt treatment at different times (staggered rollout). For example, US states adopt a minimum wage increase in different years.

The two-way FE estimator in this setting can be written as:

\[
Y_{it} = \alpha_i + \gamma_t + \delta D_{it} + u_{it}
\]

where $D_{it} = 1$ if unit $i$ is treated by period $t$.

\textbf{Problem:} Recent work (Callaway and Sant'Anna, 2021; Goodman-Bacon, 2021; de Chaisemartin and D'Haultfoeuille, 2020) has shown that in staggered DiD settings, the two-way FE $\delta$ is a weighted average of 2$\times$2 DiD estimates, but with \textbf{potentially negative weights}. If treatment effects are heterogeneous (vary over time or across groups), the two-way FE estimate can be a misleading average or even have the wrong sign.

\subsection{11.8.2 Solutions}

Modern approaches decompose the staggered DiD into a series of 2$\times$2 comparisons:

\begin{itemize}
\item \textbf{Callaway and Sant'Anna (2021):} Compute group-time average treatment effects $ATT(g,t)$ for each adoption cohort $g$ and time period $t$, then aggregate.
\item \textbf{Sun and Abraham (2021):} Introduce a fully saturated interaction model.
\item \textbf{de Chaisemartin and D'Haultfoeuille (2020):} Propose the "did\_multiplegt" estimator robust to heterogeneous effects.
\end{itemize}

\medskip\hrule\medskip

\section{11.9 Empirical Illustration: Card and Krueger (1994)}

\subsection{11.9.1 The Research Question}

In April 1992, New Jersey raised its minimum wage from \$4.25 to \$5.05 per hour. Pennsylvania did not change its minimum wage. Card and Krueger (1994) use this natural experiment to study the effect of the minimum wage increase on fast-food employment.

\textbf{Treatment group:} Fast-food restaurants in New Jersey.  
\textbf{Control group:} Fast-food restaurants in eastern Pennsylvania.  
\textbf{Pre-treatment period:} February 1992.  
\textbf{Post-treatment period:} November 1992.

\subsection{11.9.2 Identification}

The parallel trends assumption requires that, absent the New Jersey minimum wage increase, employment trends in New Jersey and Pennsylvania fast-food restaurants would have been similar.

\textbf{Plausibility:} New Jersey and Pennsylvania are geographically contiguous, compete in similar labor markets, and have similar industry mix. This makes parallel trends more plausible than comparing New Jersey to, say, California.

\subsection{11.9.3 Results}

\begin{center}
\begin{tabular}{@{} l l l l @{}}
\toprule
\textbf{} & \textbf{New Jersey (treatment)} & \textbf{Pennsylvania (control)} & \textbf{DiD} \\
\midrule
Feb (pre) & 20.44 employees & 23.33 employees & --- \\
Nov (post) & 21.03 employees & 21.17 employees & --- \\
Change & +0.59 & $-$2.16 & \textbf{+2.76} \\
\bottomrule
\end{tabular}
\end{center}

\textbf{DiD estimate:} $0.59 - (-2.16) = 2.76$ additional full-time equivalent employees in New Jersey restaurants.

This suggests the minimum wage increase actually \textit{increased} employment in New Jersey fast-food restaurants --- directly contrary to the standard competitive model prediction.

\textbf{Why?} Card and Krueger argue: (1) fast-food labor markets are monopsonistic, so minimum wages can increase employment; (2) the positive effect might reflect the positive income effect on low-wage households increasing demand for fast food.

\subsection{11.9.4 Criticism and Response}

\textbf{Neumark and Wascher (1994)} challenged the results using payroll data (vs. Card-Krueger's phone survey data). Their estimates showed negative employment effects.

\textbf{Methodological lesson:} The choice of data source, control group, and period significantly affects DiD estimates. Always subject your design to sensitivity checks.

\medskip\hrule\medskip

\section{Chapter Summary}

\begin{itemize}
\item \textbf{DiD} compares the change in outcomes before and after treatment for treated units versus control units.
\item The \textbf{parallel trends assumption} is the key identifying assumption: control group trend serves as the counterfactual for the treatment group.
\item \textbf{DiD regression:} $Y = \alpha + \beta_g G + \beta_t T + \delta(G \times T) + u$; the coefficient $\delta$ is the ATT.
\item \textbf{Event study plots} assess pre-treatment parallel trends --- necessary but not sufficient evidence.
\item \textbf{Cluster standard errors} at the level of treatment assignment; ignoring clustering severely understates uncertainty.
\item Key threats: violations of parallel trends, anticipation effects, policy endogeneity, compositional changes.
\item \textbf{Staggered DiD} settings require caution: two-way FE can have wrong-sign weights under heterogeneous effects.
\item \textbf{Card and Krueger (1994)} is the canonical DiD application, showing minimum wages may not reduce fast-food employment.
\end{itemize}

\medskip\hrule\medskip

\section{Key Terms}

\textbf{Differences-in-Differences (DiD)} --- $(\bar{Y}_{11} - \bar{Y}_{10}) - (\bar{Y}_{01} - \bar{Y}_{00})$; the change in outcomes for the treated group minus the change for the control group.

\textbf{Parallel trends assumption} --- in the absence of treatment, treated and control groups would have experienced the same trend.

\textbf{ATT (DiD context)} --- $E[Y_{i1}(1) - Y_{i1}(0) \mid g_i = 1]$; effect of treatment on treated units in the post-period.

\textbf{Event study} --- regression estimating dynamic treatment effects before and after treatment.

\textbf{Pre-trends test} --- testing whether pre-treatment DiD coefficients are zero; assesses parallel trends plausibility.

\textbf{Staggered DiD} --- DiD with variation in timing of treatment adoption across units.

\medskip\hrule\medskip

\section{Exercises}

\subsection{Conceptual Questions}

\textbf{11.1} State the parallel trends assumption in your own words. Why is it called "parallel"?

\textbf{11.2} Explain why the DiD estimator is not simply the before-after change for the treatment group.

\textbf{11.3} A researcher uses DiD to study the effect of a state-level policy on individual outcomes. She has 50 states and 10,000 individual observations. Should she cluster at the individual or state level? Why?

\textbf{11.4} Describe an event study plot. What would you conclude from (a) flat pre-trends with a jump at treatment; (b) rising pre-trends in the treatment group before treatment?

\textbf{11.5} Explain the "negative weights" problem in staggered DiD. Why is the two-way FE estimator potentially misleading?

\subsection{Analytical Questions}

\textbf{11.6} Using the Card-Krueger data in the table above, compute the DiD estimate by hand.

\textbf{11.7} Show that the OLS coefficient on $G_i \times T_t$ in the regression $Y = \alpha + \beta_g G + \beta_t T + \delta (G \times T) + u$ equals the DiD estimator.

\textbf{11.8} Suppose parallel trends fails: the treatment group would have grown by 3\% in the post-period even without treatment, while the control group grows by 1\%. The observed DiD is 5\%. What is the true ATT?

\subsection{Applied Questions}

\textbf{11.9} Using publicly available Card-Krueger data, replicate their DiD estimate in R.

\textbf{11.10} Generate event study plots for a dataset with multiple pre-treatment periods.

\medskip\hrule\medskip

\section{R Lab 11: Differences-in-Differences}

\subsection{Setup}

\begin{lstlisting}
library(tidyverse)
library(fixest)   # fast fixed effects with robust SE
library(ggplot2)

# Simulate DiD data
set.seed(42)
n_units <- 200    # 100 treatment, 100 control
n_periods <- 6    # periods -3, -2, -1 (pre) and 0, 1, 2 (post)

df <- expand.grid(
  unit = 1:n_units,
  period = -3:2
) %>%
  as_tibble() %>%
  mutate(
    treated    = unit > 100,                     # 101-200 are treated
    post       = period >= 0,                    # periods 0,1,2 are post
    time_trend = 0.5 * period,                   # common linear time trend
    unit_fe    = rnorm(n_units)[unit],           # unit-specific fixed effect
    te         = ifelse(treated & post, 3, 0),   # true ATT = 3
    noise      = rnorm(n_units * n_periods, 0, 1.5),
    Y          = 10 + unit_fe + time_trend + te + noise
  )
\end{lstlisting}

\subsection{Exercise 11.1: Visualizing Parallel Trends}

\begin{lstlisting}
df %>%
  group_by(treated, period) %>%
  summarise(mean_Y = mean(Y), .groups = "drop") %>%
  ggplot(aes(x = period, y = mean_Y, color = treated, group = treated)) +
  geom_line(linewidth = 1.2) +
  geom_point(size = 3) +
  geom_vline(xintercept = -0.5, linetype = "dashed", color = "gray50") +
  labs(
    title    = "Outcome Trends: Treatment vs. Control Group",
    subtitle = "Dashed line = treatment date. Pre-trends look parallel.",
    x = "Period", y = "Mean Outcome",
    color = "Treated"
  ) +
  theme_minimal()
\end{lstlisting}

\subsection{Exercise 11.2: The 2$\times$2 DiD}

\begin{lstlisting}
# Use only pre (-1) and post (0) periods
df_2x2 <- df %>% filter(period %in% c(-1, 0))

# Cell means
cell_means <- df_2x2 %>%
  group_by(treated, post) %>%
  summarise(mean_Y = mean(Y), .groups = "drop")
print(cell_means)

# DiD by hand
did_manual <- with(cell_means,
  (mean_Y[treated == TRUE  & post == TRUE]  -
   mean_Y[treated == TRUE  & post == FALSE]) -
  (mean_Y[treated == FALSE & post == TRUE]  -
   mean_Y[treated == FALSE & post == FALSE])
)
cat("DiD estimate (manual):", round(did_manual, 3), "\n")

# DiD via regression
did_reg <- lm(Y ~ treated * post, data = df_2x2)
summary(did_reg)
\end{lstlisting}

\subsection{Exercise 11.3: Regression DiD with Fixed Effects}

\begin{lstlisting}
# Full panel DiD with unit and time fixed effects
fe_did <- feols(Y ~ i(post, treated) | unit + period,
                data = df, cluster = ~unit)
summary(fe_did)
cat("FE DiD estimate:", round(coef(fe_did)["post::TRUE:treatedTRUE"], 3), "\n")
\end{lstlisting}

\subsection{Exercise 11.4: Event Study Plot}

\begin{lstlisting}
# Event study: dynamic treatment effects by period
event_study <- feols(Y ~ i(period, treated, ref = -1) | unit + period,
                     data = df, cluster = ~unit)

iplot(event_study,
      xlab = "Periods Relative to Treatment",
      main = "Event Study: Dynamic Treatment Effects",
      sub  = "Reference period = -1. Treatment begins at period 0.")
\end{lstlisting}

\subsection{Exercise 11.5: Card and Krueger Replication}

\begin{lstlisting}
# Read Card-Krueger data (available at various sources)
# Here we use the 'wooldridge' package
# install.packages("wooldridge")
library(wooldridge)
data("fastfood", package = "wooldridge")

# DiD regression
ck_model <- lm(fte ~ nj + d + nj:d, data = fastfood)
summary(ck_model)
cat("DiD estimate (Card-Krueger):", round(coef(ck_model)["nj:d"], 3), "\n")
\end{lstlisting}

\medskip\hrule\medskip

\textit{End of Chapter 11}

\medskip\hrule\medskip

\textbf{References}

Card, D., and Krueger, A. B. (1994). Minimum wages and employment: A case study of the fast-food industry in New Jersey and Pennsylvania. \textit{American Economic Review}, 84(4), 772--793.

Bertrand, M., Duflo, E., and Mullainathan, S. (2004). How much should we trust differences-in-differences estimates? \textit{Quarterly Journal of Economics}, 119(1), 249--275.

Angrist, J. D., and Pischke, J. (2009). \textit{Mostly Harmless Econometrics}. Chapter 5.

Callaway, B., and Sant'Anna, P. H. C. (2021). Difference-in-differences with multiple time periods. \textit{Journal of Econometrics}, 225(2), 200--230.

Goodman-Bacon, A. (2021). Difference-in-differences with variation in treatment timing. \textit{Journal of Econometrics}, 225(2), 254--277.

Cunningham, S. (2021). \textit{Causal Inference: The Mixtape}. Chapter 9.

\chapter{Chapter 12: Matching and Selection on Observables}

\medskip\hrule\medskip

\section{Chapter Overview}

Matching estimators provide an alternative to regression for estimating average treatment effects when treatment is not randomly assigned. The key assumption is \textbf{selection on observables (CIA --- Conditional Independence Assumption)}: after conditioning on observed covariates, treatment is independent of potential outcomes. This chapter develops the theory and practice of matching, connects it to regression, and introduces propensity score methods.

\medskip\hrule\medskip

\section{Learning Objectives}

After completing this chapter, students will be able to:

\begin{enumerate}
\item State the \textbf{Conditional Independence Assumption (CIA)} and explain what it requires.
\item Explain the \textbf{propensity score theorem} and why it simplifies matching.
\item Describe and implement \textbf{nearest-neighbor matching}.
\item Explain the connection between regression and matching.
\item Apply the \textbf{propensity score} for matching, weighting, and regression.
\item Check the \textbf{overlap (common support)} condition and deal with violations.
\item Evaluate the quality of matching using \textbf{balance checks}.
\item Discuss when matching is more or less appropriate than other methods.
\end{enumerate}

\medskip\hrule\medskip

\section{12.1 The Conditional Independence Assumption (CIA)}

\textbf{Definition 12.1 (Conditional Independence Assumption / Ignorability):}

\[
\{Y_{0i}, Y_{1i}\} \perp\!\!\!\perp D_i \mid \mathbf{X}_i
\]

Given observed covariates $\mathbf{X}_i$, the treatment $D_i$ is independent of potential outcomes. Conditioning on $\mathbf{X}_i$ "explains away" all selection into treatment.

This implies:

\[
E[Y_{0i} \mid D_i = 1, \mathbf{X}_i = \mathbf{x}] = E[Y_{0i} \mid D_i = 0, \mathbf{X}_i = \mathbf{x}]
\]

Within cells defined by $\mathbf{X}_i = \mathbf{x}$, treatment is effectively random.

\textbf{When is CIA plausible?}

The CIA is plausible when:
\begin{itemize}
\item Treatment assignment is driven by observable characteristics (e.g., eligibility rules based on measurable criteria).
\item The researcher has rich data capturing all important determinants of selection.
\item There are no hidden selection mechanisms (no selection on unobservables).
\end{itemize}

\textbf{When is CIA implausible?}
\begin{itemize}
\item Treatment is self-selected based on unobservable motivation or ability.
\item Assignment rules are discretionary and based on unrecorded factors.
\item There is reverse causality.
\end{itemize}

\textbf{Example 12.1:} A job training program recruits participants based on age, education, and prior employment status --- all observable. If the data contains these variables, CIA may be plausible: conditional on observables, participants and non-participants are comparable.

\medskip\hrule\medskip

\section{12.2 The Overlap Condition}

\textbf{Assumption 12.1 (Overlap / Common Support):}

\[
0 < P(D_i = 1 \mid \mathbf{X}_i = \mathbf{x}) < 1 \quad \text{for all } \mathbf{x} \text{ in the support of } \mathbf{X}
\]

Every individual, regardless of their covariate values, has a positive probability of being in both the treatment and control groups.

Overlap fails when:
\begin{itemize}
\item Some covariate values are found only among treated units (no valid control group members to compare to).
\item Perfect predictors of treatment exist.
\end{itemize}

\textbf{Dealing with overlap violations:}
\begin{itemize}
\item Trim the sample to the common support region.
\item Use only observations where propensity scores overlap between treatment and control.
\item Report results with and without trimming; check sensitivity.
\end{itemize}

\medskip\hrule\medskip

\section{12.3 Sub-classification and Direct Matching}

\subsection{12.3.1 Sub-classification}

The simplest matching approach: divide the sample into cells defined by discrete values of $\mathbf{X}_i$, estimate the treatment effect within each cell, then average across cells.

\[
\hat{\tau}^{sub} = \sum_{\mathbf{x}} \hat{w}(\mathbf{x}) \left[\bar{Y}_{1}(\mathbf{x}) - \bar{Y}_{0}(\mathbf{x})\right]
\]

where $\bar{Y}_{d}(\mathbf{x})$ is the mean outcome for group $d$ in cell $\mathbf{x}$, and $\hat{w}(\mathbf{x})$ is the fraction of the treatment group in cell $\mathbf{x}$.

\textbf{Problem:} With many continuous covariates, there may be no exact matches (the "curse of dimensionality"). Propensity score matching is the solution.

\subsection{12.3.2 Nearest-Neighbor Matching}

Match each treated unit to the control unit(s) with the most similar covariates (measured by distance in $\mathbf{X}$-space):

\[
\hat{\tau}^{NNM} = \frac{1}{N_1} \sum_{i: D_i = 1} \left[Y_i - \frac{1}{|J(i)|}\sum_{j \in J(i)} Y_j\right]
\]

where $J(i)$ is the set of matched control units for treated unit $i$.

\textbf{Key choices in nearest-neighbor matching:}
\begin{itemize}
\item \textbf{Distance metric:} Euclidean? Mahalanobis?
\item \textbf{Number of matches $M$:} 1-to-1 or 1-to-$M$? More matches reduce variance but increase bias.
\item \textbf{With or without replacement:} Matching with replacement reduces bias; without replacement is simpler.
\item \textbf{Caliper:} Restrict matches to be within a maximum distance.
\end{itemize}

\medskip\hrule\medskip

\section{12.4 The Propensity Score Theorem}

With many covariates, exact matching is impossible. Rosenbaum and Rubin (1983) showed that matching on a single scalar --- the \textbf{propensity score} --- suffices.

\textbf{Definition 12.2 (Propensity Score):}

\[
p(\mathbf{x}) = P(D_i = 1 \mid \mathbf{X}_i = \mathbf{x})
\]

The propensity score is the conditional probability of receiving treatment given observed covariates.

\textbf{Theorem 12.1 (Propensity Score Theorem, Rosenbaum and Rubin 1983):} If the CIA holds with $\mathbf{X}_i$, then it also holds conditional on $p(\mathbf{X}_i)$:

\[
\{Y_{0i}, Y_{1i}\} \perp\!\!\!\perp D_i \mid p(\mathbf{X}_i)
\]

\textbf{Why this is powerful:} Instead of matching on a $k$-dimensional covariate vector, we only need to match on a single number --- the propensity score.

\subsection{12.4.1 Estimating the Propensity Score}

In practice, $p(\mathbf{x})$ is unknown and must be estimated. Common methods:
\begin{itemize}
\item \textbf{Logistic regression:} $\hat{p}(\mathbf{x}) = \frac{\exp(\mathbf{x}^T\hat{\boldsymbol{\gamma}})}{1 + \exp(\mathbf{x}^T\hat{\boldsymbol{\gamma}})}$
\item \textbf{Probit regression}
\item \textbf{Machine learning:} Random forests, LASSO (for high-dimensional $\mathbf{X}$)
\end{itemize}

\begin{figure}[htbp]
  \centering
  \begin{tikzpicture}
    \begin{axis}[
      width=10cm, height=6cm,
      xlabel={Propensity Score $\hat{p}(\mathbf{X}_i) = P(D_i=1|\mathbf{X}_i)$},
      ylabel={Density},
      xmin=0, xmax=1,
      ymin=0, ymax=4.5,
      axis lines=left,
      tick style={draw=none},
      xtick={0,0.15,0.8,1},
      xticklabels={0,,, 1},
      ytick=\empty,
      legend style={at={(0.97,0.95)}, anchor=north east, font=\small},
    ]
      % Control group distribution (left-skewed)
      \addplot[very thick, gray!60, fill=gray!15, fill opacity=0.5, domain=0:1, samples=100]
        {3.5*exp(-0.5*((x-0.25)/0.18)^2)/sqrt(2*pi*0.18^2)} \closedcycle;

      % Treatment group distribution (right-skewed)
      \addplot[very thick, econblue, fill=econblue!20, fill opacity=0.5, domain=0:1, samples=100]
        {3.5*exp(-0.5*((x-0.65)/0.18)^2)/sqrt(2*pi*0.18^2)} \closedcycle;

      % Common support region
      \draw[<->, very thick, red!70!black] (axis cs:0.28, 3.8) -- (axis cs:0.80, 3.8)
        node[midway, above, font=\small\bfseries, red!70!black] {Common Support};
      \draw[dashed, red!70!black, thick] (axis cs:0.28, 0) -- (axis cs:0.28, 3.8);
      \draw[dashed, red!70!black, thick] (axis cs:0.80, 0) -- (axis cs:0.80, 3.8);

      \legend{Control ($D_i=0$), Treated ($D_i=1$)}
    \end{axis}
  \end{tikzpicture}
  \caption{Propensity score distributions for treated and control groups. The \emph{common support} (overlap) region (between the red dashed lines) is where both treated and control units exist at similar propensity scores. Matching and IPW estimators are only valid in the common support region; observations outside it must be trimmed or reweighted.}
  \label{fig:ch12_common_support}
\end{figure}

\medskip\hrule\medskip

\section{12.5 Propensity Score Matching Estimators}

\subsection{12.5.1 Propensity Score Nearest-Neighbor Matching}

Match each treated unit to the control unit with the closest propensity score:

\[
\hat{\tau}^{PSM} = \frac{1}{N_1}\sum_{i: D_i=1}\left[Y_i - Y_{j(i)}\right]
\]

where $j(i) = \arg\min_{j: D_j = 0} |p(\hat{\mathbf{X}}_i) - p(\hat{\mathbf{X}}_j)|$.

\subsection{12.5.2 Inverse Probability Weighting (IPW)}

An elegant alternative: weight observations by the inverse of the propensity score so that the reweighted control group "looks like" the treatment group.

\[
\hat{\tau}^{IPW}_{ATT} = \frac{1}{N_1}\sum_{i=1}^n\left[D_i Y_i - \frac{(1-D_i)\hat{p}(\mathbf{X}_i)Y_i}{1-\hat{p}(\mathbf{X}_i)}\right]
\]

\textbf{Intuition:} Control units with high propensity scores (similar to treated units) receive higher weight; those with very different characteristics receive low weight.

\subsection{12.5.3 Doubly Robust Estimation}

Modern practice combines propensity score weighting with regression --- the \textbf{doubly robust} (DR) estimator:

\[
\hat{\tau}^{DR}_{ATT} = \frac{1}{N_1}\sum_{i: D_i=1}\left[(Y_i - \hat{\mu}_0(\mathbf{X}_i))\right] + \frac{1}{N_1}\sum_{i: D_i=0}\frac{\hat{p}(\mathbf{X}_i)}{1-\hat{p}(\mathbf{X}_i)}\left[(Y_i - \hat{\mu}_0(\mathbf{X}_i))\right]
\]

where $\hat{\mu}_0(\mathbf{X}_i)$ is the predicted untreated outcome from a regression model.

\textbf{Double robustness:} The DR estimator is consistent if \textbf{either} the propensity score model \textbf{or} the outcome regression model is correctly specified --- but not necessarily both.

\medskip\hrule\medskip

\section{12.6 Regression as a Form of Matching}

Regression and matching are deeply related. Under the CIA, both identify the ATT (or ATE). The difference lies in how they handle the comparison:

\begin{itemize}
\item \textbf{Regression} fits a global linear model and uses extrapolation.
\item \textbf{Matching} compares treated units only to "nearby" controls (local comparison).
\end{itemize}

\textbf{Angrist and Pischke (2009)} show that OLS is a form of matching where each control observation receives a weight proportional to how much its $\mathbf{X}$ differs from the treatment group distribution.

\textbf{When does regression fail relative to matching?}

When there is poor overlap --- the treatment and control groups have very different covariate distributions --- regression extrapolates from the control group to predict untreated potential outcomes for treated units. This extrapolation can be badly wrong if the CEF is non-linear.

Matching avoids this by restricting comparisons to the common support.

\medskip\hrule\medskip

\section{12.7 Balance Checks and Assessment}

The CIA requires that treatment and control groups are comparable after matching. \textbf{Balance checks} verify that the matched groups have similar covariate distributions.

\textbf{Formal balance check:} Test whether the mean (and sometimes distribution) of each covariate differs between treated and matched controls.

\textbf{Standardized Difference:}

\[
\text{SDiff}_j = \frac{\bar{X}_{1j} - \bar{X}_{0j}}{\sqrt{(\hat{\sigma}^2_{1j} + \hat{\sigma}^2_{0j})/2}}
\]

Values below 0.1 are generally considered good balance.

\textbf{Love plot:} A visual display of standardized differences before and after matching.

\medskip\hrule\medskip

\section{12.8 Practical Guidance: When to Match?}

\begin{center}
\begin{tabular}{@{} l l @{}}
\toprule
\textbf{Situation} & \textbf{Recommended approach} \\
\midrule
Rich observable data, CIA plausible & Propensity score matching or IPW \\
Concern about model extrapolation & Match first, then regress on matched sample \\
Many covariates, CIA uncertain & Add controls in regression; test sensitivity \\
CIA implausible (unobservables matter) & Use IV, DiD, or RDD instead \\
Extreme overlap problems & Match with caliper; report trimmed results \\
\bottomrule
\end{tabular}
\end{center}

\medskip\hrule\medskip

\section{Chapter Summary}

\begin{itemize}
\item \textbf{Matching} identifies causal effects by comparing treated and control units with similar observed covariates.
\item The \textbf{CIA} (selection on observables) is the key identifying assumption: conditional on $\mathbf{X}$, treatment is independent of potential outcomes.
\item The \textbf{overlap condition} requires that for every covariate value, both treated and control units exist.
\item The \textbf{propensity score theorem} allows matching on a single scalar (the probability of treatment) instead of the full covariate vector.
\item \textbf{Doubly robust estimators} are consistent if either the propensity score model or the outcome model is correct.
\item \textbf{Regression is a form of matching} but extrapolates; matching restricts comparisons to the common support.
\item \textbf{Balance checks} verify that matching has achieved covariate balance.
\item CIA is a strong assumption. When unobservables matter, use IV, DiD, or RDD.
\end{itemize}

\medskip\hrule\medskip

\section{Key Terms}

\textbf{CIA (Conditional Independence Assumption)} --- treatment is independent of potential outcomes given observed covariates.

\textbf{Overlap} --- every unit has positive probability of receiving each treatment level.

\textbf{Propensity score} --- $p(\mathbf{x}) = P(D=1 \mid \mathbf{X}=\mathbf{x})$.

\textbf{Propensity score theorem} --- matching on the propensity score achieves CIA when matching on $\mathbf{X}$ does.

\textbf{Nearest-neighbor matching} --- each treated unit is matched to the control unit(s) with the most similar propensity score.

\textbf{Inverse Probability Weighting (IPW)} --- reweight control units by $p(\mathbf{X})/(1-p(\mathbf{X}))$.

\textbf{Doubly robust estimator} --- consistent if either the propensity score model or the outcome model is correct.

\textbf{Balance check} --- assessment of whether treatment and control groups have similar covariate distributions after matching.

\textbf{Standardized difference} --- normalized measure of covariate balance.

\medskip\hrule\medskip

\section{Exercises}

\subsection{Conceptual Questions}

\textbf{12.1} State the CIA. What does it require about the selection process? Give an example where it is plausible and one where it is not.

\textbf{12.2} Explain the propensity score theorem. Why is matching on the propensity score sufficient when CIA holds for the full vector $\mathbf{X}$?

\textbf{12.3} What is the overlap condition? What goes wrong when it fails?

\textbf{12.4} Explain double robustness. Why is it a desirable property?

\textbf{12.5} A researcher argues that regression and matching give the same estimates under CIA. Is this true? When might they differ?

\subsection{Applied Questions}

\textbf{12.6} In R, estimate the average treatment effect of union membership on wages using propensity score matching. Steps:
(a) Estimate a logistic regression for treatment (union = 1) on observable controls.
(b) Compute the propensity score.
(c) Check overlap.
(d) Implement nearest-neighbor matching.
(e) Assess balance.
(f) Compute the ATT.

\medskip\hrule\medskip

\section{R Lab 12: Matching and Propensity Scores}

\subsection{Setup}

\begin{lstlisting}
library(tidyverse)
library(MatchIt)    # matching
library(cobalt)     # balance assessment
library(WeightIt)   # IPW

data("lalonde", package = "MatchIt")  # LaLonde (1986) job training data
\end{lstlisting}

\subsection{Exercise 12.1: Descriptive Analysis}

\begin{lstlisting}
# Summary statistics by treatment status
lalonde %>%
  group_by(treat) %>%
  summarise(
    n       = n(),
    age     = mean(age),
    educ    = mean(educ),
    re78    = mean(re78),
    .groups = "drop"
  )
\end{lstlisting}

\subsection{Exercise 12.2: Propensity Score Estimation}

\begin{lstlisting}
# Estimate propensity score via logistic regression
ps_model <- glm(treat ~ age + educ + race + married + nodegree + re74 + re75,
                 data = lalonde, family = binomial)

lalonde$pscore <- predict(ps_model, type = "response")

# Check overlap
ggplot(lalonde, aes(x = pscore, fill = factor(treat))) +
  geom_density(alpha = 0.5) +
  labs(title = "Propensity Score Overlap", x = "Propensity Score", fill = "Treated") +
  theme_minimal()
\end{lstlisting}

\subsection{Exercise 12.3: Nearest-Neighbor Matching}

\begin{lstlisting}
m_out <- matchit(treat ~ age + educ + race + married + nodegree + re74 + re75,
                 data = lalonde, method = "nearest", ratio = 1)
summary(m_out)

# Balance
bal.tab(m_out)
love.plot(m_out, threshold = 0.1)
\end{lstlisting}

\subsection{Exercise 12.4: ATT from Matched Sample}

\begin{lstlisting}
# Compute ATT from matched data
matched_data <- match.data(m_out)
att_matching <- mean(matched_data$re78[matched_data$treat == 1]) -
                mean(matched_data$re78[matched_data$treat == 0])
cat("ATT (matching):", round(att_matching, 2), "\n")

# Compare to OLS
ols_att <- lm(re78 ~ treat + age + educ + race + married + nodegree + re74 + re75,
               data = lalonde)
cat("OLS estimate of treatment:", round(coef(ols_att)["treat"], 2), "\n")
\end{lstlisting}

\subsection{Exercise 12.5: IPW Estimator}

\begin{lstlisting}
w_out <- weightit(treat ~ age + educ + race + married + nodegree + re74 + re75,
                   data = lalonde, method = "ps", estimand = "ATT")
summary(w_out)

# ATT via IPW regression
library(jtools)
lm_ipw <- lm(re78 ~ treat, data = lalonde, weights = w_out$weights)
cat("IPW ATT:", round(coef(lm_ipw)["treat"], 2), "\n")
\end{lstlisting}

\medskip\hrule\medskip

\textit{End of Chapter 12}

\medskip\hrule\medskip

\textbf{References}

Rosenbaum, P. R., and Rubin, D. B. (1983). The central role of the propensity score in observational studies for causal effects. \textit{Biometrika}, 70(1), 41--55.

LaLonde, R. J. (1986). Evaluating the econometric evaluations of training programs with experimental data. \textit{American Economic Review}, 76(4), 604--620.

Angrist, J. D., and Pischke, J. (2009). \textit{Mostly Harmless Econometrics}. Chapter 3.

Cunningham, S. (2021). \textit{Causal Inference: The Mixtape}. Chapter 5.

Wooldridge, J. M. (2019). \textit{Introductory Econometrics}. Chapter 21.

\chapter{Chapter 13: Regression Discontinuity Design}

\medskip\hrule\medskip

\section{Chapter Overview}

\textbf{Regression Discontinuity Design (RDD)} exploits sharp thresholds in treatment assignment rules to identify causal effects. When treatment is determined by whether a continuous "running variable" crosses a threshold --- a test score above 70\% grants a scholarship; income below a cutoff qualifies for a social program --- we can compare outcomes just above and just below the threshold to estimate a causal effect. The key insight is that individuals just above and just below the threshold are nearly identical in all characteristics, except their treatment status.

RDD has become one of the most widely used and credible identification strategies in applied economics, partly because its identifying assumptions are relatively transparent and can be subjected to rigorous validity tests.

\medskip\hrule\medskip

\section{Learning Objectives}

After completing this chapter, students will be able to:

\begin{enumerate}
\item Explain the \textbf{sharp RDD} setup and the identifying assumption of \textbf{continuity}.
\item Identify the \textbf{treatment effect estimated by RDD} and explain why it is a local effect.
\item Derive the \textbf{RDD estimator} as a difference in regression limits at the cutoff.
\item Describe and implement \textbf{local linear regression} for RDD estimation.
\item Explain \textbf{bandwidth selection} and the bias-variance tradeoff.
\item Describe the \textbf{fuzzy RDD} and its relationship to IV.
\item Conduct \textbf{validity tests}: density test, covariate balance, placebo cutoffs.
\item Apply RDD to canonical examples: Thistlethwaite-Campbell (1960), regression in test-score-based policy contexts.
\end{enumerate}

\medskip\hrule\medskip

\section{13.1 The Regression Discontinuity Idea}

\subsection{13.1.1 Setup}

Let $R_i$ be the \textbf{running variable} (also called the "forcing variable" or "assignment variable"). Treatment $D_i$ is assigned based on whether $R_i$ crosses a threshold $c$:

\[
D_i = \mathbb{1}[R_i \geq c] \quad \text{(Sharp RDD)}
\]

\textbf{Examples:}
\begin{itemize}
\item Test scores (R) determine scholarship eligibility (D): threshold at 70\%.
\item Age (R) determines pension eligibility (D): threshold at age 65.
\item Birth weight (R) determines neonatal intensive care admission (D): threshold at 1500g.
\item Vote share (R) determines election winner (D): threshold at 50\%.
\end{itemize}

\subsection{13.1.2 The Local Nature of RDD}

RDD identifies a \textbf{local} causal effect --- the treatment effect for individuals near the threshold, not for the full population. This is both a limitation and a strength:

\begin{itemize}
\item \textbf{Limitation:} May not generalize to individuals far from the cutoff.
\item \textbf{Strength:} Near the cutoff, individuals just above and below are nearly identical --- making the identifying assumption very credible.
\end{itemize}

\medskip\hrule\medskip

\section{13.2 The Sharp RDD}

\subsection{13.2.1 The CEF Near the Cutoff}

For the sharp RDD, the average treatment effect at the threshold $c$ is:

\[
\tau_{RDD} = \lim_{r \downarrow c} E[Y_i \mid R_i = r] - \lim_{r \uparrow c} E[Y_i \mid R_i = r]
\]

This is the \textbf{jump} in the conditional mean of $Y$ at the threshold $c$.

\textbf{Definition 13.1 (Sharp RDD Estimand):} 

\[
\tau_{RDD} = E[Y_{1i} - Y_{0i} \mid R_i = c]
\]

This is the ATT for individuals at the threshold --- those right on the margin of receiving treatment.

\subsection{13.2.2 The Identifying Assumption: Continuity}

\textbf{Assumption 13.1 (Continuity):} The conditional expectations $E[Y_{0i} \mid R_i = r]$ and $E[Y_{1i} \mid R_i = r]$ are continuous in $r$ at $r = c$.

\textbf{Interpretation:} In the absence of the threshold, the outcome would change smoothly as the running variable changes. There is no other discontinuity in outcomes at $c$.

This rules out:
\begin{itemize}
\item Other policies that change at exactly $r = c$.
\item Individuals precisely manipulating $R_i$ to sort above or below $c$.
\item Discontinuous changes in any confounder at $c$.
\end{itemize}

\subsection{13.2.3 Identification}

Under Assumption 13.1:

\[
\lim_{r \downarrow c} E[Y_i \mid R_i = r] - \lim_{r \uparrow c} E[Y_i \mid R_i = r]
\]

\[
= E[Y_{1i} \mid R_i = c] - E[Y_{0i} \mid R_i = c] = \tau_{RDD}
\]

The key insight: by continuity, the average untreated potential outcome $E[Y_{0i} \mid R_i = r]$ is approximately the same for individuals just above and just below the threshold. So the observed jump in $Y$ at the cutoff isolates the treatment effect.

\begin{figure}[htbp]
  \centering
  \begin{tikzpicture}
    \begin{axis}[
      width=11cm, height=7cm,
      xlabel={Running Variable $R_i$},
      ylabel={Outcome $Y_i$},
      xmin=-5, xmax=5,
      ymin=0, ymax=6,
      axis lines=left,
      tick style={draw=none},
      xtick={0},
      xticklabels={$c$ (cutoff)},
      ytick=\empty,
      legend style={at={(0.05,0.95)}, anchor=north west, font=\small},
    ]
      % Cutoff line
      \draw[gray!60, dashed, thick] (axis cs:0,0) -- (axis cs:0,6);

      % Left of cutoff: control (D=0), smooth curve
      \addplot[very thick, gray!70, domain=-5:-0.05, samples=50]
        {1.5 + 0.2*x + 0.02*x^2};

      % Right of cutoff: treated (D=1), shifted up by ~1.5
      \addplot[very thick, econblue, domain=0.05:5, samples=50]
        {3.1 + 0.2*x + 0.02*x^2};

      % Scatter (control side)
      \addplot[only marks, mark=o, mark size=1.5pt, color=gray!50, opacity=0.7]
        coordinates {
          (-4.5,0.9)(-4,1.0)(-3.5,1.2)(-3,1.3)(-2.5,1.4)(-2,1.5)(-1.5,1.7)(-1,1.8)(-0.5,1.9)
          (-4.2,1.1)(-3.8,0.8)(-2.8,1.6)(-1.8,1.4)(-0.8,2.1)
        };
      % Scatter (treated side)
      \addplot[only marks, mark=*, mark size=1.5pt, color=econblue, opacity=0.7]
        coordinates {
          (0.5,3.2)(1,3.5)(1.5,3.6)(2,3.8)(2.5,3.9)(3,4.1)(3.5,4.2)(4,4.4)(4.5,4.6)
          (0.8,2.9)(1.8,4.0)(2.8,3.7)(3.8,4.5)
        };

      % Jump annotation
      \draw[<->, very thick, red!70!black] (axis cs:0.2, 1.95) -- (axis cs:0.2, 3.1)
        node[midway, right, font=\small\bfseries, red!70!black] {$\hat{\tau}_{RD}$};

      \node[font=\footnotesize, gray] at (axis cs:-3, 0.4) {$D_i=0$ (untreated)};
      \node[font=\footnotesize, econblue] at (axis cs:3, 2.5) {$D_i=1$ (treated)};
    \end{axis}
  \end{tikzpicture}
  \caption{A sharp Regression Discontinuity Design. Treatment $D_i = \mathbf{1}[R_i \geq c]$ switches on at the cutoff $c$. Units just left of $c$ are nearly identical to units just right, except for treatment status. The RDD estimate $\hat{\tau}_{RD}$ equals the jump in the outcome at the cutoff, computed by comparing the right- and left-hand limits of the regression function.}
  \label{fig:ch13_rdd_plot}
\end{figure}

\medskip\hrule\medskip

\section{13.3 Estimation: Local Linear Regression}

\subsection{13.3.1 Fitting Regression on Each Side}

The naive approach: fit a separate polynomial regression on each side of the cutoff, then compare the intercepts at $c$:

\textbf{Left side} ($R_i < c$): $Y_i = f_L(R_i) + u_i$

\textbf{Right side} ($R_i \geq c$): $Y_i = f_R(R_i) + u_i$

\[
\hat{\tau}_{RDD} = \hat{f}_R(c) - \hat{f}_L(c)
\]

\subsection{13.3.2 Local Linear Regression}

The most common approach: use \textbf{local linear regression} (LLR), which fits a linear regression within a bandwidth $h$ on each side of the cutoff.

\textbf{Definition 13.2 (Local Linear Regression Estimator):} 

For the right limit: $\hat{\mu}_R = \hat{\beta}_{R0}$ from OLS of $Y_i$ on $(1, R_i - c)$ using only observations with $c \leq R_i \leq c + h$.

For the left limit: $\hat{\mu}_L = \hat{\beta}_{L0}$ from OLS of $Y_i$ on $(1, R_i - c)$ using only observations with $c - h \leq R_i < c$.

\[
\hat{\tau}_{RDD}^{LLR} = \hat{\mu}_R - \hat{\mu}_L
\]

\textbf{In one regression (common implementation):}

\[
Y_i = \alpha + \tau D_i + \beta_1 (R_i - c) + \beta_2 D_i(R_i - c) + u_i
\]

using only observations with $|R_i - c| \leq h$. The coefficient $\hat{\tau}$ is the RDD estimate.

This specification allows the slope of $Y$ on $R$ to differ on either side of the cutoff, which is important for flexibility.

\subsection{13.3.3 Kernels and Weighting}

To give more weight to observations near the threshold, use \textbf{kernel regression}. A common choice is the \textbf{triangular kernel}: $K(r) = (1 - |r|)\mathbb{1}[|r| \leq 1]$, which gives zero weight to observations at the bandwidth edge and maximum weight to observations at the cutoff.

Implementing triangular kernel in practice: compute weights $w_i = 1 - |R_i - c|/h$ and use WLS.

\medskip\hrule\medskip

\section{13.4 Bandwidth Selection}

\subsection{13.4.1 The Bias-Variance Tradeoff}

The choice of bandwidth $h$ involves a fundamental tradeoff:

\begin{itemize}
\item \textbf{Larger $h$:} More observations, lower variance, but more bias (farther observations are less comparable).
\item \textbf{Smaller $h$:} Fewer observations, higher variance, but less bias (observations are more comparable).
\end{itemize}

\subsection{13.4.2 Optimal Bandwidth Selection}

The \textbf{Mean Squared Error (MSE)}-optimal bandwidth balances bias and variance:

\[
h^* = \arg\min_h \left[\text{Bias}(\hat{\tau})^2 + \text{Var}(\hat{\tau})\right]
\]

\textbf{Calonico, Cattaneo, and Titiunik (2014)} (CCT) propose a data-driven bandwidth selector based on minimizing the asymptotic MSE of the RDD estimator. The \texttt{rdrobust} package in R implements this.

\textbf{Good practice:} Report results under multiple bandwidths. If estimates are sensitive to bandwidth choice, this is a concern. The CCT optimal bandwidth provides a principled default.

\medskip\hrule\medskip

\section{13.5 The Fuzzy RDD and Its Relationship to IV}

\subsection{13.5.1 Setup}

In the \textbf{fuzzy RDD}, crossing the threshold changes the \textit{probability} of treatment but does not perfectly determine it:

\[
P(D_i = 1 \mid R_i = r) = \begin{cases} g_1(r) & \text{if } r \geq c \\ g_0(r) & \text{if } r < c \end{cases}
\]

where there is a jump $g_1(c) - g_0(c) \neq 0$ but this jump is not a full $0 \to 1$ transition.

\textbf{Example:} A test score above a threshold makes individuals \textit{eligible} for a program, but they must actively enroll. Some who score above the threshold don't enroll; some below find other ways to enroll.

\subsection{13.5.2 Fuzzy RDD as IV}

The fuzzy RDD estimator is:

\[
\hat{\tau}^{fuzzy} = \frac{\lim_{r\downarrow c}E[Y\mid R=r] - \lim_{r\uparrow c}E[Y\mid R=r]}{\lim_{r\downarrow c}P(D=1\mid R=r) - \lim_{r\uparrow c}P(D=1\mid R=r)}
\]

This is exactly the \textbf{Wald estimator} with the threshold indicator $\mathbb{1}[R_i \geq c]$ as the instrument:
\begin{itemize}
\item Numerator: reduced form (effect of threshold on $Y$).
\item Denominator: first stage (effect of threshold on $D$).
\end{itemize}

The fuzzy RDD therefore inherits the LATE interpretation: it estimates the average effect of treatment for compliers --- individuals near the threshold who are induced to take (or not take) treatment by the threshold rule.

\begin{figure}[htbp]
  \centering
  \begin{tikzpicture}
    \begin{axis}[
      width=10cm, height=6cm,
      xlabel={Running Variable $R_i$},
      ylabel={$P(D_i=1 \mid R_i=r)$},
      xmin=-5, xmax=5,
      ymin=-0.05, ymax=1.15,
      axis lines=left,
      tick style={draw=none},
      xtick={0},
      xticklabels={$c$},
      ytick={0,0.2,0.75,1},
      yticklabels={0, $p^-$, $p^+$, 1},
    ]
      % Sharp RDD (for reference, dashed)
      \addplot[thick, gray!50, dashed, domain=-5:-0.05] {0.05};
      \addplot[thick, gray!50, dashed, domain=0.05:5] {0.98};

      % Fuzzy RDD
      \addplot[very thick, econblue, domain=-5:-0.05, samples=50]
        {0.2 + 0.01*x};
      \addplot[very thick, econblue, domain=0.05:5, samples=50]
        {0.75 + 0.01*x};

      % Jump arrow
      \draw[<->, very thick, red!70!black] (axis cs:0.3, 0.2) -- (axis cs:0.3, 0.75)
        node[midway, right, font=\small, red!70!black] {$p^+ - p^-$};

      % Cutoff
      \draw[gray!50, dashed, thick] (axis cs:0,-0.05) -- (axis cs:0,1.15);

      \node[font=\footnotesize, econblue] at (axis cs:-3, 0.35) {Fuzzy RDD};
      \node[font=\footnotesize, gray!60] at (axis cs:3, 0.85) {Sharp RDD (ref.)};
      \node[font=\footnotesize, gray, align=center] at (axis cs:-3, -0.03)
        {$\hat{\tau}_{FRD} = \dfrac{\text{jump in }E[Y|R=r]}{\text{jump in }P(D=1|R=r)}$};
    \end{axis}
  \end{tikzpicture}
  \caption{Fuzzy RDD: the probability of treatment $P(D_i=1|R_i=r)$ jumps discontinuously at the cutoff $c$ but does not go from 0 to 1 (as in sharp RDD). The fuzzy RDD estimator is the ratio of the jump in outcomes to the jump in treatment probability, and is equivalent to a local IV estimator at the cutoff.}
  \label{fig:ch13_fuzzy_rdd}
\end{figure}

\medskip\hrule\medskip

\section{13.6 Validity Tests}

\subsection{13.6.1 The Density Test (McCrary, 2008)}

If individuals can manipulate the running variable to sort above the threshold (e.g., to get a scholarship), the assumption of quasi-random assignment near the threshold is violated.

\textbf{Test:} If there is manipulation, there will be a discontinuity in the \textbf{density of $R_i$} at $c$ --- a "bunching" above the threshold.

The \textbf{McCrary (2008) density test} tests for a discontinuity in the density of the running variable at the cutoff. Rejection is evidence of manipulation and undermines the RDD design.

In R: \texttt{rddensity} package implements this test.

\subsection{13.6.2 Covariate Balance}

If the RDD is valid, covariates that are predetermined relative to the running variable (age, education, gender, baseline outcomes) should not jump at the threshold. \textbf{Test:} Run the same RDD regression using predetermined covariates as the outcome. Non-zero results indicate a problem.

\subsection{13.6.3 Placebo Cutoffs}

If we use an artificial cutoff away from the true cutoff, we should find no discontinuity in outcomes (since there is no treatment at the placebo cutoff). Significant effects at placebo cutoffs suggest confounds.

\medskip\hrule\medskip

\section{13.7 Empirical Illustration: Thistlethwaite and Campbell (1960)}

\subsection{13.7.1 The Original RDD Paper}

Thistlethwaite and Campbell (1960) studied whether receiving a National Merit Award affected students' career aspirations. Awards were granted based on test scores exceeding a threshold.

\textbf{Running variable:} Test score.  
\textbf{Treatment:} Receiving a National Merit Award.  
\textbf{Outcome:} Career and educational aspirations.

Students just above and just below the threshold had nearly identical academic preparation (the difference in test score is negligible), but those above received recognition and financial support. Any discontinuity in outcomes at the threshold can be attributed to the award.

\textbf{This paper introduced RDD to social science.} While Thistlethwaite and Campbell did not use modern estimation methods, their design intuition was correct.

\subsection{13.7.2 Modern Applications of RDD}

RDD has been applied to a vast range of questions:

\begin{center}
\begin{tabular}{@{} l l l l @{}}
\toprule
\textbf{Study} & \textbf{Running Variable} & \textbf{Threshold} & \textbf{Outcome} \\
\midrule
Lee (2008) & Vote share & 50\% & Re-election incumbent advantage \\
Card et al. (2008) & Age & 65 (Medicare eligibility) & Health insurance, health outcomes \\
Angrist and Lavy (1999) & Class size formula & Maimonides rule & Student achievement \\
Dell (2010) & Poverty index & Eligibility cutoff & Economic development \\
\bottomrule
\end{tabular}
\end{center}

\medskip\hrule\medskip

\section{Chapter Summary}

\begin{itemize}
\item \textbf{RDD} exploits threshold-based treatment assignment to identify causal effects.
\item The \textbf{continuity assumption} (outcomes are smooth in the running variable absent treatment) is the key identifying condition.
\item The \textbf{RDD estimand} is the treatment effect at the threshold --- a local effect for individuals right at the margin.
\item \textbf{Local linear regression} within a bandwidth estimates the jump in conditional means at the cutoff.
\item \textbf{Bandwidth selection} involves a bias-variance tradeoff; the CCT method provides a principled choice.
\item The \textbf{fuzzy RDD} is an IV estimator using the threshold indicator as an instrument; it estimates the LATE for compliers near the threshold.
\item \textbf{Validity tests} (density, covariate balance, placebo cutoffs) assess the plausibility of the continuity assumption.
\item RDD estimates are local and may not generalize to populations far from the threshold.
\end{itemize}

\medskip\hrule\medskip

\section{Key Terms}

\textbf{Running variable} --- the continuous variable whose value determines treatment.

\textbf{Cutoff} --- the threshold value of the running variable at which treatment changes.

\textbf{Sharp RDD} --- treatment is perfectly determined by the running variable crossing the threshold.

\textbf{Fuzzy RDD} --- crossing the threshold changes the probability but not the certainty of treatment.

\textbf{Continuity assumption} --- conditional means are continuous in the running variable at the cutoff.

\textbf{Local linear regression (LLR)} --- fitting a linear regression within a bandwidth on each side of the cutoff.

\textbf{Bandwidth} --- the window around the cutoff used for estimation.

\textbf{McCrary density test} --- test for a discontinuity in the density of the running variable at the cutoff.

\textbf{Placebo cutoffs} --- testing for discontinuities at false cutoffs to assess confounds.

\medskip\hrule\medskip

\section{Exercises}

\subsection{Conceptual Questions}

\textbf{13.1} Explain the continuity assumption in plain language. What does it mean about individuals near the threshold?

\textbf{13.2} Why is the effect estimated by RDD "local"? Does this make RDD less useful than DiD or IV?

\textbf{13.3} Explain why the fuzzy RDD is equivalent to IV with the threshold indicator as the instrument. What is the LATE interpretation of the fuzzy RDD?

\textbf{13.4} Describe three validity tests for RDD. What does each test check?

\textbf{13.5} A researcher studies the effect of receiving a college scholarship (awarded to students scoring above 80 on a national exam) on college enrollment. The density of test scores shows a sharp spike just above 80. What does this suggest? How does it affect the validity of the RDD?

\subsection{Analytical Questions}

\textbf{13.6} Suppose the true outcome equation is $Y_i = \tau D_i + g(R_i) + u_i$, where $g(r)$ is a smooth function. Derive the RDD estimator and show it identifies $\tau$ under the continuity assumption.

\textbf{13.7} In a fuzzy RDD, the reduced form estimate is 0.08 and the first stage estimate is 0.40. Compute the LATE. What is the LATE interpretation?

\subsection{Applied Questions}

\textbf{13.8} Using the \texttt{rdrobust} package in R, implement a sharp RDD:
(a) Plot the running variable and outcomes around the cutoff.
(b) Conduct a density test for manipulation.
(c) Estimate the RDD effect using local linear regression with the CCT optimal bandwidth.
(d) Check covariate balance at the cutoff.
(e) Assess sensitivity to bandwidth choice.

\medskip\hrule\medskip

\section{R Lab 13: Regression Discontinuity Design}

\subsection{Setup}

\begin{lstlisting}
library(tidyverse)
library(rdrobust)    # RDD estimation and inference
library(rddensity)   # McCrary density test
library(ggplot2)

# Simulate RDD data
set.seed(42)
n <- 2000
c <- 0  # cutoff

R <- runif(n, -1, 1)    # running variable
D <- as.integer(R >= c)  # sharp treatment

# True treatment effect: tau = 0.5
tau <- 0.5
Y <- 2 + tau * D + 0.8 * R + rnorm(n, sd = 0.5)

df_rdd <- data.frame(R = R, D = D, Y = Y)
\end{lstlisting}

\subsection{Exercise 13.1: Visualizing the Discontinuity}

\begin{lstlisting}
# Bin means plot
df_rdd <- df_rdd %>%
  mutate(bin = cut(R, breaks = 40))

bin_means <- df_rdd %>%
  group_by(bin, D) %>%
  summarise(Y_mean = mean(Y), R_mean = mean(R), .groups = "drop")

ggplot(bin_means, aes(x = R_mean, y = Y_mean, color = factor(D))) +
  geom_point(size = 2) +
  geom_vline(xintercept = 0, linetype = "dashed") +
  geom_smooth(method = "lm", se = FALSE) +
  labs(
    title = "RDD: Outcome vs. Running Variable",
    subtitle = "Clear jump at cutoff (R = 0)",
    x = "Running Variable (R)", y = "Mean Outcome", color = "Treated"
  ) +
  theme_minimal()
\end{lstlisting}

\subsection{Exercise 13.2: Density Test}

\begin{lstlisting}
# McCrary density test
density_test <- rddensity(R, c = c)
summary(density_test)
rdplotdensity(density_test, R)
\end{lstlisting}

\subsection{Exercise 13.3: Local Linear RDD Estimation}

\begin{lstlisting}
# RDD estimate with optimal bandwidth
rdd_est <- rdrobust(Y, R, c = c, kernel = "triangular", bwselect = "mserd")
summary(rdd_est)

cat("RDD estimate:", round(rdd_est$coef["Conventional"], 3), "\n")
cat("95% CI: [", round(rdd_est$ci["Robust", 1], 3), ",",
                 round(rdd_est$ci["Robust", 2], 3), "]\n")
cat("Optimal bandwidth:", round(rdd_est$bws["h", "left"], 3), "\n")
\end{lstlisting}

\subsection{Exercise 13.4: Covariate Balance}

\begin{lstlisting}
# Generate a predetermined covariate
X_cov <- 0.3 * R + rnorm(n)  # correlated with R but not jumping at cutoff

# RDD of covariate on running variable (should find no discontinuity)
rdd_cov <- rdrobust(X_cov, R, c = c)
summary(rdd_cov)
cat("Covariate RDD estimate (should be ~0):", round(rdd_cov$coef["Conventional"], 3), "\n")
\end{lstlisting}

\subsection{Exercise 13.5: Sensitivity to Bandwidth}

\begin{lstlisting}
# Estimate over a range of bandwidths
h_grid <- seq(0.1, 0.8, by = 0.05)
results <- map_dfr(h_grid, function(h) {
  est <- rdrobust(Y, R, c = c, h = h)
  data.frame(h = h, estimate = est$coef["Conventional"],
             ci_low = est$ci["Robust", 1], ci_high = est$ci["Robust", 2])
})

ggplot(results, aes(x = h, y = estimate)) +
  geom_ribbon(aes(ymin = ci_low, ymax = ci_high), alpha = 0.2, fill = "steelblue") +
  geom_line(color = "steelblue", linewidth = 1) +
  geom_hline(yintercept = tau, color = "red", linetype = "dashed") +
  labs(
    title = "RDD Estimate Across Bandwidths",
    subtitle = "Red dashed = true effect. Estimate is stable across bandwidths.",
    x = "Bandwidth (h)", y = "RDD Estimate"
  ) +
  theme_minimal()
\end{lstlisting}

\medskip\hrule\medskip

\textit{End of Chapter 13}

\medskip\hrule\medskip

\textbf{References}

Thistlethwaite, D. L., and Campbell, D. T. (1960). Regression-discontinuity analysis: An alternative to the ex post facto experiment. \textit{Journal of Educational Psychology}, 51(6), 309--317.

Lee, D. S., and Lemieux, T. (2010). Regression discontinuity designs in economics. \textit{Journal of Economic Literature}, 48(2), 281--355.

Calonico, S., Cattaneo, M. D., and Titiunik, R. (2014). Robust nonparametric confidence intervals for regression-discontinuity designs. \textit{Econometrica}, 82(6), 2295--2326.

McCrary, J. (2008). Manipulation of the running variable in the regression discontinuity design: A density test. \textit{Journal of Econometrics}, 142(2), 698--714.

Angrist, J. D., and Lavy, V. (1999). Using Maimonides' rule to estimate the effect of class size on scholastic achievement. \textit{Quarterly Journal of Economics}, 114(2), 533--575.

Cunningham, S. (2021). \textit{Causal Inference: The Mixtape}. Chapter 6.


%% ── Back matter ───────────────────────────────────────────
\backmatter

%% ── Index (placeholder) ───────────────────────────────────
\chapter*{Notation Index}
\addcontentsline{toc}{chapter}{Notation Index}

\begin{tabular}{ll}
\toprule
\textbf{Symbol} & \textbf{Meaning} \\
\midrule
$Y_i, y_i$              & Outcome variable for unit $i$ \\
$D_i$                   & Binary treatment indicator \\
$Y_{0i}, Y_{1i}$        & Potential outcomes (untreated, treated) \\
$\boldsymbol{\beta}$    & Parameter vector \\
$\hat{\boldsymbol{\beta}}^{OLS}$ & OLS estimator \\
$\mathbf{X}$            & $n \times k$ design matrix \\
$\mathbf{P_X}$          & Projection matrix: $\mathbf{X}(\mathbf{X}^T\mathbf{X})^{-1}\mathbf{X}^T$ \\
$\mathbf{M_X}$          & Annihilator: $\mathbf{I} - \mathbf{P_X}$ \\
$\hat{\mathbf{u}}$      & OLS residual vector \\
$\sigma^2$              & Error variance \\
$\hat{\sigma}^2$        & Estimated error variance: $SSR/(n-k)$ \\
$R^2$                   & Coefficient of determination \\
$\boldsymbol{\Omega}$   & Sandwich ``meat'': $E[u_i^2 \mathbf{x}_i\mathbf{x}_i^T]$ \\
$p(\mathbf{x})$         & Propensity score: $P(D_i=1|\mathbf{X}_i=\mathbf{x})$ \\
$\xrightarrow{p}$       & Convergence in probability \\
$\xrightarrow{d}$       & Convergence in distribution \\
ATT                     & Average Treatment Effect on the Treated \\
ATE                     & Average Treatment Effect \\
LATE                    & Local Average Treatment Effect \\
IV                      & Instrumental Variable \\
2SLS                    & Two-Stage Least Squares \\
FE                      & Fixed Effects \\
DiD                     & Differences-in-Differences \\
RDD                     & Regression Discontinuity Design \\
FWL                     & Frisch-Waugh-Lovell (Theorem) \\
OVB                     & Omitted Variable Bias \\
CIA                     & Conditional Independence Assumption \\
\bottomrule
\end{tabular}

\end{document}
